\input{preamble}

% OK, start here.
%
\begin{document}

\title{Faisceaux différentiels gradués}


\maketitle

\phantomsection
\label{section-phantom}

\tableofcontents

\section{Introduction}
\label{section-introduction}

\noindent
Ce chapitre poursuit la discussion commencée dans
Algèbre différentielle graduée, section \ref{dga-section-introduction}.
On trouvera un article de synthèse dans \cite{Keller-survey}.




\section{Conventions}
\label{section-conventions}

\noindent
Dans ce chapitre, nous conservons la convention selon laquelle {\it anneau} désigne un
anneau commutatif possédant un élément unité $1$. Si $R$ est un anneau, une {\it $R$-algèbre $A$}
sera un $R$-module $A$ muni d'une application $R$-bilinéaire $A \times A \to A$
(la multiplication) telle que la multiplication soit associative et admette un
élément unité.
Autrement dit, ce sont des $R$-algèbres associatives unitaires
telles que l'homomorphisme $R \to A$ définissant la structure d'algèbre a son image dans le centre de $A$.








\section{Faisceaux d'algèbres graduées}
\label{section-ga}

\noindent
On pourra omettre cette section.

\begin{definition}
\label{definition-ga}
Soit $(\mathcal{C}, \mathcal{O})$ un site annelé. Un
{\it faisceau d'$\mathcal{O}$-algèbres graduées}
ou un {\it faisceau d'algèbres graduées} sur $(\mathcal{C}, \mathcal{O})$
est la donnée d'une famille $\mathcal{A}^n$ indexée par $n \in \mathbf{Z}$
d'$\mathcal{O}$-modules munis d'applications $\mathcal{O}$-bilinéaires
$$
\mathcal{A}^n \times \mathcal{A}^m \to \mathcal{A}^{n + m},\quad
(a, b) \longmapsto ab
$$
appelées applications de multiplication et vérifiant les propriétés suivantes :
\begin{enumerate}
\item la multiplication est associative, et
\item il existe une section globale $1$ de $\mathcal{A}^0$
qui est un élément unité bilatère pour la multiplication.
\end{enumerate}
Nous notons souvent une telle structure $\mathcal{A}$.
Un {\it homomorphisme d'$\mathcal{O}$-algèbres graduées}
$f : \mathcal{A} \to \mathcal{B}$ est une famille d'homomorphismes
$f^n : \mathcal{A}^n \to \mathcal{B}^n$
d'$\mathcal{O}$-modules compatibles avec les applications de multiplication.
\end{definition}

\noindent
Étant donnés une $\mathcal{O}$-algèbre graduée $\mathcal{A}$
et un objet $U \in \Ob(\mathcal{C})$, nous employons la
notation $$
\mathcal{A}(U)
= \Gamma(U, \mathcal{A})
= \bigoplus\nolimits_{n \in \mathbf{Z}} \mathcal{A}^n(U)
$$
C'est une $\mathcal{O}(U)$-algèbre graduée.

\begin{remark}
\label{remark-functoriality-ga}
Soit $(f, f^\sharp) : (\Sh(\mathcal{C}), \mathcal{O}_\mathcal{C})
\to (\Sh(\mathcal{D}), \mathcal{O}_\mathcal{D})$ un
morphisme de topos annelés. Nous avons :
\begin{enumerate}
\item Soit $\mathcal{A}$ une $\mathcal{O}_\mathcal{C}$-algèbre graduée. Les
applications de multiplication de $\mathcal{A}$ induisent des applications
$f_*\mathcal{A}^n \times f_*\mathcal{A}^m \to f_*\mathcal{A}^{n + m}$ que,
grâce à $f^\sharp$, nous pouvons considérer comme $\mathcal{O}_\mathcal{D}$-bilinéaires.
Nous noterons $f_*\mathcal{A}$ la
$\mathcal{O}_\mathcal{D}$-algèbre graduée ainsi obtenue.
\item Soit $\mathcal{B}$ une
$\mathcal{O}_\mathcal{D}$-algèbre graduée. Les
applications de multiplication de $\mathcal{B}$ induisent des applications
$f^*\mathcal{B}^n \times f^*\mathcal{B}^m \to f^*\mathcal{B}^{n + m}$ que,
grâce à $f^\sharp$, nous pouvons considérer comme $\mathcal{O}_\mathcal{C}$-bilinéaires.
Nous noterons $f^*\mathcal{B}$
la $\mathcal{O}_\mathcal{C}$-algèbre graduée ainsi obtenue.
\item L'ensemble des homomorphismes $f^*\mathcal{B} \to \mathcal{A}$
de $\mathcal{O}_\mathcal{C}$-algèbres graduées est en
correspondance biunivoque ($1$ à $1$) avec l'ensemble des homomorphismes
$\mathcal{B} \to f_*\mathcal{A}$ de $\mathcal{O}_\mathcal{C}$-algèbres graduées.
\end{enumerate} La
partie (3) résulte immédiatement de l'adjonction usuelle entre $f^*$ et $f_*$ sur les
faisceaux de modules.
\end{remark}




\section{Faisceaux de modules gradués}
\label{section-graded-modules}

\noindent
On pourra omettre cette section.

\begin{definition}
\label{definition-gm} Soit
$(\mathcal{C}, \mathcal{O})$ un site annelé. Soit
$\mathcal{A}$ un faisceau d'algèbres graduées sur
$(\mathcal{C}, \mathcal{O})$. Un
{\it $\mathcal{A}$-module gradué (à droite)}, ou un
{\it module gradué (à droite)} sur $\mathcal{A}$, est la
donnée d'une famille $\mathcal{M}^n$, indexée par $n \in \mathbf{Z}$,
de $\mathcal{O}$-modules munis d'applications
$\mathcal{O}$-bilinéaires
$$
\mathcal{M}^n \times \mathcal{A}^m \to \mathcal{M}^{n + m},\quad (x,
a) \longmapsto xa
$$
appelées applications de multiplication et vérifiant les propriétés
suivantes : \begin{enumerate}
\item la multiplication vérifie $(xa)a' = x(aa')$,
\item la section unité $1$ de $\mathcal{A}^0$
agit comme l'identité sur $\mathcal{M}^n$ pour tout $n$.
\end{enumerate}
Nous dirons souvent « soit $\mathcal{M}$ un $\mathcal{A}$-module gradué »
pour désigner cette situation.
Un {\it homomorphisme de $\mathcal{A}$-modules gradués}
$f : \mathcal{M} \to \mathcal{N}$ est une famille d'homomorphismes
$f^n : \mathcal{M}^n \to \mathcal{N}^n$
de $\mathcal{O}$-modules, compatible avec les applications de multiplication.
La catégorie des $\mathcal{A}$-modules gradués (à
droite) est notée $\textit{Mod}(\mathcal{A})$.
\end{definition}

\noindent On
définit exactement de la même manière les {\it modules gradués à gauche}, mais, dans
ce chapitre, les modules seront par défaut des modules à droite.

\medskip\noindent
Étant donnés un $\mathcal{A}$-module gradué $\mathcal{M}$
et un objet $U \in \Ob(\mathcal{C})$, nous employons la
notation $$
\mathcal{M}(U)
= \Gamma(U, \mathcal{M})
= \bigoplus\nolimits_{n \in \mathbf{Z}} \mathcal{M}^n(U)
$$
C'est un $\mathcal{A}(U)$-module gradué (à droite).

\begin{lemma}
\label{lemma-gm-abelian} Soit
$(\mathcal{C}, \mathcal{O})$ un site annelé. Soit
$\mathcal{A}$ une $\mathcal{O}$-algèbre graduée. La catégorie
$\textit{Mod}(\mathcal{A})$ est une catégorie abélienne possédant les
propriétés suivantes :
\begin{enumerate}
\item $\textit{Mod}(\mathcal{A})$ admet les sommes directes arbitraires,
\item $\textit{Mod}(\mathcal{A})$ admet les limites inductives arbitraires,
\item les limites inductives filtrantes dans $\textit{Mod}(\mathcal{A})$ sont exactes,
\item $\textit{Mod}(\mathcal{A})$ admet les produits arbitraires,
\item $\textit{Mod}(\mathcal{A})$ admet les limites arbitraires.
\end{enumerate}
Le foncteur
$$
\textit{Mod}(\mathcal{A}) \longrightarrow \textit{Mod}(\mathcal{O}),\quad
\mathcal{M} \longmapsto \mathcal{M}^n
$$ qui
envoie un $\mathcal{A}$-module gradué sur son terme de degré $n$ commute à
toutes les limites et limites inductives.
\end{lemma}

\noindent
Le lemme affirme que les limites et limites inductives se calculent terme à terme. Il
affirme aussi (ou l'implique, si l'on préfère) que le foncteur d'oubli
$$
\textit{Mod}(\mathcal{A})
\longrightarrow
\text{les }\mathcal{O}\text{-modules gradués}
$$
commute à toutes les limites et limites inductives.

\begin{proof}
Notons
$\text{gr}^n : \textit{Mod}(\mathcal{A}) \to \textit{Mod}(\mathcal{O})$ le foncteur
de l'énoncé. Considérons un homomorphisme
$f : \mathcal{M} \to \mathcal{N}$ de
$\mathcal{A}$-modules gradués. Le noyau et le conoyau de
$f$, calculés comme morphismes de $\mathcal{O}$-modules gradués, sont eux
aussi munis d'applications de multiplication comme dans la
définition \ref{definition-gm}. Ce sont donc également
le noyau et le conoyau dans $\textit{Mod}(\mathcal{A})$.
Ainsi, $\textit{Mod}(\mathcal{A})$ est une catégorie abélienne,
et la prise des noyaux et conoyaux commute à $\text{gr}^n$.

\medskip\noindent
Pour démontrer l'existence des limites et limites inductives, il suffit de
démontrer l'existence des produits et des sommes directes, voir
Catégories, lemmes \ref{categories-lemma-limits-products-equalizers} et
\ref{categories-lemma-colimits-coproducts-coequalizers}. Les mêmes
lemmes montrent que, pour
établir que $\text{gr}^n$ commute aux limites et limites inductives, il suffit
d'établir que $\text{gr}^n$ commute aux sommes directes et aux produits.

\medskip\noindent
Soit $\mathcal{M}_t$, $t \in T$, une famille de $\mathcal{A}$-modules gradués.
Considérons le $\mathcal{A}$-module gradué dont le terme de degré $n$ est
$\bigoplus_{t \in T} \mathcal{M}_t^n$ (muni des applications de multiplication évidentes).
Le lecteur vérifie aisément qu'il s'agit d'une somme directe dans
$\textit{Mod}(\mathcal{A})$. Il en va de même des produits.

\medskip\noindent
Observons que $\text{gr}^n$ est un foncteur exact pour tout $n$ et qu'un
complexe $\mathcal{M}_1 \to \mathcal{M}_2 \to \mathcal{M}_3$ de
$\textit{Mod}(\mathcal{A})$ est exact si et seulement si
$\text{gr}^n\mathcal{M}_1 \to \text{gr}^n\mathcal{M}_2 \to
\text{gr}^n\mathcal{M}_3$ est exact dans $\textit{Mod}(\mathcal{O})$ pour
tout $n$. On en déduit (3), puisque les limites inductives filtrantes sont
exactes dans $\textit{Mod}(\mathcal{O})$ ; c'est une
catégorie abélienne de Grothendieck, voir
Cohomologie sur les sites, section \ref{sites-cohomology-section-unbounded}.
\end{proof}






\section{La catégorie graduée des faisceaux de modules gradués}
\label{section-gm-gr-cat}

\noindent
On pourra omettre cette section. Cette section est
l'analogue de Algèbre différentielle graduée, exemple \ref{dga-example-gm-gr-cat}.
Pour nos conventions sur les catégories graduées,
voir Algèbre différentielle graduée, section \ref{dga-section-graded}.

\medskip\noindent
Soit $(\mathcal{C}, \mathcal{O})$ un site annelé. Soit
$\mathcal{A}$ un faisceau d'algèbres graduées sur
$(\mathcal{C}, \mathcal{O})$. Nous allons construire une
catégorie graduée $\textit{Mod}^{gr}(\mathcal{A})$ sur
$R = \Gamma(\mathcal{C}, \mathcal{O})$ dont
la catégorie associée $(\textit{Mod}^{gr}(\mathcal{A}))^0$ est la
catégorie des $\mathcal{A}$-modules gradués. Les objets de
$\textit{Mod}^{gr}(\mathcal{A})$ sont les
$\mathcal{A}$-modules gradués à droite (voir la section
\ref{section-graded-modules}). Étant donnés des
$\mathcal{A}$-modules gradués $\mathcal{L}$ et $\mathcal{M}$, posons
$$
\Hom_{\textit{Mod}^{gr}(\mathcal{A})}(\mathcal{L}, \mathcal{M}) =
\bigoplus\nolimits_{n \in \mathbf{Z}} \Hom^n(\mathcal{L}, \mathcal{M})
$$
où
$\Hom^n(\mathcal{L}, \mathcal{M})$ est l'ensemble
des homomorphismes de $\mathcal{A}$-modules à droite
$f : \mathcal{L} \to \mathcal{M}$ homogènes
de degré $n$. Plus précisément, $f$ est la donnée d'une
famille d'homomorphismes $f : \mathcal{L}^i \to \mathcal{M}^{i + n}$, pour
$i \in \mathbf{Z}$, compatible avec les applications de multiplication.
En termes de composantes,
$$
\Hom^n(\mathcal{L}, \mathcal{M})
\subset
\prod\nolimits_{p + q = n}
\Hom_\mathcal{O}(\mathcal{L}^{-q}, \mathcal{M}^p)
$$ (remarquer l'inversion
des indices) est le sous-ensemble des
$f = (f_{p, q})$ tels que
$$
f_{p, q}(m a) = f_{p - i, q + i}(m)a
$$ pour toutes sections locales
$a$ de $\mathcal{A}^i$ et
$m$ de $\mathcal{L}^{-q - i}$. Pour des $\mathcal{A}$-modules gradués
$\mathcal{K}$, $\mathcal{L}$ et $\mathcal{M}$, nous définissons la
composition dans $\textit{Mod}^{gr}(\mathcal{A})$ au moyen des
applications
$$
\Hom^m(\mathcal{L}, \mathcal{M}) \times
\Hom^n(\mathcal{K}, \mathcal{L}) \longrightarrow
\Hom^{n + m}(\mathcal{K}, \mathcal{M})
$$ par
simple composition des applications
de $\mathcal{A}$-modules : $(g, f) \mapsto g \circ f$.





\section{Produit tensoriel de faisceaux de modules gradués}
\label{section-tensor-product}

\noindent
On pourra omettre cette section. Cette section est analogue à une partie
d'Algèbre différentielle graduée, section \ref{dga-section-tensor-product}.

\medskip\noindent
Soit $(\mathcal{C}, \mathcal{O})$ un site annelé. Soit $\mathcal{A}$ un
faisceau d'algèbres graduées sur $(\mathcal{C}, \mathcal{O})$.
Soit $\mathcal{M}$ un $\mathcal{A}$-module gradué à droite
et soit $\mathcal{N}$ un $\mathcal{A}$-module gradué à gauche.
Nous définissons alors le {\it produit tensoriel}
$\mathcal{M} \otimes_\mathcal{A} \mathcal{N}$
comme le $\mathcal{O}$-module gradué dont le terme de degré $n$
est $$
(\mathcal{M} \otimes_\mathcal{A} \mathcal{N})^n
= \Coker\left(
\bigoplus\nolimits_{r + s + t = n} \mathcal{M}^r \otimes_\mathcal{O}
\mathcal{A}^s \otimes_\mathcal{O} \mathcal{N}^t
\longrightarrow
\bigoplus\nolimits_{p + q = n} \mathcal{M}^p \otimes_\mathcal{O} \mathcal{N}^q
\right)
$$
où l'application envoie la section locale $x \otimes a \otimes y$
de $\mathcal{M}^r \otimes_\mathcal{O} \mathcal{A}^s
\otimes_\mathcal{O} \mathcal{N}^t$
sur $xa \otimes y - x \otimes ay$.
Avec cette définition,
$(\mathcal{M} \otimes_\mathcal{A} \mathcal{N})^n$ est
le faisceau associé au préfaisceau
$U \mapsto (\mathcal{M}(U) \otimes_{\mathcal{A}(U)} \mathcal{N}(U))^n$, où le
produit tensoriel des modules gradués est celui défini dans
Algèbre différentielle graduée, section \ref{dga-section-tensor-product}.

\medskip\noindent
Si l'on fixe le $\mathcal{A}$-module gradué à gauche $\mathcal{N}$,
on obtient un foncteur
$$
- \otimes_\mathcal{A} \mathcal{N} :
\textit{Mod}(\mathcal{A})
\longrightarrow
\text{Gr}(\textit{Mod}(\mathcal{O})) =
\text{les }\mathcal{O}\text{-modules gradués}
$$
Pour la notation $\text{Gr}(-)$, voir
Homologie, définition \ref{homology-definition-graded}.
La catégorie graduée des $\mathcal{O}$-modules gradués est notée
$\text{Gr}^{gr}(\textit{Mod}(\mathcal{O}))$, voir
Algèbre différentielle graduée, exemple
\ref{dga-example-graded-category-graded-objects}.
Le foncteur ci-dessus se relève en un foncteur de catégories graduées
$$
- \otimes_\mathcal{A} \mathcal{N} :
\textit{Mod}^{gr}(\mathcal{A})
\longrightarrow
\text{Gr}^{gr}(\textit{Mod}(\mathcal{O}))
$$
en envoyant les homomorphismes de degré $n$ de $\mathcal{M} \to \mathcal{M}'$
sur l'application induite de degré $n$ de
$\mathcal{M} \otimes_\mathcal{A} \mathcal{N}$ à
$\mathcal{M}' \otimes_\mathcal{A} \mathcal{N}$.






\section{Hom interne de faisceaux de modules gradués}
\label{section-internal-hom-graded}

\noindent Nous
recommandons vivement au lecteur d'omettre cette section.

\medskip\noindent
Nous aurons besoin de la version faisceautisée de la construction de
la section \ref{section-gm-gr-cat}. Soient
$(\mathcal{C}, \mathcal{O})$, $\mathcal{A}$,
$\mathcal{M}$ et $\mathcal{L}$ comme dans la section \ref{section-gm-gr-cat}. Nous
définissons
$$
\SheafHom^{gr}_\mathcal{A}(\mathcal{M}, \mathcal{L})
$$ comme
le $\mathcal{O}$-module gradué dont le terme de degré $n$
$$
\SheafHom^n_\mathcal{A}(\mathcal{M}, \mathcal{L})
\subset
\prod\nolimits_{p + q = n}
\SheafHom_\mathcal{O}(\mathcal{L}^{-q}, \mathcal{M}^p)
$$ est le
sous-faisceau constitué des sections locales $f = (f_{p, q})$ telles que
$$
f_{p, q}(m a) = f_{p - i, q + i}(m)a
$$ pour toutes
sections locales $a$ de $\mathcal{A}^i$ et
$m$ de $\mathcal{L}^{-q - i}$. Comme dans la section \ref{section-gm-gr-cat}, il
existe une application de
composition $$
\SheafHom^{gr}_\mathcal{A}(\mathcal{L}, \mathcal{M}) \otimes_\mathcal{O}
\SheafHom^{gr}_\mathcal{A}(\mathcal{K}, \mathcal{L})
\longrightarrow
\SheafHom^{gr}_\mathcal{A}(\mathcal{K}, \mathcal{M})
$$ où
le membre de gauche est le produit tensoriel de $\mathcal{O}$-modules gradués défini
dans la section \ref{section-tensor-product}. Cette
application est donnée par l'application de
composition $$
\SheafHom^m_\mathcal{A}(\mathcal{L}, \mathcal{M}) \otimes_\mathcal{O}
\SheafHom^n_\mathcal{A}(\mathcal{K}, \mathcal{L}) \longrightarrow
\SheafHom^{n + m}_\mathcal{A}(\mathcal{K}, \mathcal{M})
$$
définie par simple composition (localement).

\medskip\noindent
Avec ces définitions, nous avons
$$
\Hom_{\textit{Mod}^{gr}(\mathcal{A})}(\mathcal{L}, \mathcal{M}) =
\Gamma(\mathcal{C}, \SheafHom^{gr}_\mathcal{A}(\mathcal{L}, \mathcal{M}))
$$ comme modules
gradués sur $R$, de manière compatible à la composition.







\section{Faisceaux de bimodules gradués et adjonction tensor-Hom}
\label{section-graded-bimodules}

\noindent
On pourra omettre cette section.

\begin{definition}
\label{definition-bimodule}
Soit $(\mathcal{C}, \mathcal{O})$ un site annelé. Soient $\mathcal{A}$ et
$\mathcal{B}$ des faisceaux d'algèbres graduées sur
$(\mathcal{C}, \mathcal{O})$. Un
{\it $(\mathcal{A}, \mathcal{B})$-bimodule gradué} est la
donnée d'une famille $\mathcal{M}^n$, indexée par $n \in \mathbf{Z}$, de
$\mathcal{O}$-modules munis d'applications $\mathcal{O}$-bilinéaires
$$
\mathcal{M}^n \times \mathcal{B}^m \to \mathcal{M}^{n + m},\quad (x,
b) \longmapsto xb
$$
et
$$
\mathcal{A}^n \times \mathcal{M}^m \to \mathcal{M}^{n + m},\quad (a, x)
\longmapsto ax
$$ appelées
applications de multiplication et vérifiant les propriétés
suivantes : \begin{enumerate}
\item la multiplication vérifie $a(a'x) = (aa')x$
et $(xb)b' = x(bb')$,
\item $(ax)b = a(xb)$,
\item la section unité $1$ de $\mathcal{A}^0$ agit comme
l'identité par multiplication, et
\item la section unité $1$ de
$\mathcal{B}^0$ agit comme l'identité par multiplication.
\end{enumerate} Nous notons
souvent une telle structure $\mathcal{M}$. Un
{\it homomorphisme de $(\mathcal{A}, \mathcal{B})$-bimodules gradués}
$f : \mathcal{M} \to \mathcal{N}$ est une famille d'homomorphismes
$f^n : \mathcal{M}^n \to \mathcal{N}^n$ de
$\mathcal{O}$-modules, compatible avec les applications de multiplication.
\end{definition}

\noindent
Étant donnés un $(\mathcal{A}, \mathcal{B})$-bimodule gradué $\mathcal{M}$
et un objet $U \in \Ob(\mathcal{C})$, nous employons la
notation $$
\mathcal{M}(U)
= \Gamma(U, \mathcal{M})
= \bigoplus\nolimits_{n \in \mathbf{Z}} \mathcal{M}^n(U)
$$
C'est un $(\mathcal{A}(U), \mathcal{B}(U))$-bimodule gradué.

\medskip\noindent
Soit $(\mathcal{C}, \mathcal{O})$ un site annelé. Soient $\mathcal{A}$ et
$\mathcal{B}$ des faisceaux d'algèbres graduées sur
$(\mathcal{C}, \mathcal{O})$. Soit $\mathcal{M}$ un
$\mathcal{A}$-module gradué à droite et soit $\mathcal{N}$ un
$(\mathcal{A}, \mathcal{B})$-bimodule gradué. Dans ce cas, le produit tensoriel
gradué défini dans la section \ref{section-tensor-product}
$$
\mathcal{M} \otimes_\mathcal{A} \mathcal{N}
$$ est un
$\mathcal{B}$-module gradué à droite, muni des applications de multiplication évidentes.
Cette construction définit un foncteur et un foncteur de catégories
graduées $$
\otimes_\mathcal{A} \mathcal{N}
: \textit{Mod}(\mathcal{A})
\longrightarrow
\textit{Mod}(\mathcal{B})
\quad\text{et}\quad
\otimes_\mathcal{A} \mathcal{N} :
\textit{Mod}^{gr}(\mathcal{A})
\longrightarrow
\textit{Mod}^{gr}(\mathcal{B})
$$
en envoyant les homomorphismes de degré $n$ de $\mathcal{M} \to \mathcal{M}'$ sur
les applications induites de degré $n$
de $\mathcal{M} \otimes_\mathcal{A} \mathcal{N}$
vers $\mathcal{M}' \otimes_\mathcal{A} \mathcal{N}$.

\medskip\noindent
Soit $(\mathcal{C}, \mathcal{O})$ un site annelé. Soient $\mathcal{A}$ et
$\mathcal{B}$ des faisceaux d'algèbres graduées sur
$(\mathcal{C}, \mathcal{O})$. Soit $\mathcal{N}$ un
$(\mathcal{A}, \mathcal{B})$-bimodule gradué. Soit
$\mathcal{L}$ un $\mathcal{B}$-module gradué à droite. Dans ce
cas, le Hom interne gradué défini dans
la section \ref{section-internal-hom-graded}
$$
\SheafHom_\mathcal{B}^{gr}(\mathcal{N}, \mathcal{L})
$$ est un
$\mathcal{A}$-module gradué à droite, muni des
applications de multiplication\footnote{Selon nos conventions,
aucun signe n'intervient ici.}
$$
\SheafHom^n_\mathcal{B}(\mathcal{N}, \mathcal{L})
\times \mathcal{A}^m
\longrightarrow
\SheafHom^{n + m}_\mathcal{B}(\mathcal{N}, \mathcal{L})
$$ qui envoient
une section $f = (f_{p,q})$ de
$\SheafHom^n_\mathcal{B}(\mathcal{N}, \mathcal{L})$ sur $U$ et une
section $a$ de $\mathcal{A}^m$ sur $U$ sur la section
$f a$ de $\SheafHom^{n + m}_\mathcal{B}(\mathcal{N}, \mathcal{L})$ sur $U$ définie par la
famille d'applications
$$
\mathcal{N}^{-q - m}|_U \xrightarrow{a \cdot -}
\mathcal{N}^{-q}|_U \xrightarrow{f_{p, q}}
\mathcal{M}^p|_U
$$ Nous omettons
la vérification que cette définition est correcte. Cette construction
définit un foncteur et un foncteur de catégories graduées
$$
\SheafHom_\mathcal{B}^{gr}(\mathcal{N}, -)
: \textit{Mod}(\mathcal{B})
\longrightarrow
\textit{Mod}(\mathcal{A})
\quad\text{et}\quad
\SheafHom_\mathcal{B}^{gr}(\mathcal{N}, -) :
\textit{Mod}^{gr}(\mathcal{B})
\longrightarrow
\textit{Mod}^{gr}(\mathcal{A})
$$ en
envoyant les homomorphismes de degré $n$ de $\mathcal{L} \to \mathcal{L}'$ sur
les applications induites de degré $n$
de $\SheafHom_\mathcal{B}^{gr}(\mathcal{N}, \mathcal{L})$
vers $\SheafHom_\mathcal{B}^{gr}(\mathcal{N}, \mathcal{L}')$.

\begin{lemma}
\label{lemma-tensor-hom-adjunction-gr}
Soit $(\mathcal{C}, \mathcal{O})$ un site annelé. Soient $\mathcal{A}$ et
$\mathcal{B}$ des faisceaux d'algèbres graduées sur
$(\mathcal{C}, \mathcal{O})$. Soit $\mathcal{M}$ un
$\mathcal{A}$-module gradué à droite. Soit $\mathcal{N}$ un
$(\mathcal{A}, \mathcal{B})$-bimodule gradué. Soit $\mathcal{L}$ un
$\mathcal{B}$-module gradué à droite. Avec les conventions précédentes,
nous
avons $$
\Hom_{\textit{Mod}^{gr}(\mathcal{B})}(
\mathcal{M} \otimes_\mathcal{A} \mathcal{N}, \mathcal{L})
= \Hom_{\textit{Mod}^{gr}(\mathcal{A})}(
\mathcal{M}, \SheafHom_\mathcal{B}^{gr}(\mathcal{N}, \mathcal{L}))
$$
et
$$
\SheafHom_\mathcal{B}^{gr}(
\mathcal{M} \otimes_\mathcal{A} \mathcal{N}, \mathcal{L}) =
\SheafHom_\mathcal{A}^{gr}(
\mathcal{M}, \SheafHom_\mathcal{B}^{gr}(\mathcal{N}, \mathcal{L}))
$$
fonctoriellement en $\mathcal{M}$, $\mathcal{N}$ et $\mathcal{L}$.
\end{lemma}

\begin{proof}
Omis. Indication : cela résulte de l'interprétation des deux membres
comme applications graduées $\mathcal{A}$-bilinéaires
$\psi : \mathcal{M} \times \mathcal{N} \to \mathcal{L}$,
qui sont $\mathcal{B}$-linéaires à droite.
\end{proof}

\noindent
Soit $(\mathcal{C}, \mathcal{O})$ un site annelé. Soient $\mathcal{A}$ et
$\mathcal{B}$ des faisceaux d'algèbres graduées sur
$(\mathcal{C}, \mathcal{O})$.
Comme cas particulier de ce qui précède, supposons donné
un homomorphisme $\varphi : \mathcal{A} \to \mathcal{B}$
de $\mathcal{O}$-algèbres graduées. On obtient alors un foncteur
et un foncteur de catégories
graduées $$
\otimes_{\mathcal{A}, \varphi} \mathcal{B}
: \textit{Mod}(\mathcal{A})
\longrightarrow
\textit{Mod}(\mathcal{B})
\quad\text{et}\quad
\otimes_{\mathcal{A}, \varphi} \mathcal{B} :
\textit{Mod}^{gr}(\mathcal{A})
\longrightarrow
\textit{Mod}^{gr}(\mathcal{B})
$$
D'autre part, nous avons les foncteurs de restriction
$$
res_\varphi :
\textit{Mod}(\mathcal{B})
\longrightarrow
\textit{Mod}(\mathcal{A})
\quad\text{et}\quad
res_\varphi :
\textit{Mod}^{gr}(\mathcal{B})
\longrightarrow
\textit{Mod}^{gr}(\mathcal{A})
$$ Le
lemme précédent permet de montrer que ces foncteurs sont adjoints
l'un à l'autre (comme d'habitude pour la restriction et le changement
de base). Notons ${}_\mathcal{A}\mathcal{B}_\mathcal{B}$
le faisceau $\mathcal{B}$ considéré comme $(\mathcal{A}, \mathcal{B})$-bimodule
gradué. Pour tout $\mathcal{B}$-module gradué à droite $\mathcal{L}$,
nous
avons alors $$
\SheafHom_\mathcal{B}^{gr}({}_\mathcal{A}\mathcal{B}_\mathcal{B}, \mathcal{L})
= res_\varphi(\mathcal{L})
$$
comme $\mathcal{A}$-modules
gradués à droite. Le lemme \ref{lemma-tensor-hom-adjunction-gr}
fournit donc un isomorphisme
fonctoriel $$
\Hom_{\textit{Mod}^{gr}(\mathcal{B})}(
\mathcal{M} \otimes_{\mathcal{A}, \varphi} \mathcal{B}, \mathcal{L})
= \Hom_{\textit{Mod}^{gr}(\mathcal{A})}(
\mathcal{M}, res_\varphi(\mathcal{L}))
$$
Lorsque le contexte est clair, nous omettons généralement $\varphi$
dans cette formule. De la même manière, nous
obtenons l'égalité
$$
\SheafHom^{gr}_\mathcal{B}(
\mathcal{M} \otimes_\mathcal{A} \mathcal{B}, \mathcal{L})
= \SheafHom_\mathcal{A}^{gr}(\mathcal{M}, \mathcal{L})
$$
de $\mathcal{O}$-modules gradués.




\section{Image inverse et image directe des faisceaux de modules gradués}
\label{section-functoriality-graded}

\noindent
Nous conseillons au lecteur d'omettre cette section.

\medskip\noindent
Soit $(f, f^\sharp) : (\Sh(\mathcal{C}), \mathcal{O}_\mathcal{C})
\to (\Sh(\mathcal{D}), \mathcal{O}_\mathcal{D})$ un morphisme
de topos annelés. Soit $\mathcal{A}$ une
$\mathcal{O}_\mathcal{C}$-algèbre graduée. Soit $\mathcal{B}$ une
$\mathcal{O}_\mathcal{D}$-algèbre graduée. Supposons qu'on nous donne
une application
$$
\varphi : f^{-1}\mathcal{B} \to \mathcal{A}
$$ d'algèbres
graduées sur $f^{-1}\mathcal{O}_\mathcal{D}$. Par l'adjonction
entre restriction et extension des scalaires, cela revient à se
donner un morphisme $\varphi : f^*\mathcal{B} \to \mathcal{A}$ de
$\mathcal{O}_\mathcal{C}$-algèbres graduées ou, de manière équivalente,
$\varphi$ peut être vu comme un morphisme
$$
\varphi : \mathcal{B} \to f_*\mathcal{A}
$$
d'algèbres graduées sur $\mathcal{O}_\mathcal{D}$.
Voir la remarque \ref{remark-functoriality-ga}.

\medskip\noindent
Définissons maintenant un
foncteur $$
f_*
: \textit{Mod}(\mathcal{A})
\longrightarrow
\textit{Mod}(\mathcal{B})
$$
Étant donné un $\mathcal{A}$-module gradué $\mathcal{M}$, nous
définissons $f_*\mathcal{M}$ comme le $\mathcal{B}$-module
gradué dont le terme de degré $n$ est $f_*\mathcal{M}^n$. Comme
multiplication,
nous utilisons $$
f_*\mathcal{M}^n \times \mathcal{B}^m
\xrightarrow{(\text{id}, \varphi^m)}
f_*\mathcal{M}^n \times f_*\mathcal{A}^m
\xrightarrow{f_*\mu_{n, m}}
f_*\mathcal{M}^{n + m}
$$
où $\mu_{n, m} : \mathcal{M}^n \times \mathcal{A}^m
\to \mathcal{M}^{n + m}$ est l'application de multiplication pour $\mathcal{M}$ sur
$\mathcal{A}$. Nous utilisons ici le fait que $f_*$ commute aux produits.
La construction est clairement fonctorielle en
$\mathcal{M}$ et nous obtenons notre foncteur.

\medskip\noindent
Définissons maintenant un
foncteur $$
f^*
: \textit{Mod}(\mathcal{B})
\longrightarrow
\textit{Mod}(\mathcal{A})
$$
Nous définirons ce foncteur comme un composé de foncteurs
$$
\textit{Mod}(\mathcal{B})
\xrightarrow{f^{-1}}
\textit{Mod}(f^{-1}\mathcal{B})
\xrightarrow{ - \otimes_{f^{-1}\mathcal{B}} \mathcal{A}}
\textit{Mod}(\mathcal{A})
$$ Tout d'abord,
étant donné un $\mathcal{B}$-module gradué $\mathcal{N}$, nous définissons
$f^{-1}\mathcal{N}$ comme le $f^{-1}\mathcal{B}$-module gradué
dont le terme de degré $n$ est $f^{-1}\mathcal{N}^n$. Comme multiplication,
nous
utilisons $$
f^{-1}\nu_{n, m}
: f^{-1}\mathcal{N}^n \times f^{-1}\mathcal{B}^m
\longrightarrow
f^{-1}\mathcal{N}^{n + m}
$$
où $\nu_{n, m} : \mathcal{N}^n \times \mathcal{B}^m
\to \mathcal{N}^{n + m}$ est l'application de multiplication pour $\mathcal{N}$ sur
$\mathcal{B}$. Nous utilisons ici le fait que $f^{-1}$ commute aux produits.
La construction est clairement fonctorielle en
$\mathcal{N}$ et nous obtenons notre foncteur $f^{-1}$.
Ceci étant dit, nous pouvons utiliser la
discussion sur le produit tensoriel dans la section \ref{section-graded-bimodules}
pour définir
le foncteur $$
- \otimes_{f^{-1}\mathcal{B}} \mathcal{A}
: \textit{Mod}(f^{-1}\mathcal{B})
\longrightarrow
\textit{Mod}(\mathcal{A})
$$
Enfin, nous
posons $$
f^*\mathcal{N}
= f^{-1}\mathcal{N} \otimes_{f^{-1}\mathcal{B}, \varphi} \mathcal{A}
$$
comme déjà prédit ci-dessus.

\medskip\noindent
Les foncteurs $f_*$ et $f^*$ se relèvent naturellement en des
foncteurs de catégories graduées
$$
f_*
: \textit{Mod}^{gr}(\mathcal{A})
\longrightarrow
\textit{Mod}^{gr}(\mathcal{B})
\quad\text{et}\quad
f^* :
\textit{Mod}^{gr}(\mathcal{B})
\longrightarrow
\textit{Mod}^{gr}(\mathcal{A})
$$
qui font la même chose sur les objets sous-jacents et sont définis
par la fonctorialité des constructions sur les morphismes homogènes
de degré $n$.

\begin{lemma}
\label{lemma-adjunction-push-pull-gr} Dans
la situation ci-dessus nous
avons $$
\Hom_{\textit{Mod}^{gr}(\mathcal{B})}(
\mathcal{N}, f_*\mathcal{M})
= \Hom_{\textit{Mod}^{gr}(\mathcal{A})}(
f^*\mathcal{N}, \mathcal{M})
$$
\end{lemma}

\begin{proof}
Omis. Indications : on commence par prouver que $f^{-1}$ et $f_*$ sont adjoints
en tant que foncteurs entre $\textit{Mod}(\mathcal{B})$
et $\textit{Mod}(f^{-1}\mathcal{B})$ en utilisant l'adjonction entre
$f^{-1}$ et $f_*$ sur les faisceaux de groupes abéliens.
On utilise ensuite l'adjonction entre le changement de base et la restriction
donnée à la section \ref{section-graded-bimodules}.
\end{proof}





\section{Localisation et faisceaux de modules gradués}
\label{section-localize-graded}

\noindent
Nous conseillons au lecteur d'omettre cette section.

\medskip\noindent
Soit $(\mathcal{C}, \mathcal{O})$ un site annelé. Soit
$U \in \Ob(\mathcal{C})$ et désignons
$$ j
:
(\Sh(\mathcal{C}/U), \mathcal{O}_U)
\longrightarrow
(\Sh(\mathcal{C}), \mathcal{O})
$$ le
morphisme de localisation correspondant (Modules
sur les sites, section \ref{sites-modules-section-localize}).
Ci-dessous, nous utiliserons le fait suivant : pour les $\mathcal{O}_U$-modules
$\mathcal{M}_i$, $i = 1, 2$ et un $\mathcal{O}$-module $\mathcal{A}$,
il existe une application
canonique $$
j_!
: \Hom_{\mathcal{O}_U}(
\mathcal{M}_1 \otimes_{\mathcal{O}_U} \mathcal{A}|_U, \mathcal{M}_2)
\longrightarrow
\Hom_\mathcal{O}(
j_!\mathcal{M}_1 \otimes_\mathcal{O} \mathcal{A}, j_!\mathcal{M}_2)
$$
À savoir, nous
avons $j_!(\mathcal{M}_1 \otimes_{\mathcal{O}_U} \mathcal{A}|_U)
= j_!\mathcal{M}_1 \otimes_\mathcal{O} \mathcal{A}$ par
Modules sur les sites, lemme \ref{sites-modules-lemma-j-shriek-and-tensor}.

\medskip\noindent
Soit $\mathcal{A}$ une $\mathcal{O}$-algèbre graduée. Nous
noterons $\mathcal{A}_U$ la restriction de $\mathcal{A}$ à
$\mathcal{C}/U$, en d'autres termes, nous avons
$\mathcal{A}_U = j^*\mathcal{A} = j^{-1}\mathcal{A}$. Dans la section
\ref{section-functoriality-graded} nous avons construit
des foncteurs adjoints
$$ j_*
:
\textit{Mod}^{gr}(\mathcal{A}_U)
\longrightarrow
\textit{Mod}^{gr}(\mathcal{A})
\quad\text{et}\quad j^*
:
\textit{Mod}^{gr}(\mathcal{A})
\longrightarrow
\textit{Mod}^{gr}(\mathcal{A}_U)
$$ avec
$j^*$ adjoint à gauche à $j_*$. Nous affirmons qu'il existe en outre un foncteur
exact
$$ j_!
:
\textit{Mod}^{gr}(\mathcal{A}_U)
\longrightarrow
\textit{Mod}^{gr}(\mathcal{A})
$$ adjoint à gauche à
$j^*$. À savoir, étant donné un $\mathcal{A}_U$-module gradué
$\mathcal{M}$, nous définissons $j_!\mathcal{M}$ comme le $\mathcal{A}$-module gradué dont le terme de degré
$n$ est $j_!\mathcal{M}^n$. Comme application de multiplication, nous
utilisons
$$
j_!\mu_{n, m} :
j_!\mathcal{M}^n \times \mathcal{A}^m \to
j_!\mathcal{M}^{n + m}
$$ où
$\mu_{m, n} : \mathcal{M}^n \times \mathcal{A}^m \to \mathcal{M}^{n + m}$ est l'application
de multiplication donnée. Étant donnée une application homogène
$f : \mathcal{M} \to \mathcal{M}'$ de degré $n$ de
$\mathcal{A}_U$-modules gradués, nous obtenons une application homogène
$j_!f : j_!\mathcal{M} \to j_!\mathcal{M}'$ de degré $n$. Ainsi
nous obtenons notre foncteur.

\begin{lemma}
\label{lemma-extension-by-zero-graded} Dans
la situation ci-dessus nous
avons $$
\Hom_{\textit{Mod}^{gr}(\mathcal{A})}(
j_!\mathcal{M}, \mathcal{N})
= \Hom_{\textit{Mod}^{gr}(\mathcal{A}_U)}(
\mathcal{M}, j^*\mathcal{N})
$$
\end{lemma}

\begin{proof}
D'après la discussion
dans Modules sur les sites, section \ref{sites-modules-section-localize},
les foncteurs $j_!$ et $j^*$ sur les $\mathcal{O}$-modules sont
adjoints. Ainsi, en ne considérant que les structures de $\mathcal{O}$-modules,
nous
savons que $$
\Hom_{\text{Gr}^{gr}(\textit{Mod}(\mathcal{O}))}(
j_!\mathcal{M}, \mathcal{N})
= \Hom_{\text{Gr}^{gr}(\textit{Mod}(\mathcal{O}_U))}(
\mathcal{M}, j^*\mathcal{N})
$$
(Rappelons que $\text{Gr}^{gr}(\textit{Mod}(\mathcal{O}))$
désigne la catégorie graduée des $\mathcal{O}$-modules
gradués.) Il faut ensuite vérifier que ces identifications
envoient les homomorphismes de $\mathcal{A}$-modules du côté
gauche vers les homomorphismes de $\mathcal{A}_U$-modules du
côté droit. Pour le vérifier, étant donnés des morphismes $\mathcal{O}_U$-linéaires
$f^n : \mathcal{M}^n \to j^*\mathcal{N}^{n + d}$
correspondant à des morphismes $\mathcal{O}$-linéaires
$g^n : j_!\mathcal{M}^n \to \mathcal{N}^{n + d}$,
il
suffit de montrer que $$
\xymatrix{
\mathcal{M}^n \otimes_{\mathcal{O}_U} \mathcal{A}_U^m
\ar[r]_{f^n \otimes 1} \ar[d]
& j^*\mathcal{N}^{n + d} \otimes_{\mathcal{O}_U} \mathcal{A}_U^m \ar[d] \\
\mathcal{M}^{n + m} \ar[r]^{f^{n + m}}
& j^*\mathcal{N}^{n + m + d}
}
$$
commute
si et seulement si $$
\xymatrix{
j_!\mathcal{M}^n \otimes_\mathcal{O} \mathcal{A}^m
\ar[r]_{g^n \otimes 1} \ar[d]
& \mathcal{N}^{n + d} \otimes_\mathcal{O} \mathcal{A}_U^m \ar[d] \\
j_!\mathcal{M}^{n + m} \ar[r]^{g^{n + m}}
& \mathcal{N}^{n + m + d}
}
$$
est commutatif.
Cependant, nous savons que \begin{align*}
\Hom_{\mathcal{O}_U}(\mathcal{M}^n \otimes_{\mathcal{O}_U} \mathcal{A}_U^m,
j^*\mathcal{N}^{n + d + m})
&
= \Hom_\mathcal{O}(j_!(\mathcal{M}^n \otimes_{\mathcal{O}_U} \mathcal{A}_U^m),
\mathcal{N}^{n + d + m}) \\
&
= \Hom_\mathcal{O}(j_!\mathcal{M}^n \otimes_\mathcal{O} \mathcal{A}^m,
\mathcal{N}^{n + d + m})
\end{align*}
d'après Modules
sur les sites, lemme \ref{sites-modules-lemma-j-shriek-and-tensor}, déjà
utilisé. Nous omettons la vérification montrant que l'obstruction à la
commutativité du premier diagramme dans le premier groupe est
envoyée sur l'obstruction à la commutativité du second diagramme dans
le dernier groupe.
\end{proof}

\begin{lemma}
\label{lemma-tensor-with-extension-by-zero} Dans
la situation ci-dessus, soit $\mathcal{M}$ un
$\mathcal{A}_U$-module gradué à droite et soit $\mathcal{N}$ un
$\mathcal{A}$-module gradué à gauche. Alors
$$
j_!\mathcal{M} \otimes_\mathcal{A} \mathcal{N} =
j_!(\mathcal{M} \otimes_{\mathcal{A}_U} \mathcal{N}|_U)
$$ en tant
que $\mathcal{O}$-modules gradués, fonctoriellement en $\mathcal{M}$
et $\mathcal{N}$.
\end{lemma}

\begin{proof}
Rappelons que la composante de degré $n$ de
$j_!\mathcal{M} \otimes_\mathcal{A} \mathcal{N}$ est le conoyau de
l'application canonique
$$
\bigoplus\nolimits_{r + s + t = n}
j_!\mathcal{M}^r \otimes_\mathcal{O}
\mathcal{A}^s \otimes_\mathcal{O}
\mathcal{N}^t
\longrightarrow
\bigoplus\nolimits_{p + q = n}
j_!\mathcal{M}^p \otimes_\mathcal{O} \mathcal{N}^q
$$ Voir
la section \ref{section-tensor-product}.
Par Modules sur les sites, lemme \ref{sites-modules-lemma-j-shriek-and-tensor}
ceci est la même chose que le conoyau de
$$
\bigoplus\nolimits_{r + s + t = n}
j_!(\mathcal{M}^r \otimes_{\mathcal{O}_U}
\mathcal{A}^s|_U \otimes_{\mathcal{O}_U}
\mathcal{N}^t|_U)
\longrightarrow
\bigoplus\nolimits_{p + q = n}
j_!(\mathcal{M}^p \otimes_{\mathcal{O}_U} \mathcal{N}^q|_U)
$$ ce qui
conclut. Une preuve alternative serait de refaire l'argument de Yoneda
donné dans la démonstration du lemme cité ci-dessus.
\end{proof}








\section{Foncteurs de décalage sur les faisceaux de modules gradués}
\label{section-shift}

\noindent
Nous recommandons vivement au lecteur d'omettre cette section. Il s'avère que les
faisceaux de modules gradués sur une algèbre graduée sont un exemple du phénomène
discuté dans
Algèbre différentielle graduée, remarque \ref{dga-remark-graded-shift-functors}.

\medskip\noindent
Soit $(\mathcal{C}, \mathcal{O})$ un site annelé. Soit
$\mathcal{A}$ un faisceau d'algèbres graduées sur
$(\mathcal{C}, \mathcal{O})$. Soit
$\mathcal{M}$ un $\mathcal{A}$-module gradué. Soit $k \in \mathbf{Z}$. Nous définissons
le {\it $k$-ième décalage de} $\mathcal{M}$, noté $\mathcal{M}[k]$, comme le
$\mathcal{A}$-module gradué dont la composante de degré $n$ est donnée par
$$
(\mathcal{M}[k])^n = \mathcal{M}^{n + k}
$$
c'est-à-dire la composante de degré $(n + k)$ de $\mathcal{M}$. Pour les
applications de multiplication $$
(\mathcal{M}[k])^n \times \mathcal{A}^m
\longrightarrow
(\mathcal{M}[k])^{n + m}
$$
nous utilisons
simplement les applications de multiplication $$
\mathcal{M}^{n + k} \times \mathcal{A}^m
\longrightarrow
\mathcal{M}^{n + m + k}
$$
de $\mathcal{M}$. Il est clair que nous avons défini un foncteur $[k]$,
que nous avons $[k + l] = [k] \circ [l]$,
et que nous avons $$
\Hom_{\textit{Mod}^{gr}(\mathcal{A})}(\mathcal{L}, \mathcal{M}[k])
= \Hom_{\textit{Mod}^{gr}(\mathcal{A})}(\mathcal{L}, \mathcal{M})[k]
$$
sans intervention
de signes, fonctoriellement en $\mathcal{M}$ et $\mathcal{L}$.
Ainsi, la catégorie graduée des $\mathcal{A}$-modules gradués se retrouve
bien à partir de la catégorie ordinaire des $\mathcal{A}$-modules gradués et
des foncteurs de décalage, comme expliqué dans
Algèbre différentielle graduée, remarque \ref{dga-remark-graded-shift-functors}.

\begin{lemma}
\label{lemma-gm-grothendieck-abelian}
Soit $(\mathcal{C}, \mathcal{O})$ un site annelé. Soit
$\mathcal{A}$ une $\mathcal{O}$-algèbre graduée. La
catégorie $\textit{Mod}(\mathcal{A})$ est une catégorie abélienne de Grothendieck.
\end{lemma}

\begin{proof}
D'après le lemme \ref{lemma-gm-abelian} et la définition d'une catégorie
abélienne de Grothendieck
(Injectifs, définition \ref{injectives-definition-grothendieck-conditions}),
il suffit de
montrer que $\textit{Mod}(\mathcal{A})$ possède un générateur. Nous affirmons
que $$
\mathcal{G} = \bigoplus\nolimits_{k, U} j_{U!}\mathcal{A}_U[k]
$$
est un générateur où la somme est effectuée sur tous les objets $U$ de $\mathcal{C}$
et $k \in \mathbf{Z}$. En effet, étant donné un $\mathcal{A}$-module
gradué $\mathcal{M}$, s'il n'existe aucun morphisme non nul de $\mathcal{G}$ vers $\mathcal{M}$,
alors, pour tout $k$ et tout $U$,
on a $$
\Hom_{\textit{Mod}(\mathcal{A})}(j_{U!}\mathcal{A}_U[k], \mathcal{M})
= \Hom_{\textit{Mod}(\mathcal{A}_U)}(\mathcal{A}_U[k], \mathcal{M}|_U) =
\Gamma(U, \mathcal{M}^{-k})
$$ et
ce groupe est nul. Ainsi, $\mathcal{M}$ est nul.
\end{proof}









\section{Faisceaux d'algèbres différentielles graduées}
\label{section-dga}

\noindent
Cette section est l'analogue
d'Algèbre différentielle graduée, section \ref{dga-section-dga}.

\begin{definition}
\label{definition-dga} Soit
$(\mathcal{C}, \mathcal{O})$ un site annelé. Un
{\it faisceau de $\mathcal{O}$-algèbres différentielles graduées}, ou
un {\it faisceau d'algèbres différentielles graduées},
sur $(\mathcal{C}, \mathcal{O})$ est un complexe de
cochaînes $\mathcal{A}^\bullet$ de $\mathcal{O}$-modules
muni d'applications $\mathcal{O}$-bilinéaires
$$
\mathcal{A}^n \times \mathcal{A}^m \to \mathcal{A}^{n + m},\quad
(a, b) \longmapsto ab
$$
appelées applications de multiplication et vérifiant les propriétés
suivantes : \begin{enumerate}
\item la multiplication est
associative, \item il existe une section globale $1$ de $\mathcal{A}^0$
qui est une unité bilatère pour la
multiplication, \item pour $U \in \Ob(\mathcal{C})$, $a \in \mathcal{A}^n(U)$
et $b \in \mathcal{A}^m(U)$, on a
$$
\text{d}^{n + m}(ab) = \text{d}^n(a)b + (-1)^n a\text{d}^m(b)
$$
\end{enumerate} Nous
notons souvent une telle structure $(\mathcal{A}, \text{d})$. Un
{\it homomorphisme d'algèbres différentielles graduées sur $\mathcal{O}$}
de $(\mathcal{A}, \text{d})$ vers $(\mathcal{B}, \text{d})$ est un morphisme
$f : \mathcal{A}^\bullet \to \mathcal{B}^\bullet$ de complexes de
$\mathcal{O}$-modules compatible avec les applications de multiplication.
\end{definition}

\noindent
Si l'on donne une $\mathcal{O}$-algèbre différentielle graduée $(\mathcal{A}, \text{d})$
et un objet $U \in \Ob(\mathcal{C})$, on utilise la notation
$$
\mathcal{A}(U)
= \Gamma(U, \mathcal{A}) =
\bigoplus\nolimits_{n \in \mathbf{Z}} \mathcal{A}^n(U)
$$ Il
s'agit d'une $\mathcal{O}(U)$-algèbre différentielle graduée.

\medskip\noindent
Dans la mesure du possible, nous considérerons une
$\mathcal{O}$-algèbre différentielle graduée $(\mathcal{A}, \text{d})$
comme une $\mathcal{O}$-algèbre graduée $\mathcal{A}$ munie
de l'opérateur $\text{d} : \mathcal{A} \to \mathcal{A}$ de degré $1$
(où $\mathcal{A}$ est considérée comme un $\mathcal{O}$-module gradué),
satisfaisant à la règle de Leibniz de la définition.

\begin{remark}
\label{remark-functoriality-dga} Soit
$(f, f^\sharp) : (\Sh(\mathcal{C}), \mathcal{O}_\mathcal{C})
\to (\Sh(\mathcal{D}), \mathcal{O}_\mathcal{D})$ un morphisme
de topos annelés.
\begin{enumerate}
\item Soit $(\mathcal{A}, \text{d})$ une
$\mathcal{O}_\mathcal{C}$-algèbre différentielle graduée. L'image directe est la
$\mathcal{O}_\mathcal{D}$-algèbre différentielle graduée
$(f_*\mathcal{A}, \text{d})$, où $f_*\mathcal{A}$ est comme dans la remarque
\ref{remark-functoriality-ga} et
$\text{d} = f_*\text{d}$ comme morphisme $f_*\mathcal{A}^n \to f_*\mathcal{A}^{n + 1}$. Nous omettons
la vérification que la règle de Leibniz est satisfaite.
\item Soit $\mathcal{B}$ une
$\mathcal{O}_\mathcal{D}$-algèbre différentielle graduée. L'image inverse est la
$\mathcal{O}_\mathcal{C}$-algèbre différentielle graduée
$(f^*\mathcal{B}, \text{d})$, où $f^*\mathcal{B}$ est comme dans la remarque
\ref{remark-functoriality-ga} et
$\text{d} = f^*\text{d}$ comme morphisme $f^*\mathcal{B}^n \to f^*\mathcal{B}^{n + 1}$. Nous omettons
la vérification que la règle de Leibniz est satisfaite.
\item L'ensemble des homomorphismes $f^*\mathcal{B} \to \mathcal{A}$ de
$\mathcal{O}_\mathcal{C}$-algèbres différentielles graduées est en
correspondance biunivoque ($1$ à $1$) avec l'ensemble des homomorphismes
$\mathcal{B} \to f_*\mathcal{A}$ de
$\mathcal{O}_\mathcal{D}$-algèbres différentielles graduées.
\end{enumerate} La
partie (3) découle immédiatement de l'adjonction usuelle entre $f^*$ et $f_*$ sur les
faisceaux de modules.
\end{remark}
















\section{Faisceaux de modules différentiels gradués}
\label{section-modules}

\noindent
Cette section est l'analogue de
l'Algèbre différentielle graduée, section \ref{dga-section-modules}.

\begin{definition}
\label{definition-dgm}
Soit $(\mathcal{C}, \mathcal{O})$ un site annelé. Soit
$(\mathcal{A}, \text{d})$ un faisceau d'algèbres différentielles graduées sur
$(\mathcal{C}, \mathcal{O})$. Un
{\it $\mathcal{A}$-module différentiel gradué à droite}, ou
{\it module différentiel gradué à droite} sur $\mathcal{A}$, est un
complexe de cochaînes $\mathcal{M}^\bullet$ muni d'applications
$\mathcal{O}$-bilinéaires
$$
\mathcal{M}^n \times \mathcal{A}^m \to \mathcal{M}^{n + m},\quad (x,
a) \longmapsto xa
$$
appelées applications de multiplication, satisfaisant aux propriétés
suivantes : \begin{enumerate}
\item la multiplication vérifie $(xa)a' = x(aa')$,
\item la section unité $1$ de $\mathcal{A}^0$
agit comme l'identité sur $\mathcal{M}^n$ pour tout $n$,
\item pour $U \in \Ob(\mathcal{C})$, $x \in \mathcal{M}^n(U)$ et
$a \in \mathcal{A}^m(U)$, on a
$$
\text{d}^{n + m}(xa) = \text{d}^n(x)a + (-1)^n x\text{d}^m(a)
$$
\end{enumerate} Nous
disons souvent « soit $\mathcal{M}$ un
$\mathcal{A}$-module différentiel gradué » pour indiquer cette situation. Un
{\it homomorphisme de $\mathcal{A}$-modules différentiels gradués} de
$\mathcal{M}$ vers $\mathcal{N}$ est un morphisme
$f : \mathcal{M}^\bullet \to \mathcal{N}^\bullet$ de complexes de
$\mathcal{O}$-modules compatible avec les applications de multiplication.
La catégorie des $\mathcal{A}$-modules différentiels gradués (à
droite) est notée $\textit{Mod}(\mathcal{A}, \text{d})$.
\end{definition}

\noindent
On définit de même les {\it modules différentiels gradués à gauche}, mais, sauf mention
contraire, les modules de ce chapitre seront des modules à droite.

\medskip\noindent
Étant donnés un $\mathcal{A}$-module différentiel gradué $\mathcal{M}$
et un objet $U \in \Ob(\mathcal{C})$, nous utilisons la notation
$$
\mathcal{M}(U)
= \Gamma(U, \mathcal{M})
= \bigoplus\nolimits_{n \in \mathbf{Z}} \mathcal{M}^n(U)
$$
Il s'agit d'un $\mathcal{A}(U)$-module différentiel gradué à droite.

\begin{lemma}
\label{lemma-dgm-abelian} Soit
$(\mathcal{C}, \mathcal{O})$ un site annelé. Soit
$(\mathcal{A}, \text{d})$ une $\mathcal{O}$-algèbre différentielle graduée. La catégorie
$\textit{Mod}(\mathcal{A}, \text{d})$ est une catégorie abélienne qui possède les
propriétés suivantes :
\begin{enumerate}
\item $\textit{Mod}(\mathcal{A}, \text{d})$ possède des sommes directes arbitraires,
\item $\textit{Mod}(\mathcal{A}, \text{d})$ possède des limites inductives arbitraires,
\item les limites inductives filtrantes dans $\textit{Mod}(\mathcal{A}, \text{d})$ sont exactes,
\item $\textit{Mod}(\mathcal{A}, \text{d})$ possède des produits
arbitraires, \item $\textit{Mod}(\mathcal{A}, \text{d})$ possède des limites
arbitraires. \end{enumerate}
Le foncteur
d'oubli $$
\textit{Mod}(\mathcal{A}, \text{d})
\longrightarrow
\textit{Mod}(\mathcal{A})
$$
qui envoie un $\mathcal{A}$-module différentiel gradué sur son module
gradué sous-jacent commute à toutes les limites et limites inductives.
\end{lemma}

\begin{proof}
Notons
$F : \textit{Mod}(\mathcal{A}, \text{d}) \to \textit{Mod}(\mathcal{A})$
le foncteur de l'énoncé. La catégorie
$\textit{Mod}(\mathcal{A})$ possède les propriétés (1) -- (5) ; voir le lemme
\ref{lemma-gm-abelian}.

\medskip\noindent
Considérons un homomorphisme $f : \mathcal{M} \to \mathcal{N}$
de $\mathcal{A}$-modules gradués. Le noyau et
le conoyau de $f$, comme morphisme de $\mathcal{A}$-modules
gradués, sont en outre munis de différentiels comme dans
la définition \ref{definition-dgm}. Ce sont donc aussi le
noyau et le conoyau dans $\textit{Mod}(\mathcal{A}, \text{d})$. Ainsi,
$\textit{Mod}(\mathcal{A}, \text{d})$ est une catégorie abélienne, et
la formation des noyaux et conoyaux commute avec $F$.

\medskip\noindent
Pour prouver l'existence de limites et de limites inductives, il suffit de
prouver l'existence de produits et de sommes directes, voir
Catégories, lemmes \ref{categories-lemma-limits-products-equalizers} et
\ref{categories-lemma-colimits-coproducts-coequalizers}. Les mêmes lemmes
montrent que, pour prouver
que $F$ commute aux limites et aux limites inductives, il suffit de
prouver que $F$ commute aux sommes directes et aux produits.

\medskip\noindent
Soit $\mathcal{M}_t$, $t \in T$, une famille de
$\mathcal{A}$-modules différentiels gradués. On peut considérer la somme directe
$\bigoplus \mathcal{M}_t$ comme un $\mathcal{A}$-module gradué.
Puisque la somme directe des modules gradués se calcule terme à terme,
$\bigoplus \mathcal{M}_t$ est naturellement muni d'un différentiel. Le lecteur
vérifie facilement que c'est une somme directe en
$\textit{Mod}(\mathcal{A}, \text{d})$. Il en va de même pour les produits.

\medskip\noindent
Observons que $F$ est exact et qu'un complexe
$\mathcal{M}_1 \to \mathcal{M}_2 \to \mathcal{M}_3$ de
$\textit{Mod}(\mathcal{A}, \text{d})$ est exact si et seulement si
$F(\mathcal{M}_1) \to F(\mathcal{M}_2) \to F(\mathcal{M}_3)$ est
exact dans $\textit{Mod}(\mathcal{A})$. On en déduit (3), puisque les
limites inductives filtrantes sont exactes dans
$\textit{Mod}(\mathcal{A})$.
\end{proof}

\noindent
En combinant les lemmes \ref{lemma-dgm-abelian} et \ref{lemma-gm-abelian},
on obtient un foncteur exact et fidèle
$$
\textit{Mod}(\mathcal{A}, \text{d}) \longrightarrow \text{Comp}(\mathcal{O})
$$ entre catégories
abéliennes. Pour un $\mathcal{A}$-module différentiel gradué
$\mathcal{M}$, les $\mathcal{O}$-modules de cohomologie, notés $H^i(\mathcal{M})$, sont
définis comme les cohomologies du complexe sous-jacent
de $\mathcal{O}$-modules associé à $\mathcal{M}$. Par
conséquent, une suite exacte courte
$0 \to \mathcal{K} \to \mathcal{L} \to \mathcal{M} \to 0$ de
$\mathcal{A}$-modules différentiels gradués
donne lieu à une suite exacte longue
\begin{equation}
\label{equation-les}
H^n(\mathcal{K}) \to H^n(\mathcal{L}) \to H^n(\mathcal{M}) \to
H^{n + 1}(\mathcal{K})
\end{equation} des
modules de cohomologie, voir
Homologie, lemme \ref{homology-lemma-long-exact-sequence-cochain}.

\medskip\noindent En
outre, à partir de maintenant, nous empruntons toute la terminologie
utilisée pour les complexes de modules. Par exemple, nous disons
qu'un $\mathcal{A}$-module différentiel gradué $\mathcal{M}$ est {\it acyclique}
si $H^k(\mathcal{M}) = 0$ pour tous les $k \in \mathbf{Z}$.
Un homomorphisme $\mathcal{M} \to \mathcal{N}$
de $\mathcal{A}$-modules différentiels gradués est un {\it quasi-isomorphisme}
s'il induit des isomorphismes $H^k(\mathcal{M}) \to H^k(\mathcal{N})$
pour tout $k \in \mathbf{Z}$. Et ainsi de suite.


















\section{La catégorie différentielle graduée des modules}
\label{section-dgm-dg-cat}

\noindent
Cette section est l'analogue
d'Algèbre différentielle graduée, exemple \ref{dga-example-dgm-dg-cat}.
Pour nos conventions sur les catégories différentielles graduées,
voir Algèbre différentielle graduée, section \ref{dga-section-dga-categories}.

\medskip\noindent
Soit $(\mathcal{C}, \mathcal{O})$ un site annelé. Soit
$(\mathcal{A}, \text{d})$ un faisceau d'algèbres différentielles graduées sur
$(\mathcal{C}, \mathcal{O})$. Nous allons construire une
catégorie différentielle graduée
$$
\textit{Mod}^{dg}(\mathcal{A}, \text{d})
$$
sur $R = \Gamma(\mathcal{C}, \mathcal{O})$, dont la catégorie de complexes
associée est la catégorie des $\mathcal{A}$-modules différentiels gradués :
$$
\textit{Mod}(\mathcal{A}, \text{d}) =
\text{Comp}(\textit{Mod}^{dg}(\mathcal{A}, \text{d}))
$$ Comme objets
de $\textit{Mod}^{dg}(\mathcal{A}, \text{d})$, nous prenons les
$\mathcal{A}$-modules différentiels gradués à droite, voir
la section \ref{section-modules}. Étant donnés des
$\mathcal{A}$-modules différentiels gradués $\mathcal{L}$ et $\mathcal{M}$,
posons $$
\Hom_{\textit{Mod}^{dg}(\mathcal{A}, \text{d})}(\mathcal{L}, \mathcal{M})
= \Hom_{\textit{Mod}^{gr}(\mathcal{A})}(\mathcal{L}, \mathcal{M}) =
\bigoplus\nolimits_{n \in \mathbf{Z}} \Hom^n(\mathcal{L}, \mathcal{M})
$$ comme
$R$-module gradué, voir la section \ref{section-gm-gr-cat}. Autrement
dit, la composante graduée de degré $n$,
$\Hom^n(\mathcal{L}, \mathcal{M})$, est le $R$-module des morphismes
de $\mathcal{A}$-modules à droite homogènes de degré $n$.
Pour $f \in \Hom^n(\mathcal{L}, \mathcal{M})$, on pose
$$
\text{d}(f) =
\text{d}_\mathcal{M} \circ f - (-1)^n f \circ \text{d}_\mathcal{L}
$$ Pour donner
un sens à cette formule, on considère $\text{d}_\mathcal{M}$ et
$\text{d}_\mathcal{L}$ comme des morphismes gradués de $\mathcal{O}$-modules
et on utilise la composition de morphismes gradués de $\mathcal{O}$-modules.
Il est clair que $\text{d}(f)$
est homogène de degré $n + 1$ comme morphisme gradué de $\mathcal{O}$-modules,
et qu'il est $\mathcal{A}$-linéaire, car, pour des sections locales homogènes $x$ et $a$
de $\mathcal{M}$ et $\mathcal{A}$,
on a \begin{align*}
\text{d}(f)(xa)
&
= \text{d}_\mathcal{M}(f(x) a) - (-1)^n f (\text{d}_\mathcal{L}(xa)) \\
&
= \text{d}_\mathcal{M}(f(x)) a + (-1)^{\deg(x) + n} f(x) \text{d}(a)
- (-1)^n f(\text{d}_\mathcal{L}(x)) a - (-1)^{n + \deg(x)} f(x) \text{d}(a) \\
& = \text{d}(f)(x)
a \end{align*}
comme souhaité (observons que ce calcul ne fonctionnerait pas sans le
signe dans la définition de notre différentiel sur $\Hom$).

\medskip\noindent
Pour des $\mathcal{A}$-modules différentiels
gradués $\mathcal{K}$, $\mathcal{L}$ et $\mathcal{M}$,
nous avons déjà défini la composition
$$
\Hom^m(\mathcal{L}, \mathcal{M}) \times
\Hom^n(\mathcal{K}, \mathcal{L}) \longrightarrow
\Hom^{n + m}(\mathcal{K}, \mathcal{M})
$$
dans la section \ref{section-gm-gr-cat} par composition usuelle
de morphismes de faisceaux. Cela définit un morphisme de modules différentiels
gradués $$
\Hom_{\textit{Mod}^{dg}(\mathcal{A}, \text{d})}(\mathcal{L}, \mathcal{M})
\otimes_R
\Hom_{\textit{Mod}^{dg}(\mathcal{A}, \text{d})}(\mathcal{K}, \mathcal{L})
\longrightarrow
\Hom_{\textit{Mod}^{dg}(\mathcal{A}, \text{d})}(\mathcal{K}, \mathcal{M})
$$
comme
l'exige l'Algèbre différentielle graduée, définition \ref{dga-definition-dga-category},
comme le montre le calcul suivant :
\begin{align*}
\text{d}(g \circ
f) & = \text{d}_\mathcal{M} \circ g \circ
f - (-1)^{n + m} g \circ f \circ \text{d}_\mathcal{K} \\
&
= \left(\text{d}_\mathcal{M} \circ g - (-1)^m g \circ \text{d}_L\right) \circ
f + (-1)^m g \circ \left(\text{d}_\mathcal{L} \circ
f - (-1)^n f \circ \text{d}_\mathcal{K}\right) \\
& =
\text{d}(g) \circ f + (-1)^m g \circ \text{d}(f)
\end{align*}
si $f$ est de degré $n$ et $g$ de degré $m$, comme souhaité.






\section{Produit tensoriel de faisceaux de modules différentiels gradués}
\label{section-tensor-product-dg}

\noindent
Cette section est l'analogue d'une partie de
l'Algèbre différentielle graduée, section \ref{dga-section-tensor-product}.

\medskip\noindent
Soit $(\mathcal{C}, \mathcal{O})$ un site annelé. Soit
$(\mathcal{A}, \text{d})$ un
faisceau d'algèbres différentielles graduées sur $(\mathcal{C}, \mathcal{O})$.
Soient $\mathcal{M}$ un $\mathcal{A}$-module différentiel gradué à
droite et $\mathcal{N}$ un $\mathcal{A}$-module différentiel gradué à gauche.
Dans cette situation, nous définissons le {\it produit tensoriel}
$\mathcal{M} \otimes_\mathcal{A} \mathcal{N}$ comme
suit. Comme $\mathcal{O}$-module gradué, il est donné par la
construction de la section \ref{section-tensor-product}. Il est
doté d'un différentiel
$$
\text{d}_{\mathcal{M} \otimes_\mathcal{A} \mathcal{N}}
: (\mathcal{M} \otimes_\mathcal{A} \mathcal{N})^n
\longrightarrow
(\mathcal{M} \otimes_\mathcal{A} \mathcal{N})^{n + 1}
$$
défini par la règle
$$
\text{d}_{\mathcal{M} \otimes_\mathcal{A} \mathcal{N}}(x \otimes y) =
\text{d}_\mathcal{M}(x) \otimes y +
(-1)^{\deg(x)}x \otimes \text{d}_\mathcal{N}(y)
$$ pour
les sections locales homogènes $x$ et $y$ de $\mathcal{M}$ et $\mathcal{N}$. Pour
vérifier que cette définition est correcte, montrons
que $\text{d}_{\mathcal{M} \otimes_\mathcal{A} \mathcal{N}}$ annule les
éléments de la forme $xa \otimes y - x \otimes ay$, pour des sections
locales homogènes $x$, $a$, $y$ de $\mathcal{M}$, $\mathcal{A}$, $\mathcal{N}$.
Nous
calculons \begin{align*}
&
\text{d}_{\mathcal{M} \otimes_\mathcal{A} \mathcal{N}}(
xa \otimes y - x \otimes ay) \\
&
= \text{d}_\mathcal{M}(xa) \otimes y + (-1)^{\deg(x) + \deg(a)}
xa \otimes \text{d}_\mathcal{N}(y)
-\text{d}_\mathcal{M}(x) \otimes ay - (-1)^{\deg(x)}
x \otimes \text{d}_\mathcal{N}(ay) \\
&
= \text{d}_\mathcal{M}(x)a \otimes y + (-1)^{\deg(x)}x\text{d}(a) \otimes
y + (-1)^{\deg(x) + \deg(a)} xa \otimes \text{d}_\mathcal{N}(y) \\
&
-\text{d}_\mathcal{M}(x) \otimes ay - (-1)^{\deg(x)}
x \otimes \text{d}(a)y - (-1)^{\deg(x) +
\deg(a)} x\otimes a\text{d}_\mathcal{N}(y)
\end{align*}
Nous observons que les éléments
$$
\text{d}_\mathcal{M}(x)a \otimes y - \text{d}_\mathcal{M}(x) \otimes ay,\quad
x\text{d}(a) \otimes y - x \otimes \text{d}(a)y,\quad
\text{et}\quad xa
\otimes \text{d}_\mathcal{N}(y) - x\otimes a\text{d}_\mathcal{N}(y)
$$ ont une
image nulle dans $\mathcal{M} \otimes_\mathcal{A} \mathcal{N}$,
ce qui conclut. Nous omettons la vérification
que $\text{d}_{\mathcal{M} \otimes_\mathcal{A} \mathcal{N}} \circ
\text{d}_{\mathcal{M} \otimes_\mathcal{A} \mathcal{N}} = 0$.

\medskip\noindent
Si nous fixons le $\mathcal{A}$-module différentiel gradué à gauche $\mathcal{N}$,
nous obtenons un
foncteur $$
- \otimes_\mathcal{A} \mathcal{N}
: \textit{Mod}(\mathcal{A}, \text{d})
\longrightarrow
\text{Comp}(\mathcal{O})
$$
où le membre de droite est la catégorie des complexes de
$\mathcal{O}$-modules. Ce foncteur se relève en un foncteur de
catégories différentielles graduées
$$
- \otimes_\mathcal{A} \mathcal{N} :
\textit{Mod}^{dg}(\mathcal{A}, \text{d})
\longrightarrow
\text{Comp}^{dg}(\mathcal{O})
$$ Sur
les objets gradués sous-jacents, nous
envoyons un homomorphisme $f : \mathcal{M} \to \mathcal{M}'$ de degré $n$
sur l'application de degré $n$,
$f \otimes \text{id}_\mathcal{N}
: \mathcal{M} \otimes_\mathcal{A} \mathcal{N} \to
\mathcal{M}' \otimes_\mathcal{A} \mathcal{N}$,
conformément à la construction de la section \ref{section-tensor-product}.
Pour montrer que cela fonctionne, nous devons
vérifier que l'application $$
\Hom_{\textit{Mod}^{dg}(\mathcal{A}, \text{d})}(\mathcal{M}, \mathcal{M}')
\longrightarrow
\Hom_{\text{Comp}^{dg}(\mathcal{O})}(
\mathcal{M} \otimes_\mathcal{A} \mathcal{N},
\mathcal{M}' \otimes_\mathcal{A} \mathcal{N})
$$
est compatible avec les différentiels. Pour voir cela pour $f$
comme
ci-dessus nous devons montrer que $$
(\text{d}_{\mathcal{M}'} \circ f - (-1)^n f \circ \text{d}_\mathcal{M})
\otimes \text{id}_\mathcal{N}
$$
est
égal à $$
\text{d}_{\mathcal{M}' \otimes_\mathcal{A} \mathcal{N}}
\circ (f \otimes \text{id}_\mathcal{N})
- (-1)^n (f \otimes \text{id}_\mathcal{N}) \circ
\text{d}_{\mathcal{M} \otimes_\mathcal{A} \mathcal{N}}
$$
Calculons donc l'effet de ces opérateurs sur
une section locale de la forme $x \otimes y$, où $x$ et $y$ sont des sections
locales homogènes de $\mathcal{M}$ et $\mathcal{N}$. Pour le
premier, nous obtenons $$
(\text{d}_{\mathcal{M}'}(f(x)) - (-1)^n f(\text{d}_\mathcal{M}(x))) \otimes
y $$
et pour
le second nous obtenons \begin{align*}
&\text{d}_{\mathcal{M}' \otimes_\mathcal{A} \mathcal{N}}(f(x) \otimes
y) - (-1)^n (f \otimes \text{id}_\mathcal{N})(
\text{d}_{\mathcal{M} \otimes_\mathcal{A} \mathcal{N}}(x \otimes y) \\
&
= \text{d}_{\mathcal{M}'}(f(x)) \otimes
y + (-1)^{\deg(x) + n}f(x) \otimes \text{d}_\mathcal{N}(y) \\
&
-(-1)^n f(\text{d}_\mathcal{M}(x)) \otimes
y -(-1)^n (-1)^{\deg(x)}f(x) \otimes \text{d}_\mathcal{N}(y)
\end{align*}
qui est bien la même section locale.










\section{Hom interne de faisceaux de modules différentiels gradués}
\label{section-internal-hom-dg}

\noindent
Nous aurons besoin de la version faisceautisée de la construction de la
section \ref{section-dgm-dg-cat}. Soient
$(\mathcal{C}, \mathcal{O})$, $\mathcal{A}$,
$\mathcal{M}$ et $\mathcal{L}$ comme dans la section \ref{section-dgm-dg-cat}. Nous
définissons
$$
\SheafHom^{dg}_\mathcal{A}(\mathcal{M}, \mathcal{L}) =
\SheafHom^{gr}_\mathcal{A}(\mathcal{M}, \mathcal{L}) =
\bigoplus\nolimits_{n \in \mathbf{Z}}
\SheafHom^n_\mathcal{A}(\mathcal{M}, \mathcal{L})
$$ comme
$\mathcal{O}$-module gradué, voir la section
\ref{section-internal-hom-graded}. Autrement dit, une section
$f$ de la composante graduée de degré $n$,
$\SheafHom^n_\mathcal{A}(\mathcal{L}, \mathcal{M})$, sur $U$ est un morphisme de
$\mathcal{A}_U$-modules à droite
$\mathcal{L}|_U \to \mathcal{M}|_U$, homogène de degré $n$. Pour un tel
$f$, on pose
$$
\text{d}(f) =
\text{d}_\mathcal{M}|_U \circ f - (-1)^n f \circ \text{d}_\mathcal{L}|_U
$$ Pour donner un sens à cette formule,
on considère $\text{d}_\mathcal{M}|_U$ et
$\text{d}_\mathcal{L}|_U$ comme des morphismes gradués de $\mathcal{O}_U$-modules et on utilise la
composition de morphismes gradués de $\mathcal{O}_U$-modules. Il est
clair que $\text{d}(f)$ est homogène de degré
$n + 1$ comme morphisme gradué de $\mathcal{O}_U$-modules. En
utilisant exactement le même calcul que dans la section \ref{section-dgm-dg-cat},
nous voyons que $\text{d}(f)$ est $\mathcal{A}_U$-linéaire.

\medskip\noindent
Comme dans la section \ref{section-dgm-dg-cat}, il y a une application de
composition $$
\SheafHom^{dg}_\mathcal{A}(\mathcal{L}, \mathcal{M}) \otimes_\mathcal{O}
\SheafHom^{dg}_\mathcal{A}(\mathcal{K}, \mathcal{L})
\longrightarrow
\SheafHom^{dg}_\mathcal{A}(\mathcal{K}, \mathcal{M})
$$ où le
côté gauche est le produit tensoriel de
$\mathcal{O}$-modules différentiels gradués défini à la section
\ref{section-tensor-product-dg}. Cette application est donnée
par l'application de composition
$$
\SheafHom^m(\mathcal{L}, \mathcal{M}) \otimes_\mathcal{O}
\SheafHom^n(\mathcal{K}, \mathcal{L}) \longrightarrow
\SheafHom^{n + m}(\mathcal{K}, \mathcal{M})
$$ définie par simple
composition (localement). En appliquant exactement le même calcul
que dans la section \ref{section-dgm-dg-cat} aux sections locales, nous voyons
que l'application de composition est un morphisme de
$\mathcal{O}$-modules différentiels gradués.

\medskip\noindent
Avec ces définitions, nous avons
$$
\Hom_{\textit{Mod}^{dg}(\mathcal{A})}(\mathcal{L}, \mathcal{M}) =
\Gamma(\mathcal{C}, \SheafHom^{dg}_\mathcal{A}(\mathcal{L}, \mathcal{M}))
$$ comme
$R$-modules gradués, compatiblement à la composition.









\section{Faisceaux de bimodules différentiels gradués et adjonction tensor-Hom}
\label{section-dg-bimodules}

\noindent
Cette section est l'analogue d'une partie de
l'Algèbre différentielle graduée, section \ref{dga-section-tensor-product}.

\begin{definition}
\label{definition-dg-bimodule}
Soit $(\mathcal{C}, \mathcal{O})$ un site annelé. Soient $\mathcal{A}$ et
$\mathcal{B}$ des faisceaux d'algèbres différentielles graduées sur
$(\mathcal{C}, \mathcal{O})$. Un
{\it $(\mathcal{A}, \mathcal{B})$-bimodule différentiel gradué} est
la donnée d'un complexe $\mathcal{M}^\bullet$
de $\mathcal{O}$-modules muni d'applications $\mathcal{O}$-bilinéaires
$$
\mathcal{M}^n \times \mathcal{B}^m \to \mathcal{M}^{n + m},\quad
(x, b) \longmapsto xb
$$
et
$$
\mathcal{A}^n \times \mathcal{M}^m \to \mathcal{M}^{n + m},\quad (a,
x) \longmapsto ax
$$
appelées applications de multiplication, avec les propriétés suivantes
\begin{enumerate}
\item la multiplication satisfait $a(a'x) = (aa')x$
et $(xb)b' = x(bb')$,
\item $(ax)b = a(xb)$,
\item $\text{d}(ax) = \text{d}(a) x + (-1)^{\deg(a)}a \text{d}(x)$ et
$\text{d}(xb) = \text{d}(x) b + (-1)^{\deg(x)}x \text{d}(b)$,
\item la section unité $1$ de $\mathcal{A}^0$ agit comme l'identité par
multiplication, et
\item la section unité $1$ de
$\mathcal{B}^0$ agit comme l'identité par multiplication.
\end{enumerate} Nous désignons
souvent une telle structure $\mathcal{M}$ et parfois nous écrivons
${}_\mathcal{A}\mathcal{M}_\mathcal{B}$. Un
{\it homomorphisme de
$(\mathcal{A}, \mathcal{B})$-bimodules différentiels gradués}
$f : \mathcal{M} \to \mathcal{N}$ est un morphisme de complexes
$f : \mathcal{M}^\bullet \to \mathcal{N}^\bullet$ de
$\mathcal{O}$-modules compatible avec les applications de multiplication.
\end{definition}

\noindent
Étant donnés un $(\mathcal{A}, \mathcal{B})$-bimodule différentiel gradué $\mathcal{M}$
et un objet $U \in \Ob(\mathcal{C})$, nous utilisons la
notation $$
\mathcal{M}(U)
= \Gamma(U, \mathcal{M})
= \bigoplus\nolimits_{n \in \mathbf{Z}} \mathcal{M}^n(U)
$$
C'est un $(\mathcal{A}(U), \mathcal{B}(U))$-bimodule différentiel gradué.

\medskip\noindent
Observons qu'un $(\mathcal{A}, \mathcal{B})$-bimodule différentiel gradué
$\mathcal{M}$ revient à se donner un
$\mathcal{B}$-module différentiel gradué à droite qui est aussi un
$\mathcal{A}$-module différentiel gradué à gauche, de sorte que les graduations et les
différentiels coïncident et que les structures de $\mathcal{A}$-module
et de $\mathcal{B}$-module commutent. Voici un énoncé précis.

\begin{lemma}
\label{lemma-what-makes-a-bimodule-dg} Soit
$(\mathcal{C}, \mathcal{O})$ un site annelé. Soient $\mathcal{A}$ et
$\mathcal{B}$ des faisceaux d'algèbres différentielles graduées sur
$(\mathcal{C}, \mathcal{O})$. Soit $\mathcal{N}$ un
$\mathcal{B}$-module différentiel gradué à droite. Il existe une correspondance biunivoque ($1$ à $1$) entre les
structures de $(\mathcal{A}, \mathcal{B})$-bimodule sur
$\mathcal{N}$ compatibles avec la structure
donnée de $\mathcal{B}$-module différentiel gradué et les homomorphismes
$$
\mathcal{A}
\longrightarrow
\SheafHom^{dg}_\mathcal{B}(\mathcal{N}, \mathcal{N})
$$ de
$\mathcal{O}$-algèbres différentielles graduées.
\end{lemma}

\begin{proof}
Démonstration
omise. \end{proof}

\noindent
Soit $(\mathcal{C}, \mathcal{O})$ un site annelé. Soient $\mathcal{A}$ et
$\mathcal{B}$ des faisceaux d'algèbres différentielles graduées sur
$(\mathcal{C}, \mathcal{O})$. Soient $\mathcal{M}$ un
$\mathcal{A}$-module différentiel gradué à droite et $\mathcal{N}$ un
$(\mathcal{A}, \mathcal{B})$-bimodule différentiel gradué. Dans ce cas, le produit
tensoriel différentiel gradué défini à la section
\ref{section-tensor-product-dg}
$$
\mathcal{M} \otimes_\mathcal{A} \mathcal{N}
$$ est un
$\mathcal{B}$-module différentiel gradué à droite, avec les
applications de multiplication de la section \ref{section-graded-bimodules}.
Cette construction définit un foncteur et un foncteur de catégories
graduées $$
\otimes_\mathcal{A} \mathcal{N}
: \textit{Mod}(\mathcal{A}, \text{d})
\longrightarrow
\textit{Mod}(\mathcal{B}, \text{d})
\quad\text{et}\quad
\otimes_\mathcal{A} \mathcal{N}
: \textit{Mod}^{dg}(\mathcal{A}, \text{d})
\longrightarrow
\textit{Mod}^{dg}(\mathcal{B}, \text{d})
$$
en envoyant les homomorphismes de degré $n$ de $\mathcal{M} \to \mathcal{M}'$
sur l'application induite de degré $n$
de $\mathcal{M} \otimes_\mathcal{A} \mathcal{N}$ vers
$\mathcal{M}' \otimes_\mathcal{A} \mathcal{N}$.

\medskip\noindent
Soit $(\mathcal{C}, \mathcal{O})$ un site annelé. Soient $\mathcal{A}$ et
$\mathcal{B}$ des faisceaux d'algèbres différentielles graduées sur
$(\mathcal{C}, \mathcal{O})$. Soit $\mathcal{N}$ un
$(\mathcal{A}, \mathcal{B})$-bimodule différentiel gradué et soit
$\mathcal{L}$ un $\mathcal{B}$-module différentiel gradué à droite. Dans
ce cas, le Hom interne différentiel gradué défini à
la section \ref{section-internal-hom-dg}
$$
\SheafHom_\mathcal{B}^{dg}(\mathcal{N}, \mathcal{L})
$$ est un
$\mathcal{A}$-module différentiel gradué à droite dont la structure
graduée de $\mathcal{A}$-module à droite est celle de la section
\ref{section-graded-bimodules}. Une autre définition de la multiplication
consiste à utiliser la composition
$$
\SheafHom_\mathcal{B}^{dg}(\mathcal{N}, \mathcal{L})
\otimes_\mathcal{O} \mathcal{A}
\to
\SheafHom_\mathcal{B}^{dg}(\mathcal{N}, \mathcal{L})
\otimes_\mathcal{O} \SheafHom_\mathcal{B}^{dg}(\mathcal{N}, \mathcal{N})
\to
\SheafHom_\mathcal{B}^{dg}(\mathcal{N}, \mathcal{L})
$$ où la première
flèche provient du lemme \ref{lemma-what-makes-a-bimodule-dg} et la seconde est
l'application de composition de la section
\ref{section-internal-hom-dg}. Puisque ces deux flèches sont
compatibles avec les différentiels, nous concluons que nous obtenons
en effet un $\mathcal{A}$-module différentiel gradué.
Cette construction
définit un foncteur et un foncteur de catégories différentielles
graduées $$
\SheafHom_\mathcal{B}^{dg}(\mathcal{N},
-) : \textit{Mod}(\mathcal{B}, \text{d})
\longrightarrow
\textit{Mod}(\mathcal{A})
\quad\text{et}\quad
\SheafHom_\mathcal{B}^{dg}(\mathcal{N}, -)
: \textit{Mod}^{dg}(\mathcal{B}, \text{d})
\longrightarrow
\textit{Mod}^{dg}(\mathcal{A}, \text{d})
$$
en envoyant les homomorphismes de degré $n$ de $\mathcal{L} \to \mathcal{L}'$
sur l'application induite de degré $n$
de $\SheafHom_\mathcal{B}^{dg}(\mathcal{N}, \mathcal{L})$ vers
$\SheafHom_\mathcal{B}^{dg}(\mathcal{N}, \mathcal{L}')$.

\begin{lemma}
\label{lemma-tensor-hom-adjunction-dg}
Soit $(\mathcal{C}, \mathcal{O})$ un site annelé. Soient $\mathcal{A}$
et $\mathcal{B}$ des faisceaux d'algèbres différentielles graduées
sur $(\mathcal{C}, \mathcal{O})$. Soient $\mathcal{M}$ un
$\mathcal{A}$-module différentiel gradué à droite, $\mathcal{N}$ un
$(\mathcal{A}, \mathcal{B})$-bimodule différentiel gradué et $\mathcal{L}$ un
$\mathcal{B}$-module différentiel gradué à droite. Avec les conventions ci-dessus, nous
avons
$$
\Hom_{\textit{Mod}^{dg}(\mathcal{B}, \text{d})}(
\mathcal{M} \otimes_\mathcal{A} \mathcal{N}, \mathcal{L}) =
\Hom_{\textit{Mod}^{dg}(\mathcal{A}, \text{d})}(
\mathcal{M}, \SheafHom_\mathcal{B}^{dg}(\mathcal{N}, \mathcal{L}))
$$
et
$$
\SheafHom_\mathcal{B}^{dg}(
\mathcal{M} \otimes_\mathcal{A} \mathcal{N}, \mathcal{L}) =
\SheafHom_\mathcal{A}^{dg}(
\mathcal{M}, \SheafHom_\mathcal{B}^{dg}(\mathcal{N}, \mathcal{L}))
$$
fonctoriellement en $\mathcal{M}$, $\mathcal{N}$ et $\mathcal{L}$.
\end{lemma}

\begin{proof}
Omis. Indication : au niveau gradué, nous avons établi ce résultat
dans le lemme \ref{lemma-tensor-hom-adjunction-gr}. Il suffit
donc de vérifier que les isomorphismes sont compatibles avec les
différentiels, ce qui se fait par un calcul sur les sections locales.
\end{proof}

\noindent
Soit $(\mathcal{C}, \mathcal{O})$ un site annelé. Soient $\mathcal{A}$
et $\mathcal{B}$ des faisceaux d'algèbres différentielles graduées
sur $(\mathcal{C}, \mathcal{O})$.
Comme cas particulier de ce qui précède, supposons
donné un homomorphisme $\varphi : \mathcal{A} \to \mathcal{B}$
de $\mathcal{O}$-algèbres différentielles graduées. On obtient alors
un foncteur et un foncteur de catégories différentielles
graduées $$
\otimes_{\mathcal{A}, \varphi} \mathcal{B}
: \textit{Mod}(\mathcal{A}, \text{d})
\longrightarrow
\textit{Mod}(\mathcal{B}, \text{d})
\quad\text{et}\quad
\otimes_{\mathcal{A}, \varphi} \mathcal{B}
: \textit{Mod}^{dg}(\mathcal{A}, \text{d})
\longrightarrow
\textit{Mod}^{dg}(\mathcal{B}, \text{d})
$$
D'autre part, nous avons les foncteurs de
restriction $$
res_\varphi
: \textit{Mod}(\mathcal{B}, \text{d})
\longrightarrow
\textit{Mod}(\mathcal{A}, \text{d})
\quad\text{et}\quad
res_\varphi
: \textit{Mod}^{dg}(\mathcal{B}, \text{d})
\longrightarrow
\textit{Mod}^{dg}(\mathcal{A}, \text{d})
$$
Le lemme ci-dessus montre que ces foncteurs sont adjoints
(comme d'habitude pour la restriction et le changement de base).
À savoir, écrivons ${}_\mathcal{A}\mathcal{B}_\mathcal{B}$
pour $\mathcal{B}$ considéré comme un
$(\mathcal{A}, \mathcal{B})$-bimodule différentiel
gradué. Alors, pour tout $\mathcal{B}$-module différentiel gradué $\mathcal{L}$,
nous
avons $$
\SheafHom_\mathcal{B}^{dg}({}_\mathcal{A}\mathcal{B}_\mathcal{B}, \mathcal{L})
= res_\varphi(\mathcal{L})
$$
comme $\mathcal{A}$-modules différentiels
gradués. Ainsi, le lemme \ref{lemma-tensor-hom-adjunction-gr}
fournit un isomorphisme
fonctoriel $$
\Hom_{\textit{Mod}^{dg}(\mathcal{B}, \text{d})}(
\mathcal{M} \otimes_{\mathcal{A}, \varphi} \mathcal{B}, \mathcal{L})
= \Hom_{\textit{Mod}^{dg}(\mathcal{A}, \text{d})}(
\mathcal{M}, res_\varphi(\mathcal{L}))
$$
Lorsque le contexte est clair, nous omettons généralement $\varphi$
dans cette formule. De la même manière, nous obtenons
l'égalité
$$
\SheafHom^{dg}_\mathcal{B}(
\mathcal{M} \otimes_\mathcal{A} \mathcal{B}, \mathcal{L}) =
\SheafHom_\mathcal{A}^{dg}(\mathcal{M}, \mathcal{L})
$$
de $\mathcal{O}$-modules gradués.












\section{Image inverse et image directe des faisceaux de modules différentiels gradués}
\label{section-functoriality-dg}

\noindent
Soit $(f, f^\sharp) : (\Sh(\mathcal{C}), \mathcal{O}_\mathcal{C})
\to (\Sh(\mathcal{D}), \mathcal{O}_\mathcal{D})$ un morphisme
de topos annelés. Soient $\mathcal{A}$ une
$\mathcal{O}_\mathcal{C}$-algèbre différentielle graduée et $\mathcal{B}$ une
$\mathcal{O}_\mathcal{D}$-algèbre différentielle graduée. Supposons qu'on nous donne
une application
$$
\varphi : f^{-1}\mathcal{B} \to \mathcal{A}
$$ de
$f^{-1}\mathcal{O}_\mathcal{D}$-algèbres différentielles graduées. Par l'adjonction
entre restriction et extension des scalaires, cela revient à se donner
un morphisme $\varphi : f^*\mathcal{B} \to \mathcal{A}$ de
$\mathcal{O}_\mathcal{C}$-algèbres différentielles graduées ou, de manière équivalente,
$\varphi$ peut être considéré comme un morphisme
$$
\varphi : \mathcal{B} \to f_*\mathcal{A}
$$
de $\mathcal{O}_\mathcal{D}$-algèbres différentielles graduées.
Voir la remarque \ref{remark-functoriality-dga}.

\medskip\noindent
Définissons maintenant un
foncteur $$
f_*
: \textit{Mod}(\mathcal{A}, \text{d})
\longrightarrow
\textit{Mod}(\mathcal{B}, \text{d})
$$
Étant donné un $\mathcal{A}$-module différentiel gradué $\mathcal{M}$, nous
définissons $f_*\mathcal{M}$ comme le $\mathcal{B}$-module gradué construit dans
la section \ref{section-functoriality-graded}, muni du différentiel
donné par les morphismes $f_*d : f_*\mathcal{M}^n \to f_*\mathcal{M}^{n + 1}$. La
construction est clairement fonctorielle en
$\mathcal{M}$ et nous obtenons notre foncteur.

\medskip\noindent
Définissons maintenant un
foncteur $$
f^*
: \textit{Mod}(\mathcal{B}, \text{d})
\longrightarrow
\textit{Mod}(\mathcal{A}, \text{d})
$$
Étant donné un $\mathcal{B}$-module différentiel gradué $\mathcal{N}$,
nous définissons $f^*\mathcal{N}$ comme le $\mathcal{A}$-module
gradué construit dans la section \ref{section-functoriality-graded}.
Rappelons
que $$
f^*\mathcal{N} = f^{-1}\mathcal{N} \otimes_{f^{-1}\mathcal{B}} \mathcal{A}
$$
Puisque $f^{-1}\mathcal{N}$ est muni des applications
$f^{-1}\text{d} : f^{-1}\mathcal{N}^n \to f^{-1}\mathcal{N}^{n + 1}$, nous
pouvons considérer ce produit tensoriel comme un exemple
du produit tensoriel étudié dans la section \ref{section-dg-bimodules},
qui le munit d'un différentiel. La
construction est clairement fonctorielle en
$\mathcal{N}$ et nous obtenons notre foncteur $f^*$.

\medskip\noindent
Les foncteurs $f_*$ et $f^*$ se relèvent naturellement en des
foncteurs de catégories différentielles graduées
$$
f_*
: \textit{Mod}^{dg}(\mathcal{A}, \text{d})
\longrightarrow
\textit{Mod}^{dg}(\mathcal{B}, \text{d})
\quad\text{et}\quad f^*
:
\textit{Mod}^{dg}(\mathcal{B}, \text{d})
\longrightarrow
\textit{Mod}^{dg}(\mathcal{A}, \text{d})
$$ qui
agissent de la même façon sur les objets sous-jacents et sont définis
sur les morphismes homogènes de degré
$n$ par la fonctorialité des constructions.

\begin{lemma}
\label{lemma-adjunction-push-pull-dg} Dans
la situation ci-dessus, nous
avons $$
\Hom_{\textit{Mod}^{dg}(\mathcal{B}, \text{d})}(
\mathcal{N}, f_*\mathcal{M})
= \Hom_{\textit{Mod}^{dg}(\mathcal{A}, \text{d})}(
f^*\mathcal{N}, \mathcal{M})
$$
\end{lemma}

\begin{proof}
Démonstration omise. Indication : l'énoncé est vrai pour les catégories graduées
sous-jacentes par le lemme \ref{lemma-adjunction-push-pull-gr}. Un calcul
montre que ces isomorphismes sont compatibles avec les différentiels.
\end{proof}















\section{Localisation et faisceaux de modules différentiels gradués}
\label{section-localize-dg}

\noindent
Soit $(\mathcal{C}, \mathcal{O})$ un site annelé. Soit
$U \in \Ob(\mathcal{C})$ et notons
$$ j
:
(\Sh(\mathcal{C}/U), \mathcal{O}_U)
\longrightarrow
(\Sh(\mathcal{C}), \mathcal{O})
$$ le
morphisme de localisation correspondant (Modules
sur les sites, section \ref{sites-modules-section-localize}). Nous
utiliserons le fait suivant : pour des $\mathcal{O}_U$-modules
$\mathcal{M}_i$, $i = 1, 2$, et un $\mathcal{O}$-module $\mathcal{A}$, il
existe une application canonique
$$
j_!
: \Hom_{\mathcal{O}_U}(
\mathcal{M}_1 \otimes_{\mathcal{O}_U} \mathcal{A}|_U, \mathcal{M}_2)
\longrightarrow
\Hom_\mathcal{O}(
j_!\mathcal{M}_1 \otimes_\mathcal{O} \mathcal{A}, j_!\mathcal{M}_2)
$$
Nous avons notamment
$j_!(\mathcal{M}_1 \otimes_{\mathcal{O}_U} \mathcal{A}|_U) =
j_!\mathcal{M}_1 \otimes_\mathcal{O} \mathcal{A}$ par
Modules sur les sites, lemme \ref{sites-modules-lemma-j-shriek-and-tensor}.

\medskip\noindent
Soit $\mathcal{A}$ une $\mathcal{O}$-algèbre différentielle graduée.
Notons $\mathcal{A}_U$ la restriction de $\mathcal{A}$ à
$\mathcal{C}/U$, en d'autres termes, nous avons
$\mathcal{A}_U = j^*\mathcal{A} = j^{-1}\mathcal{A}$. Dans la section
\ref{section-functoriality-dg}, nous avons
construit des foncteurs adjoints
$$ j_*
:
\textit{Mod}^{dg}(\mathcal{A}_U, \text{d})
\longrightarrow
\textit{Mod}^{dg}(\mathcal{A}, \text{d})
\quad\text{et}\quad j^*
:
\textit{Mod}^{dg}(\mathcal{A}, \text{d})
\longrightarrow
\textit{Mod}^{dg}(\mathcal{A}_U, \text{d})
$$ avec
$j^*$ adjoint à gauche à $j_*$. Nous affirmons qu'il existe de plus un foncteur
exact
$$ j_!
:
\textit{Mod}^{dg}(\mathcal{A}_U, \text{d})
\longrightarrow
\textit{Mod}^{dg}(\mathcal{A}, \text{d})
$$ adjoint à droite
de $j_*$. À savoir, étant donné un
$\mathcal{A}_U$-module différentiel gradué $\mathcal{M}$, nous définissons $j_!\mathcal{M}$ comme
le $\mathcal{A}$-module gradué construit dans la section
\ref{section-localize-graded} avec les
différentiels
$j_!\text{d} : j_!\mathcal{M}^n \to j_!\mathcal{M}^{n + 1}$. Étant donnée
une application homogène
$f : \mathcal{M} \to \mathcal{M}'$ de degré $n$ de
$\mathcal{A}_U$-modules différentiels gradués, nous obtenons une application
homogène $j_!f : j_!\mathcal{M} \to j_!\mathcal{M}'$ de degré $n$
de $\mathcal{A}$-modules différentiels gradués. Nous omettons
la simple vérification que cette
construction est compatible avec les différentiels. Ainsi
nous obtenons notre foncteur.

\begin{lemma}
\label{lemma-extension-by-zero-dg}
Dans la situation ci-dessus, nous
avons $$
\Hom_{\textit{Mod}^{dg}(\mathcal{A}, \text{d})}(
j_!\mathcal{M}, \mathcal{N})
= \Hom_{\textit{Mod}^{dg}(\mathcal{A}_U, \text{d})}(
\mathcal{M}, j^*\mathcal{N})
$$
\end{lemma}

\begin{proof}
Démonstration omise. Nous avons vu dans le lemme \ref{lemma-extension-by-zero-graded}
que le lemme est vrai au niveau gradué. Ainsi, tout ce qui doit être
vérifié est que l'isomorphisme résultant est compatible avec les différentiels.
\end{proof}

\begin{lemma}
\label{lemma-tensor-with-extension-by-zero-dg} Dans
la situation ci-dessus, soit $\mathcal{M}$ un
$\mathcal{A}_U$-module différentiel gradué à droite et soit $\mathcal{N}$ un
$\mathcal{A}$-module différentiel gradué à gauche. Alors
$$
j_!\mathcal{M} \otimes_\mathcal{A} \mathcal{N} =
j_!(\mathcal{M} \otimes_{\mathcal{A}_U} \mathcal{N}|_U)
$$ comme
complexes de $\mathcal{O}$-modules,
fonctoriellement en $\mathcal{M}$ et $\mathcal{N}$.
\end{lemma}

\begin{proof} Au
niveau des modules gradués, cela découle
du lemme \ref{lemma-tensor-with-extension-by-zero}.
Nous omettons la vérification que cet isomorphisme
est compatible avec les différentiels.
\end{proof}








\section{Foncteurs de décalage sur les faisceaux de modules différentiels gradués}
\label{section-shift-dg}

\noindent
Soit $(\mathcal{C}, \mathcal{O})$ un site annelé. Soit
$\mathcal{A}$ un faisceau d'algèbres différentielles graduées
sur $(\mathcal{C}, \mathcal{O})$.
Soit $\mathcal{M}$ un $\mathcal{A}$-module différentiel gradué.
Soit $k \in \mathbf{Z}$. Nous définissons le {\it $k$-ième décalage de} $\mathcal{M}$,
noté $\mathcal{M}[k]$, comme
suit : \begin{enumerate}
\item le $\mathcal{A}$-module gradué $\mathcal{M}[k]$ est
celui défini dans la section \ref{section-shift},
\item le différentiel
$d_{\mathcal{M}[k]} : (\mathcal{M}[k])^n \to (\mathcal{M}[k])^{n + 1}$
est donné par
$(-1)^k\text{d}_\mathcal{M} : \mathcal{M}^{n + k} \to \mathcal{M}^{n + k + 1}$.
\end{enumerate} Pour
un homomorphisme $f : \mathcal{L} \to \mathcal{M}$ de $\mathcal{A}$-modules homogène
de degré $n$, définissons $f[k] : \mathcal{L}[k] \to \mathcal{M}[k]$ par les
mêmes applications composantes que $f$. Alors $f[k]$ est
un homomorphisme de $\mathcal{A}$-modules homogène de degré $n$.
Cela définit
une application $$
\Hom_{\textit{Mod}^{dg}(\mathcal{A}, \text{d})}(\mathcal{L}, \mathcal{M})
\longrightarrow
\Hom_{\textit{Mod}^{dg}(\mathcal{A}, \text{d})}
(\mathcal{L}[k], \mathcal{M}[k])
$$
compatible avec les différentiels, car les
différentiels de $\mathcal{L}$ et
$\mathcal{M}$ sont multipliés par le même signe. Ces choix sont compatibles avec le choix
en Algèbre différentielle graduée,
définition \ref{dga-definition-shift}. Il est clair que nous avons
défini un foncteur
$$ [k]
:
\textit{Mod}^{dg}(\mathcal{A}, \text{d})
\longrightarrow
\textit{Mod}^{dg}(\mathcal{A}, \text{d})
$$ Ce foncteur est un foncteur
de catégories différentielles graduées et l'on a
$[k + l] = [k] \circ [l]$.

\medskip\noindent
Nous affirmons que l'isomorphisme
$$
\Hom_{\textit{Mod}^{dg}(\mathcal{A}, \text{d})}(\mathcal{L}, \mathcal{M}[k])
=
\Hom_{\textit{Mod}^{dg}(\mathcal{A}, \text{d})}(\mathcal{L}, \mathcal{M})[k]
$$ défini à la
section \ref{section-shift} sur les modules gradués sous-jacents est
compatible avec les
différentiels. Pour le voir, soit une application $\mathcal{A}$-linéaire
$f : \mathcal{L} \to \mathcal{M}[k]$ homogène de degré $n$ ; c'est un élément de degré
$n$ du membre de gauche. Notons
$f' : \mathcal{L} \to \mathcal{M}$ l'homomorphisme homogène de $\mathcal{A}$-modules de degré
$n + k$ dont les applications composantes sont les {\bf mêmes} que celles de $f$.
Selon nos conventions, c'est l'élément correspondant de degré $n$ du membre
de droite.
Par définition du différentiel du membre de gauche, nous
obtenons $$
\text{d}_{LHS}(f)
= \text{d}_{\mathcal{M}[k]} \circ f - (-1)^n f \circ \text{d}_\mathcal{L}
= (-1)^k\text{d}_\mathcal{M} \circ f - (-1)^n f \circ \text{d}_\mathcal{L}
$$
et, pour le différentiel du membre de droite, nous
obtenons $$
\text{d}_{RHS}(f')
= (-1)^k\left(
\text{d}_\mathcal{M} \circ f' - (-1)^{n + k} f' \circ \text{d}_\mathcal{L}
\right)
= (-1)^k\text{d}_\mathcal{M} \circ f' - (-1)^n f' \circ \text{d}_\mathcal{L}
$$ Les
applications composantes sont les mêmes, ce qui achève la démonstration.










\section{La catégorie homotopique}
\label{section-homotopy}

\noindent
Cette section est l'analogue de
l'Algèbre différentielle graduée, section \ref{dga-section-homotopy}.

\begin{definition}
\label{definition-homotopy} Soit
$(\mathcal{C}, \mathcal{O})$ un site annelé. Soit
$\mathcal{A}$ un faisceau d'algèbres différentielles graduées sur
$(\mathcal{C}, \mathcal{O})$. Soient
$f, g : \mathcal{M} \to \mathcal{N}$ des
homomorphismes de modules différentiels gradués sur $\mathcal{A}$. Une
{\it homotopie entre $f$ et $g$} est un morphisme gradué de $\mathcal{A}$-modules
$h : \mathcal{M} \to \mathcal{N}$, homogène de degré $-1$, tel
que
$$ f - g
= \text{d}_\mathcal{N} \circ h + h \circ \text{d}_\mathcal{M}
$$ S'il existe
une telle homotopie, nous disons que $f$ et $g$ sont {\it homotopes}.
\end{definition}

\noindent
Dans la situation de la définition, si nous avons des
applications $a : \mathcal{K} \to \mathcal{M}$ et $c : \mathcal{N} \to \mathcal{L}$
alors nous voyons
que $$
\begin{matrix}
h\text{ est une homotopie} \\
\text{ entre }f\text{ et } g
\end{matrix}
\quad
\Rightarrow
\quad
\begin{matrix}
c \circ h \circ a\text{ est une homotopie}\\
\text{entre } c
\circ f \circ a\text{ et } c\circ g \circ a
\end{matrix}
$$ Ainsi,
nous pouvons définir la composition des classes d'homotopie des morphismes
dans $\textit{Mod}(\mathcal{A}, \text{d})$.

\begin{definition}
\label{definition-complexes-notation}
Soit $(\mathcal{C}, \mathcal{O})$ un site annelé. Soit
$\mathcal{A}$ un faisceau d'algèbres différentielles graduées
sur $(\mathcal{C}, \mathcal{O})$.
La {\it catégorie homotopique}, notée $K(\textit{Mod}(\mathcal{A}, \text{d}))$, est
la catégorie dont les objets sont ceux de
$\textit{Mod}(\mathcal{A}, \text{d})$ et dont les morphismes sont les classes d'homotopie
d'homomorphismes de $\mathcal{A}$-modules différentiels gradués.
\end{definition}

\noindent
La notation $K(\textit{Mod}(\mathcal{A}, \text{d}))$
n'est pas standard, mais elle
est compatible avec les autres usages de $K(-)$ dans le projet Champs.

\medskip\noindent
Dans Algèbre différentielle graduée, définition
\ref{dga-definition-homotopy-category-of-dga-category},
nous avons défini la catégorie des complexes
$\text{Comp}(\mathcal{S})$ et la catégorie
homotopique $K(\mathcal{S})$ d'une catégorie
différentielle graduée $\mathcal{S}$. En appliquant
cela à la catégorie différentielle graduée
$\textit{Mod}^{dg}(\mathcal{A}, \text{d})$, nous obtenons
$$
\textit{Mod}(\mathcal{A}, \text{d}) =
\text{Comp}(\textit{Mod}^{dg}(\mathcal{A}, \text{d}))
$$ (voir la
discussion de la section \ref{section-dgm-dg-cat}) et nous obtenons
$$
K(\textit{Mod}(\mathcal{A}, \text{d})) =
K(\textit{Mod}^{dg}(\mathcal{A}, \text{d}))
$$
Pour vérifier cette dernière égalité, remarquons que nous avons
$$
\text{d}_{\Hom_{\textit{Mod}^{dg}(\mathcal{A}, \text{d})}
(\mathcal{M}, \mathcal{N})}(h) =
\text{d}_\mathcal{N} \circ h + h \circ \text{d}_\mathcal{M}
$$ lorsque
$h$ est comme dans la définition \ref{definition-homotopy}. Nous omettons les
détails.

\begin{lemma}
\label{lemma-homotopy-direct-sums} Soit
$(\mathcal{C}, \mathcal{O})$ un site annelé. Soit
$\mathcal{A}$ un faisceau d'algèbres différentielles graduées sur
$(\mathcal{C}, \mathcal{O})$. La catégorie
homotopique $K(\textit{Mod}(\mathcal{A}, \text{d}))$ possède des
sommes directes et des produits.
\end{lemma}

\begin{proof}
Démonstration omise. Indication : utiliser les sommes directes et
les produits du lemme \ref{lemma-dgm-abelian}. Cela fonctionne car nous
avons vu que ces foncteurs commutent au foncteur d'oubli vers la catégorie
des modules gradués sur $\mathcal{A}$, et car $\prod$ et $\bigoplus$ sont
des foncteurs exacts sur la catégorie des familles de groupes abéliens.
\end{proof}










\section{Cônes et triangles}
\label{section-conclude-triangulated}

\noindent
Dans cette section, nous utilisons les
résultats d'Algèbre différentielle graduée, section \ref{dga-section-review},
pour conclure que la catégorie homotopique
de la catégorie des $\mathcal{A}$-modules différentiels
gradués est triangulée.

\begin{lemma}
\label{lemma-axioms-AB}
Soit $(\mathcal{C}, \mathcal{O})$ un site annelé. Soit
$\mathcal{A}$ un faisceau d'algèbres différentielles graduées
sur $(\mathcal{C}, \mathcal{O})$. La
catégorie différentielle graduée
$\textit{Mod}^{dg}(\mathcal{A}, \text{d})$
satisfait aux axiomes (A) et (B)
d'Algèbre différentielle graduée, section \ref{dga-section-review}.
\end{lemma}

\begin{proof}
Supposons donnés des $\mathcal{A}$-modules différentiels
gradués $\mathcal{M}$ et $\mathcal{N}$. Considérons
le $\mathcal{A}$-module différentiel gradué $\mathcal{M} \oplus \mathcal{N}$
défini de manière évidente. Les coprojections
$i : \mathcal{M} \to \mathcal{M} \oplus \mathcal{N}$ et
$j : \mathcal{N} \to \mathcal{M} \oplus \mathcal{N}$ ainsi que les
projections
$p : \mathcal{M} \oplus \mathcal{N} \to \mathcal{N}$ et
$q : \mathcal{M} \oplus \mathcal{N} \to \mathcal{M}$ sont des morphismes
de modules différentiels gradués sur $\mathcal{A}$. Ainsi,
$i, j, p, q$ sont homogènes de degré
$0$ et fermés, c'est-à-dire que $\text{d}(i) = 0$, etc. Ainsi,
cette somme directe est une somme différentielle graduée au sens de
l'Algèbre différentielle graduée, définition \ref{dga-definition-dg-direct-sum}.
Cela prouve l'axiome (A).

\medskip\noindent
L'axiome (B) a été vérifié dans la section \ref{section-shift-dg}.
\end{proof}

\noindent
Soit $(\mathcal{C}, \mathcal{O})$ un site annelé. Soit
$\mathcal{A}$ un faisceau d'algèbres différentielles graduées sur
$(\mathcal{C}, \mathcal{O})$.
Rappelons qu'une suite
$$ 0
\to \mathcal{K} \to \mathcal{L} \to \mathcal{N} \to 0
$$ dans
$\textit{Mod}(\mathcal{A}, \text{d})$ est appelée
une {\it suite exacte courte admissible} (voir
Algèbre différentielle graduée, section \ref{dga-section-review}) si elle
est scindée dans $\textit{Mod}(\mathcal{A})$. Autrement dit, elle est scindée comme
suite de $\mathcal{A}$-modules gradués. Notons
$s : \mathcal{N} \to \mathcal{L}$ et
$\pi : \mathcal{L} \to \mathcal{K}$ les scindages
gradués $\mathcal{A}$-linéaires. En combinant
le lemme \ref{lemma-axioms-AB} et Algèbre
différentielle graduée, lemme \ref{dga-lemma-get-triangle}, nous obtenons
un triangle
$$
\mathcal{K} \to \mathcal{L} \to \mathcal{N} \to \mathcal{K}[1]
$$ où la flèche
$\mathcal{N} \to \mathcal{K}[1]$ dans la preuve d'Algèbre
différentielle graduée, lemme \ref{dga-lemma-get-triangle} est construite
comme
$$
\delta = \pi \circ
\text{d}_{\Hom_{\textit{Mod}^{dg}(\mathcal{A}, \text{d})}
(\mathcal{L}, \mathcal{M})}(s) =
\pi \circ \text{d}_\mathcal{L} \circ s -
\pi \circ s \circ \text{d}_\mathcal{N} =
\pi \circ \text{d}_\mathcal{L} \circ s
$$ avec nos excuses pour cette
notation peu lisible. Quoi qu'il en soit, nous voyons que, dans notre
situation, l'application de bord $\delta$ construite dans Algèbre
différentielle graduée, lemme \ref{dga-lemma-get-triangle}, coïncide
sur les complexes sous-jacents de $\mathcal{O}$-modules avec
l'application de bord usuelle employée dans le projet Champs pour
les suites exactes courtes de complexes scindées terme à terme ;
voir Catégories dérivées, définition \ref{derived-definition-split-ses}.

\begin{definition}
\label{definition-cone}
Soit $(\mathcal{C}, \mathcal{O})$ un site annelé. Soit
$\mathcal{A}$ un faisceau d'algèbres différentielles graduées sur
$(\mathcal{C}, \mathcal{O})$. Soit
$f : \mathcal{K} \to \mathcal{L}$ un
homomorphisme de modules différentiels gradués sur $\mathcal{A}$. Le
{\it cône} de $f$ est le $\mathcal{A}$-module différentiel gradué
$C(f)$ défini comme suit :
\begin{enumerate}
\item le complexe sous-jacent de $\mathcal{O}$-modules est le cône
du morphisme correspondant
$f : \mathcal{K}^\bullet \to \mathcal{L}^\bullet$ de complexes de
$\mathcal{A}$-modules ; ainsi,
$C(f)^n = \mathcal{L}^n \oplus \mathcal{K}^{n + 1}$ et le différentiel
est
$$
d_{C(f)} =
\left(
\begin{matrix}
\text{d}_\mathcal{L} & f \\ 0 &
-\text{d}_\mathcal{K}
\end{matrix}
\right)
$$
\item l'application de multiplication
$$ C(f)^n
\times \mathcal{A}^m \to C(f)^{n + m}
$$ est la somme
directe de l'application de multiplication
$\mathcal{L}^n \times \mathcal{A}^m \to \mathcal{L}^{n + m}$ et
de l'application de multiplication
$\mathcal{K}^{n + 1} \times \mathcal{A}^m \to \mathcal{K}^{n + 1 + m}$.
\end{enumerate}
Il est équipé d'homomorphismes canoniques de
$\mathcal{A}$-modules différentiels gradués $i : \mathcal{L} \to C(f)$ et
$p : C(f) \to \mathcal{K}[1]$ induits par les applications évidentes.
\end{definition}

\noindent
Remarquons que, dans la situation de la définition, la suite
$$
0 \to \mathcal{L} \to C(f) \to \mathcal{K}[1] \to 0
$$
est une suite exacte courte admissible.

\begin{lemma}
\label{lemma-axiom-C}
Soit $(\mathcal{C}, \mathcal{O})$ un site annelé. Soit
$\mathcal{A}$ un faisceau d'algèbres différentielles graduées sur
$(\mathcal{C}, \mathcal{O})$. La
catégorie différentielle graduée
$\textit{Mod}^{dg}(\mathcal{A}, \text{d})$
satisfait à l'axiome (C) formulé dans
Algèbre différentielle graduée, situation \ref{dga-situation-ABC}.
\end{lemma}

\begin{proof}
Soit $f : \mathcal{K} \to \mathcal{L}$ un
homomorphisme de modules différentiels gradués sur $\mathcal{A}$. D'après
ce qui précède, nous avons une suite exacte courte admissible
$$
0 \to \mathcal{L} \to C(f) \to \mathcal{K}[1] \to 0
$$ Pour
achever la démonstration, nous devons vérifier que l'application de
bord $$
\delta : \mathcal{K}[1] \to \mathcal{L}[1]
$$
qui lui est associée (voir la discussion ci-dessus) est égale à $f[1]$.
Pour la section $s : \mathcal{K}[1] \to C(f)$, nous utilisons en
degré $n$ le plongement $\mathcal{K}^{n + 1} \to C(f)^n$. En degré
$n$, l'application $\pi$ est donnée par les projections
$C(f)^n \to \mathcal{L}^n$. Alors finalement nous devons calculer
$$
\delta = \pi \circ \text{d}_{C(f)} \circ s
$$ (voir
discussion ci-dessus). En notation matricielle, cette expression est
égale à $$
\left(
\begin{matrix}
1
& 0 \end{matrix}
\right)
\left(
\begin{matrix}
\text{d}_\mathcal{L} & f \\
0 & -\text{d}_\mathcal{K}
\end{matrix}
\right)
\left(
\begin{matrix} 0
\\
1
\end{matrix}
\right) = f
$$
comme souhaité.
\end{proof}

\noindent À
ce stade, nous savons que tous les lemmes prouvés en
Algèbre différentielle graduée, section \ref{dga-section-review} sont
valables pour la catégorie différentielle graduée
$\textit{Mod}^{dg}(\mathcal{A}, \text{d})$. En particulier,
nous obtenons la proposition suivante.

\begin{proposition}
\label{proposition-homotopy-category-triangulated}
Soit $(\mathcal{C}, \mathcal{O})$ un site annelé. Soit
$\mathcal{A}$ un faisceau d'algèbres différentielles graduées
sur $(\mathcal{C}, \mathcal{O})$.
La catégorie homotopique $K(\textit{Mod}(\mathcal{A}, \text{d}))$
est une catégorie triangulée
où \begin{enumerate}
\item les foncteurs de décalage sont ceux de la
section \ref{section-shift-dg},
\item les triangles distingués sont les triangles de
$K(\textit{Mod}(\mathcal{A}, \text{d}))$ isomorphes, en
tant que triangles, à un triangle
$$
\mathcal{K} \to \mathcal{L} \to \mathcal{N}
\xrightarrow{\delta} \mathcal{K}[1],\quad\quad
\delta = \pi \circ \text{d}_\mathcal{L} \circ s
$$ construit
à partir d'une suite exacte courte admissible
$0 \to \mathcal{K} \to \mathcal{L} \to \mathcal{N} \to 0$ dans
$\textit{Mod}(\mathcal{A}, \text{d})$ ci-dessus.
\end{enumerate}
\end{proposition}

\begin{proof}
Rappelons que $K(\textit{Mod}(\mathcal{A}, \text{d}))
= K(\textit{Mod}^{dg}(\mathcal{A}, \text{d}))$,
voir la section \ref{section-homotopy}.
La proposition découle
alors des lemmes \ref{lemma-axioms-AB} et \ref{lemma-axiom-C} et
d'Algèbre
différentielle graduée, proposition
\ref{dga-proposition-ABC-homotopy-category-triangulated}.
\end{proof}

\begin{remark}
\label{remark-cone-identity}
Soit $(\mathcal{C}, \mathcal{O})$ un site annelé. Soit
$\mathcal{A}$ un faisceau d'algèbres différentielles graduées
sur $(\mathcal{C}, \mathcal{O})$. Soit $C = C(\text{id}_\mathcal{A})$
le cône de l'application identité $\mathcal{A} \to \mathcal{A}$, considérée
comme un morphisme de $\mathcal{A}$-modules différentiels
gradués.
Alors $$
\Hom_{\textit{Mod}(\mathcal{A}, \text{d})}(C, \mathcal{M})
= \{(x, y) \in
\Gamma(\mathcal{C}, \mathcal{M}^0) \times
\Gamma(\mathcal{C}, \mathcal{M}^{-1}) \mid
x = \text{d}(y)\}
$$
où l'application de gauche à droite envoie $f$ à la paire $(x, y)$
où $x$ est l'image de la section globale $(0, 1)$
de $C^{-1} = \mathcal{A}^{-1} \oplus \mathcal{A}^0$ et
où $y$ est l'image de la section globale $(1, 0)$ de
$C^0 = \mathcal{A}^0 \oplus \mathcal{A}^1$.
\end{remark}

\begin{lemma}
\label{lemma-dgm-grothendieck-abelian}
Soit $(\mathcal{C}, \mathcal{O})$ un site annelé. Soit
$(\mathcal{A}, \text{d})$ une $\mathcal{O}$-algèbre différentielle graduée. La
catégorie $\textit{Mod}(\mathcal{A}, \text{d})$ est une
catégorie abélienne de Grothendieck.
\end{lemma}

\begin{proof}
Par le lemme \ref{lemma-dgm-abelian} et la définition d'une catégorie abélienne
de Grothendieck
(Injectifs, définition \ref{injectives-definition-grothendieck-conditions}), il
suffit de montrer
que $\textit{Mod}(\mathcal{A}, \text{d})$ possède
un générateur. Pour chaque objet $U$ de $\mathcal{C}$, notons
$C_U$ le cône du morphisme identité $\mathcal{A}_U \to \mathcal{A}_U$, comme
dans la remarque \ref{remark-cone-identity}. Nous affirmons
que $$
\mathcal{G} = \bigoplus\nolimits_{k, U} j_{U!}C_U[k]
$$
est un générateur, où la somme porte sur tous les objets $U$ de $\mathcal{C}$
et tous les $k \in \mathbf{Z}$. En
effet, étant donné un $\mathcal{A}$-module différentiel gradué $\mathcal{M}$,
s'il n'existe aucun morphisme non nul de $\mathcal{G}$ vers $\mathcal{M}$,
alors, pour tous $k$ et $U$,
le groupe \begin{align*}
\Hom_{\textit{Mod}(\mathcal{A})}(j_{U!}C_U[k], \mathcal{M}) \\
&
= \Hom_{\textit{Mod}(\mathcal{A}_U)}(C_U[k], \mathcal{M}|_U) \\
&
= \{(x, y) \in \mathcal{M}^{-k}(U) \times \mathcal{M}^{-k - 1}(U) \mid
x = \text{d}(y)\}
\end{align*}
est égal à zéro. Ainsi, $\mathcal{M}$ est nul.
\end{proof}











\section{Résolutions plates}
\label{section-P-resolutions}

\noindent
Cette section est l'analogue de
l'Algèbre différentielle graduée, section \ref{dga-section-P-resolutions}.

\medskip\noindent
Soit $(\mathcal{C}, \mathcal{O})$ un site annelé. Soit
$\mathcal{A}$ un faisceau d'algèbres différentielles graduées sur
$(\mathcal{C}, \mathcal{O})$. Nous dirons qu'un
$\mathcal{A}$-module différentiel gradué à droite $\mathcal{P}$ est {\it bon}
si
\begin{enumerate}
\item le foncteur
$\mathcal{N} \mapsto \mathcal{P} \otimes_\mathcal{A} \mathcal{N}$ est
exact sur la catégorie des $\mathcal{A}$-modules gradués à gauche,
\item si $\mathcal{N}$ est un
$\mathcal{A}$-module différentiel gradué acyclique à gauche, alors
$\mathcal{P} \otimes_\mathcal{A} \mathcal{N}$ est acyclique,
\item pour tout morphisme $(f, f^\sharp) : (\Sh(\mathcal{C}'), \mathcal{O}')
\to (\Sh(\mathcal{C}), \mathcal{O})$ de topos
annelés, toute $\mathcal{O}'$-algèbre différentielle graduée
$\mathcal{A}'$ et tout morphisme $\varphi : f^{-1}\mathcal{A} \to \mathcal{A}'$ d'algèbres
différentielles graduées sur $f^{-1}\mathcal{O}_\mathcal{D}$, les propriétés
(1) et (2) valent pour l'image inverse $f^*\mathcal{P}$
(section \ref{section-functoriality-dg}),
considérée comme un $\mathcal{A}'$-module différentiel gradué.
\end{enumerate}
La première condition signifie que $\mathcal{P}$ est plat comme module gradué
à droite sur $\mathcal{A}$ ; la deuxième signifie que $\mathcal{P}$ est
K-plat au sens de Spaltenstein (voir Cohomologie
sur les sites, section \ref{sites-cohomology-section-flat}) ; la troisième
exige que ces propriétés restent vraies après tout changement de base.

\medskip\noindent Il
est peut-être surprenant qu'il y ait beaucoup de bons modules.

\begin{lemma}
\label{lemma-supply-good}
Soit $(\mathcal{C}, \mathcal{O})$ un site annelé. Soit
$\mathcal{A}$ un faisceau d'algèbres différentielles graduées
sur $(\mathcal{C}, \mathcal{O})$. Soit $U \in \Ob(\mathcal{C})$.
Alors $j_!\mathcal{A}_U$ est un bon
$\mathcal{A}$-module différentiel gradué.
\end{lemma}

\begin{proof}
Soit $\mathcal{N}$ un module gradué à gauche sur $\mathcal{A}$.
Par le lemme \ref{lemma-tensor-with-extension-by-zero} nous avons
$$
j_!\mathcal{A}_U \otimes_\mathcal{A} \mathcal{N} =
j_!(\mathcal{A}_U \otimes_{\mathcal{A}_U} \mathcal{N}|_U) =
j_!(\mathcal{N}_U)
$$ comme modules
gradués. Puisque la
restriction à $U$ et le foncteur $j_!$ sont tous deux exacts, cela
prouve (1). Le même argument est valable pour (2) en utilisant
le lemme \ref{lemma-tensor-with-extension-by-zero-dg}.

\medskip\noindent
Considérons un morphisme $(f, f^\sharp) : (\Sh(\mathcal{C}'), \mathcal{O}')
\to (\Sh(\mathcal{C}), \mathcal{O})$ de
topos annelés, une $\mathcal{O}'$-algèbre différentielle graduée
$\mathcal{A}'$ et un morphisme $\varphi : f^{-1}\mathcal{A} \to \mathcal{A}'$
de $f^{-1}\mathcal{O}$-algèbres différentielles graduées.
Nous devons montrer
que $$
f^*j_!\mathcal{A}_U
= f^{-1}j_!\mathcal{A}_U \otimes_{f^{-1}\mathcal{A}} \mathcal{A}'
$$
satisfait aux points (1) et (2) pour le
topos annelé $(\Sh(\mathcal{C}'), \mathcal{O}')$ doté
du faisceau de $\mathcal{O}'$-algèbres différentielles
graduées $\mathcal{A}'$. Pour le prouver, nous
pouvons remplacer $(\Sh(\mathcal{C}), \mathcal{O})$
et $(\Sh(\mathcal{C}'), \mathcal{O}')$ par des topos
annelés
équivalents. Ainsi, par
Modules sur les sites, lemme \ref{sites-modules-lemma-morphism-ringed-topoi-comes-from-morphism-ringed-sites},
nous pouvons supposer que $f$ provient
d'un morphisme de sites $f : \mathcal{C} \to \mathcal{C}'$
donné par un foncteur continu $u : \mathcal{C} \to \mathcal{C}'$.
Dans ce cas, posons $U' = u(U)$
et notons $j' : \Sh(\mathcal{C}'/U') \to \Sh(\mathcal{C}')$
le morphisme de localisation
correspondant. On obtient un carré commutatif de morphismes
de topos annelés $$
\xymatrix{
(\Sh(\mathcal{C}'/U'), \mathcal{O}'_{U'})
\ar[rr]_{(j', (j')^\sharp)} \ar[d]_{(f', (f')^\sharp)}
& & (\Sh(\mathcal{C}'), \mathcal{O}')
\ar[d]^{(f, f^\sharp)} \\
(\Sh(\mathcal{C}/U), \mathcal{O}_U)
\ar[rr]^{(j, j^\sharp)}
& & (\Sh(\mathcal{C}), \mathcal{O}).
}
$$
et nous avons $f'_*(j')^{-1} = j^{-1}f_*$.
Voir Modules sur
les sites, lemme \ref{sites-modules-lemma-localize-morphism-ringed-sites}.
Par l'unicité des adjoints,
on obtient $f^{-1}j_! = j'_!(f')^{-1}$.
Ainsi nous obtenons \begin{align*}
f^*j_!\mathcal{A}_U
&
= f^{-1}j_!\mathcal{A}_U \otimes_{f^{-1}\mathcal{A}} \mathcal{A}' \\
&
= j'_!(f')^{-1}\mathcal{A}_U \otimes_{f^{-1}\mathcal{A}} \mathcal{A}' \\
&
= j'_!\left(
(f')^{-1}\mathcal{A}_U \otimes_{f^{-1}\mathcal{A}|_{U'}}
\mathcal{A}'|_{U'}\right) \\
&
= j'_!\mathcal{A}'_{U'}
\end{align*}
La première équation est la définition de l'image inverse de $j_!\mathcal{A}_U$
comme module différentiel gradué sur $\mathcal{A}'$.
La deuxième égalité découle de $f^{-1}j_! = j'_!(f')^{-1}$.
La troisième découle du lemme \ref{lemma-tensor-with-extension-by-zero-dg},
appliqué au site annelé $(\mathcal{C}', f^{-1}\mathcal{O})$ muni
du faisceau d'algèbres différentielles graduées $f^{-1}\mathcal{A}$,
avec les modules différentiels gradués $(f')^{-1}\mathcal{A}_U$ sur $\mathcal{C}'/U'$
et $\mathcal{A}'$ sur $\mathcal{C}'$.
La quatrième égalité résulte de
$(f')^{-1}\mathcal{A}_U = f^{-1}\mathcal{A}|_{U'}$. Nous voyons donc que
l'image inverse est encore un module du même type, et nous
avons démontré ci-dessus les conditions (1) et (2) pour celui-ci.
\end{proof}

\begin{lemma}
\label{lemma-good-admissible-ses}
Soit $(\mathcal{C}, \mathcal{O})$ un site annelé.
Soit $\mathcal{A}$ un faisceau d'algèbres différentielles
graduées sur $(\mathcal{C}, \mathcal{O})$.
Soit $0 \to \mathcal{P} \to \mathcal{P}' \to \mathcal{P}'' \to 0$
une suite exacte courte admissible de
$\mathcal{A}$-modules différentiels gradués. Si deux modules sur trois sont
bons, le troisième aussi.
\end{lemma}

\begin{proof}
La condition (1) est immédiate, car la suite est scindée au niveau gradué.
Pour la condition (2), remarquons que, pour tout module
différentiel gradué à gauche sur $\mathcal{A}$, la suite
$$ 0
\to
\mathcal{P} \otimes_\mathcal{A} \mathcal{N} \to
\mathcal{P}' \otimes_\mathcal{A} \mathcal{N} \to
\mathcal{P}'' \otimes_\mathcal{A} \mathcal{N} \to 0
$$ est une
suite exacte courte admissible de
$\mathcal{O}$-modules différentiels gradués (puisque, après oubli du différentiel, le
produit tensoriel se calcule dans la catégorie des modules gradués). Ainsi,
si deux des trois sont exacts comme complexes de $\mathcal{O}$-modules, le
troisième aussi. Enfin, le même argument montre que, pour un
morphisme $(f, f^\sharp) : (\Sh(\mathcal{C}'), \mathcal{O}')
\to (\Sh(\mathcal{C}), \mathcal{O})$ de topos
annelés, une $\mathcal{O}'$-algèbre différentielle graduée
$\mathcal{A}'$ et un morphisme $\varphi : f^{-1}\mathcal{A} \to \mathcal{A}'$ de
$f^{-1}\mathcal{O}$-algèbres différentielles graduées,
la suite
$$ 0
\to f^*\mathcal{P} \to f^*\mathcal{P}' \to f^*\mathcal{P}'' \to 0
$$ est
une suite exacte courte admissible de
$\mathcal{A}'$-modules différentiels gradués et le même argument que ci-dessus s'applique ici.
\end{proof}

\begin{lemma}
\label{lemma-good-direct-sum} Soit
$(\mathcal{C}, \mathcal{O})$ un site annelé. Soit
$\mathcal{A}$ un faisceau d'algèbres différentielles graduées sur
$(\mathcal{C}, \mathcal{O})$. Une somme directe
arbitraire de bons $\mathcal{A}$-modules différentiels gradués est
bonne. Une limite inductive filtrante de bons
$\mathcal{A}$-modules différentiels gradués est bonne.
\end{lemma}

\begin{proof}
Démonstration omise. Indication : les sommes directes et les limites inductives
filtrantes commutent aux produits tensoriels et aux images inverses.
\end{proof}

\begin{lemma}
\label{lemma-good-quotient}
Soit $(\mathcal{C}, \mathcal{O})$ un site annelé.
Soit $\mathcal{A}$ un faisceau d'algèbres différentielles graduées
sur $(\mathcal{C}, \mathcal{O})$. Soit $\mathcal{M}$
un $\mathcal{A}$-module différentiel gradué. Il existe un homomorphisme
$\mathcal{P} \to \mathcal{M}$ de $\mathcal{A}$-modules différentiels gradués
avec les propriétés suivantes
\begin{enumerate}
\item $\mathcal{P} \to \mathcal{M}$ est surjectif,
\item $\Ker(\text{d}_\mathcal{P}) \to \Ker(\text{d}_\mathcal{M})$
est surjectif, et
\item $\mathcal{P}$ est bon.
\end{enumerate}
\end{lemma}

\begin{proof}
Considérons les triples $(U, k, x)$ où $U$ est un objet de $\mathcal{C}$,
$k \in \mathbf{Z}$, et $x$ est une section de $\mathcal{M}^k$ sur $U$ avec
$\text{d}_\mathcal{M}(x) = 0$. Nous obtenons alors un unique morphisme
de $\mathcal{A}_U$-modules différentiels gradués
$\varphi_x : \mathcal{A}_U[-k] \to \mathcal{M}|_U$
envoyant $1$ sur $x$. Ceci est adjoint à un morphisme
$\psi_x : j_{U!}\mathcal{A}_U[-k] \to \mathcal{M}$.
Remarquons que $1 \in \mathcal{A}_U(U)$ correspond à une
section $1 \in j_{U!}\mathcal{A}_U[-k](U)$ de degré $k$, dont le
différentiel est nul et qui est envoyée sur $x$ par $\psi_x$. Ainsi, si
nous considérons l'application
$$
\bigoplus\nolimits_{(U, k, x)} j_{U!}\mathcal{A}_U[-k]
\longrightarrow
\mathcal{M}
$$
les conditions (2) et (3) seront satisfaites. À savoir, les objets
$j_{U!}\mathcal{A}_U[-k]$ sont bons (lemme \ref{lemma-supply-good})
et toute somme directe de bons objets est bonne (lemme \ref{lemma-good-direct-sum}).

\medskip\noindent
Considérons ensuite les triples $(U, k, x)$ où $U$ est un objet de $\mathcal{C}$,
$k \in \mathbf{Z}$ et $x$ une section de $\mathcal{M}^k$, pas nécessairement annulée par
le différentiel. Ensuite, nous pouvons considérer le cône
$C_U$ de l'application identité $\mathcal{A}_U \to \mathcal{A}_U$ comme dans
la remarque \ref{remark-cone-identity}. L'élément $x$ déterminera une application
$\varphi_x : C_U[-k - 1] \to \mathcal{A}_U$, voir remarque
\ref{remark-cone-identity}. Maintenant, puisque nous
avons une suite exacte courte admissible
$$
0 \to \mathcal{A}_U \to C_U \to \mathcal{A}_U[1] \to 0
$$
nous concluons que $j_{U!}C_U$ est un bon module par le
lemme \ref{lemma-good-admissible-ses} et le
lemme \ref{lemma-supply-good} déjà utilisé. Comme
ci-dessus, nous concluons que la somme directe des applications
$\psi_x : j_{U!}C_U \to \mathcal{M}$ adjointes aux $\varphi_x$
$$
\bigoplus\nolimits_{(U, k, x)} j_{U!}C_U \longrightarrow \mathcal{M}
$$ est
surjective. En prenant la somme directe avec l'application produite
au premier paragraphe, nous concluons.
\end{proof}

\begin{remark}
\label{remark-sheaf-graded-sets}
Soit $(\mathcal{C}, \mathcal{O})$ un site annelé. Un
{\it faisceau d'ensembles gradués} sur $\mathcal{C}$ est un faisceau
d'ensembles $\mathcal{S}$ muni d'un morphisme
$\deg : \mathcal{S} \to \underline{\mathbf{Z}}$
de faisceaux d'ensembles. Notons $\mathcal{O}[\mathcal{S}]$
le $\mathcal{O}$-module gradué qui est le
$\mathcal{O}$-module libre sur le faisceau d'ensembles gradués $\mathcal{S}$.
Plus précisément, la composante graduée de degré $n$ de
$\mathcal{O}[\mathcal{S}]$ est le faisceau associé au préfaisceau
$$
U \longmapsto
\bigoplus\nolimits_{s \in \mathcal{S}(U),\ \deg(s) = n} s \cdot \mathcal{O}(U)
$$
Avec le différentiel nul, nous pouvons aussi le considérer comme un
$\mathcal{O}$-module différentiel gradué.
Soit $\mathcal{A}$ un faisceau d'$\mathcal{O}$-algèbres graduées.
Nous définissons de même $\mathcal{A}[\mathcal{S}]$ comme le
$\mathcal{A}$-module gradué dont la composante graduée de degré $n$ est le
faisceau associé au préfaisceau
$$
U \longmapsto
\bigoplus\nolimits_{s \in \mathcal{S}(U)} s \cdot \mathcal{A}^{n - \deg(s)}(U)
$$
Si $\mathcal{A}$ est une $\mathcal{O}$-algèbre différentielle graduée,
nous en faisons un $\mathcal{O}$-module différentiel gradué
en posant $\text{d}(s) = 0$ pour tout $s \in \mathcal{S}(U)$,
puis en prenant le faisceau associé.
\end{remark}

\begin{lemma}
\label{lemma-free-graded-module-good}
Soit $(\mathcal{C}, \mathcal{O})$ un site annelé. Soit
$\mathcal{A}$ une $\mathcal{A}$-algèbre différentielle graduée. Soit
$\mathcal{S}$ un faisceau d'ensembles gradués sur $\mathcal{C}$.
Alors le module gradué libre $\mathcal{A}[\mathcal{S}]$
sur $\mathcal{S}$, muni du différentiel de la remarque
\ref{remark-sheaf-graded-sets}, est
un bon module différentiel gradué sur $\mathcal{A}$.
\end{lemma}

\begin{proof}
Soit $\mathcal{N}$ un module gradué à gauche sur $\mathcal{A}$.
Alors nous
avons $$
\mathcal{A}[\mathcal{S}] \otimes_\mathcal{A} \mathcal{N} =
\mathcal{O}[\mathcal{S}] \otimes_\mathcal{O} \mathcal{N} =
\mathcal{N}[\mathcal{S}]
$$
où $\mathcal{N}[\mathcal{S}$ est le module
gradué sur $\mathcal{O}$ dont la composante de degré $n$
est le faisceau associé au préfaisceau
$$ U
\longmapsto
\bigoplus\nolimits_{s \in \mathcal{S}(U)} s \cdot \mathcal{N}^{n - \deg(s)}(U)
$$ Il est
clair que $\mathcal{N} \to \mathcal{N}[\mathcal{S}]$ est un foncteur
exact ; ainsi, $\mathcal{A}[\mathcal{S}$ est plat comme module
gradué sur $\mathcal{A}$. Supposons maintenant que $\mathcal{N}$ soit un
module différentiel gradué à gauche sur $\mathcal{A}$. Alors nous avons
$$
H^*(\mathcal{A}[\mathcal{S}] \otimes_\mathcal{A} \mathcal{N}) =
H^*(\mathcal{O}[\mathcal{S}] \otimes_\mathcal{O} \mathcal{N})
$$
comme faisceaux gradués de $\mathcal{O}$-modules, ce qui, par platitude
(sur $\mathcal{O})$ est égal à
$$
H^*(\mathcal{N})[\mathcal{S}]
$$ comme module
gradué sur $\mathcal{O}$. Par conséquent, si $\mathcal{N}$ est acyclique,
alors $\mathcal{A}[\mathcal{S}] \otimes_\mathcal{A} \mathcal{N}$ est
acyclique.

\medskip\noindent
Enfin, considérons un morphisme
$(f, f^\sharp) : (\Sh(\mathcal{C}'), \mathcal{O}')
\to (\Sh(\mathcal{C}), \mathcal{O})$ de
topos annelés, une $\mathcal{O}'$-algèbre différentielle graduée
$\mathcal{A}'$, et une application $\varphi : f^{-1}\mathcal{A} \to \mathcal{A}'$
de $f^{-1}\mathcal{O}$-algèbres différentielles
graduées. Il est donc facile de voir que
$$
f^*\mathcal{A}[\mathcal{S}] = \mathcal{A}'[f^{-1}\mathcal{S}]
$$ ce
qui achève la démonstration que notre module est bon.
\end{proof}

\begin{lemma}
\label{lemma-resolve}
Soit $(\mathcal{C}, \mathcal{O})$ un site annelé.
Soit $\mathcal{A}$ un faisceau d'algèbres différentielles graduées
sur $(\mathcal{C}, \mathcal{O})$. Soit $\mathcal{M}$
un $\mathcal{A}$-module différentiel gradué. Il existe un homomorphisme
$\mathcal{P} \to \mathcal{M}$ de $\mathcal{A}$-modules différentiels gradués
avec les propriétés suivantes
\begin{enumerate}
\item $\mathcal{P} \to \mathcal{M}$ est un quasi-isomorphisme, et
\item $\mathcal{P}$ est bon.
\end{enumerate}
\end{lemma}

\begin{proof}[Première démonstration]
Soit $\mathcal{S}_0$ le faisceau d'ensembles
gradués (remarque \ref{remark-sheaf-graded-sets})
dont la composante de degré $n$ est $\Ker(\text{d}_\mathcal{M}^n)$.
Considérons l'homomorphisme de modules différentiels gradués
$$
\mathcal{P}_0 = \mathcal{A}[\mathcal{S}_0] \longrightarrow \mathcal{M}
$$ où
le membre de gauche est défini dans la remarque \ref{remark-sheaf-graded-sets} et
le morphisme envoie une section locale $s$ de $\mathcal{S}_0$ sur
la section locale correspondante de $\mathcal{M}^{\deg(s)}$ ; celle-ci
appartient au noyau du différentiel, ce qui assure que nous avons
bien un morphisme de modules différentiels gradués. Par construction, les
applications induites sur les faisceaux de cohomologie $H^n(\mathcal{P}_0) \to H^n(\mathcal{M})$
sont surjectives. Nous allons construire par récurrence
des applications $$
\mathcal{P}_0 \to \mathcal{P}_1 \to \mathcal{P}_2 \to \ldots \to \mathcal{M}
$$
Notons bien sûr que $H^*(\mathcal{P}_i) \to H^*(\mathcal{M})$ sera
surjectif pour tous les $i$. Étant
donné $\mathcal{P}_i \to \mathcal{M}$, notons $\mathcal{S}_{i + 1}$ le faisceau
d'ensembles gradués dont la composante de degré $n$
est $$
\Ker(\text{d}_{\mathcal{P}_i}^{n + 1})
\times_{\mathcal{M}^{n + 1}, \text{d}}
\mathcal{M}^n
$$ Posons
alors
$$
\mathcal{P}_{i + 1} = \mathcal{P}_i \oplus \mathcal{A}[\mathcal{S}_{i + 1}]
$$ comme module
gradué sur $\mathcal{A}$, avec le différentiel et le morphisme vers
$\mathcal{M}$ définis comme suit :
\begin{enumerate}
\item sur les sections locales de $\mathcal{P}_i$, on utilise le différentiel de
$\mathcal{P}_i$ et le morphisme donné vers $\mathcal{M}$,
\item pour une section locale $s = (p, m)$ de $\mathcal{S}_{i + 1}$, on pose
$\text{d}(s)$ égal à $p$, vu comme une section de $\mathcal{P}_i$ de degré
$\deg(s) + 1$, et on envoie $s$ sur $m$ dans $\mathcal{M}$, et
\item on prolonge le différentiel de manière unique de façon que la règle de Leibniz soit satisfaite.
\end{enumerate}
Cette construction est bien définie, car $\text{d}(m)$ est l'image de $p$
et $\text{d}(p) = 0$.
Enfin, posons $\mathcal{P} = \colim \mathcal{P}_i$, muni du
morphisme induit vers $\mathcal{M}$.

\medskip\noindent
L'application $\mathcal{P} \to \mathcal{M}$ est un quasi-isomorphisme :
nous avons $H^n(\mathcal{P}) = \colim H^n(\mathcal{P}_i)$
et pour chaque $i$ l'application $H^n(\mathcal{P}_i) \to H^n(\mathcal{M})$
est surjective avec le noyau anéanti par l'application
$H^n(\mathcal{P}_i) \to H^n(\mathcal{P}_{i + 1})$ par construction.
Chaque $\mathcal{P}_i$ est bon parce que $\mathcal{P}_0$ est bon
par le lemme \ref{lemma-free-graded-module-good} et chaque $\mathcal{P}_{i + 1}$
est au milieu de la suite exacte courte admissible
$0 \to \mathcal{P}_i \to \mathcal{P}_{i + 1} \to
\mathcal{A}[\mathcal{S}_{i + 1}] \to 0$ dont les termes
extérieurs sont bons par induction. Par conséquent,
$\mathcal{P}_{i + 1}$ est bon par
le lemme \ref{lemma-good-admissible-ses}.
Enfin, nous concluons que $\mathcal{P}$ est bon par
le lemme \ref{lemma-good-direct-sum}.
\end{proof}

\begin{proof}[Seconde démonstration]
Nous invitons le lecteur à lire la preuve d'Algèbre différentielle
graduée, lemme \ref{dga-lemma-resolve} avant de lire cette preuve.
Posons $\mathcal{M} = \mathcal{M}_0$.
Nous choisissons par récurrence des suites
exactes courtes $$
0 \to \mathcal{M}_{i + 1} \to \mathcal{P}_i \to \mathcal{M}_i \to
0 $$
où
les applications $\mathcal{P}_i \to \mathcal{M}_i$ sont
choisies comme dans le lemme \ref{lemma-good-quotient}.
Cela
fournit une « résolution » $$
\ldots \to \mathcal{P}_2 \xrightarrow{f_2}
\mathcal{P}_1 \xrightarrow{f_1} \mathcal{P}_0 \to \mathcal{M} \to
0 $$
Définissons alors $\mathcal{P}$ comme le
module différentiel gradué sur $\mathcal{A}$
suivant \begin{enumerate}
\item comme $\mathcal{A}$-module gradué,
posons $\mathcal{P} = \bigoplus_{a \leq 0} \mathcal{P}_{-a}[-a]$ ; sa
composante de degré $n$ est donc
$\mathcal{P}^n = \bigoplus\nolimits_{a + b = n} \mathcal{P}_{-a}^b$,
\item le différentiel sur $\mathcal{P}$ est, comme dans la construction du complexe
total associé à un double complexe, donné par
$$
\text{d}_\mathcal{P}(x) = f_{-a}(x) + (-1)^a \text{d}_{\mathcal{P}_{-a}}(x)
$$ pour
$x$ une section locale de $\mathcal{P}_{-a}^b$.
\end{enumerate} Avec ces conventions
$\mathcal{P}$ est en effet un
$\mathcal{A}$-module différentiel gradué ; nous omettons les détails. Il existe un morphisme
$\mathcal{P} \to \mathcal{M}$ de
$\mathcal{A}$-modules différentiels gradués qui est nul sur les facteurs directs $\mathcal{P}_{-a}[-a]$ pour
$a < 0$ et égal au morphisme donné $\mathcal{P}_0 \to \mathcal{M}$ pour $a = 0$. Notons que
nous avons
$$
\mathcal{P} = \colim_i F_i\mathcal{P}
$$ où
$F_i\mathcal{P} \subset \mathcal{P}$ est le sous-module différentiel gradué sur
$\mathcal{A}$ dont le module gradué sous-jacent sur $\mathcal{A}$
est
$$
F_i\mathcal{P} =
\bigoplus\nolimits_{i \geq -a \geq 0} \mathcal{P}_{-a}[-a]
$$ Il est immédiat
que les applications
$$ 0
\to F_1\mathcal{P} \to F_2\mathcal{P} \to F_3\mathcal{P} \to
\ldots \to \mathcal{P}
$$ sont toutes des
monomorphismes admissibles, et nous avons des suites exactes courtes admissibles
$$ 0
\to F_i\mathcal{P} \to
F_{i + 1}\mathcal{P} \to
\mathcal{P}_{i + 1}[i + 1] \to 0
$$ Par
récurrence et par le lemme \ref{lemma-good-admissible-ses}, nous
obtenons que $F_i\mathcal{P}$ est un bon module différentiel
gradué sur $\mathcal{A}$. Puisque $\mathcal{P} = \colim F_i\mathcal{P}$ nous
trouvons que $\mathcal{P}$ est bon par le lemme \ref{lemma-good-direct-sum}.

\medskip\noindent
Enfin, nous devons démontrer que $\mathcal{P} \to \mathcal{M}$ est
un quasi-isomorphisme. Si $\mathcal{C}$ a suffisamment de points, cela
découle du lemme élémentaire d'Homologie
\ref{homology-lemma-good-resolution-gives-qis}, en vérifiant sur
les fibres. En général, on peut raisonner comme suit (cette
démonstration est beaucoup trop longue ; on peut aussi utiliser les sections
locales, comme dans la démonstration élémentaire, mais cet argument est
également assez long).
Puisque les limites inductives filtrantes sont exactes dans la catégorie des
faisceaux abéliens,
nous avons $$
H^d(\mathcal{P}) = \colim H^d(F_i\mathcal{P})
$$
Nous affirmons que, pour tout $i \geq 0$ et tout $d \in \mathbf{Z}$,
nous avons (a) une
suite exacte courte $$
0 \to H^d(\mathcal{M}_{i + 1}[i]) \to
H^d(F_i\mathcal{P}) \to H^d(\mathcal{M}) \to
0 $$
où la deuxième flèche
provient de $F_i\mathcal{P} \to \mathcal{P} \to \mathcal{M}$
et b) la composition
$$
H^d(\mathcal{M}_{i + 1}[i]) \to
H^d(F_i\mathcal{P}) \to
H^d(F_{i + 1}\mathcal{P})
$$ est
nulle. Cette assertion suffit manifestement à achever la démonstration.

\medskip\noindent
Démontrons l'assertion. Pour tout $i \geq 0$, il existe une
application $\mathcal{M}_{i + 1}[i] \to F_i\mathcal{P}$
provenant de l'inclusion de $\mathcal{M}_{i + 1}$ dans $\mathcal{P}_i$
en tant que noyau de $f_i$. Considérons la suite exacte
courte $$
0 \to \mathcal{M}_{i + 1}[i] \to F_i\mathcal{P} \to C_i \to
0 $$
de complexes de modules sur $\mathcal{O}$ définissant $C_i$.
Remarquons que $C_0 = \mathcal{M}_0 = \mathcal{M}$.
Observons de plus que $C_i$ est le complexe total associé
au double complexe $C_i^{\bullet, \bullet}$ avec des
colonnes $$
\mathcal{M}_i = \mathcal{P}_i/\mathcal{M}_{i + 1},
\mathcal{P}_{i - 1}, \ldots, \mathcal{P}_0
$$
en degrés $-i, -i + 1, \ldots, 0$. Il existe un morphisme de doubles
complexes $C_i^{\bullet, \bullet} \to C_{i - 1}^{\bullet, \bullet}$ qui
est $0$ sur la colonne en degré $-i$, est la surjection
$\mathcal{P}_{i - 1} \to \mathcal{M}_{i - 1}$ en degré $-i + 1$, et
est l'identité sur les autres colonnes. Nous
obtenons donc des morphismes de
complexes $$
C_i \longrightarrow C_{i - 1}
$$
Ces morphismes sont des quasi-isomorphismes surjectifs : leur
noyau est le complexe total du double complexe ayant pour
colonnes $\mathcal{M}_i, \mathcal{M}_i$ en degrés $-i, -i + 1$,
avec l'identité entre ces deux colonnes. En utilisant
les identifications qui en
résultent, $H^d(C_i) = H^d(C_{i - 1} = \ldots = H^d(\mathcal{M})$,
nous obtenons déjà une suite exacte longue
$$
H^d(\mathcal{M}_{i + 1}[i]) \to
H^d(F_i\mathcal{P}) \to H^d(\mathcal{M}) \to
H^{d + 1}(\mathcal{M}_{i + 1}[i])
$$
de la suite exacte courte de complexes ci-dessus.
Cependant, nous avons aussi le diagramme commutatif
$$
\xymatrix{
\mathcal{M}_{i + 2}[i + 1] \ar[r]_a & T_{i + 1} \ar[r] &
F_{i + 1}\mathcal{P} \ar[r] & C_{i + 1} \ar[d] \\ &
\mathcal{M}_{i + 1}[i] \ar[r] \ar[u]^b &
F_i\mathcal{P} \ar[u] \ar[r] &
C_i
}
$$ où
$T_{i + 1}$ est le complexe total du double complexe avec les
colonnes $\mathcal{P}_{i + 1}, \mathcal{M}_{i + 1}$ placées en degrés
$-i - 1$ et $-i$. Autrement dit, $T_{i + 1}$ est un décalage du
cône du morphisme
$\mathcal{P}_{i + 1} \to \mathcal{M}_{i + 1}$ ; nous voyons que
$a$ est un quasi-isomorphisme et que $a^{-1} \circ b$ est un décalage du troisième morphisme du
triangle distingué de $D(\mathcal{O})$ associé à la suite exacte
courte
$$ 0
\to \mathcal{M}_{i + 2} \to \mathcal{P}_{i + 1} \to \mathcal{M}_{i + 1} \to 0
$$ L'application
$H^d(\mathcal{P}_{i + 1}) \to H^d(\mathcal{M}_{i + 1})$ est surjective parce que nous
avons choisi nos applications de telle sorte que
$\Ker(\text{d}_{\mathcal{P}_{i + 1}}) \to
\Ker(\text{d}_{\mathcal{M}_{i + 1}})$ est surjective. Ainsi,
$a^{-1} \circ b$ induit l'application nulle sur les faisceaux de cohomologie.
Cela prouve la partie b) de
l'assertion. Puisque $T_{i + 1}$ est le noyau du morphisme surjectif de
complexes $F_{i + 1}\mathcal{P} \to C_i$, nous obtenons un
morphisme de suites exactes
longues de cohomologie $$
\xymatrix{
H^d(T_{i + 1}) \ar[r]
& H^d(F_{i + 1}\mathcal{P}) \ar[r]
& H^d(\mathcal{M}) \ar[r]
& H^{d + 1}(T_{i + 1}) \\
H^d(\mathcal{M}_{i + 1}[i]) \ar[r] \ar[u]
& H^d(F_i\mathcal{P}) \ar[r] \ar[u]
& H^d(\mathcal{M}) \ar[r] \ar[u]
& H^{d + 1}(\mathcal{M}_{i + 1}[i]) \ar[u]
}
$$
D'après la discussion ci-dessus, les flèches
verticales extérieures sont nulles. Ainsi, les
applications $H^d(F_{i + 1}\mathcal{P}) \to H^d(\mathcal{M})$
sont surjectives, ce qui prouve la partie
(a) de l'assertion. Plus précisément, cela prouve l'assertion pour $i > 0$ ;
nous laissons le cas $i = 0$ au lecteur (il suffit en
fait de traiter tous les $i \gg 0$ pour obtenir le lemme).
\end{proof}

\begin{lemma}
\label{lemma-acyclic-good}
Soit $(\mathcal{C}, \mathcal{O})$ un site annelé. Soit
$\mathcal{A}$ un faisceau d'algèbres différentielles graduées sur
$(\mathcal{C}, \mathcal{O})$. Soit $\mathcal{P}$ un bon module
différentiel gradué acyclique à droite sur $\mathcal{A}$.
\begin{enumerate}
\item pour tout module différentiel gradué à gauche sur $\mathcal{A}$,
$\mathcal{N}$, le produit tensoriel
$\mathcal{P} \otimes_\mathcal{A} \mathcal{N}$ est acyclique,
\item pour tout morphisme $(f, f^\sharp) : (\Sh(\mathcal{C}'), \mathcal{O}')
\to (\Sh(\mathcal{C}), \mathcal{O})$ de
topos annelés, toute $\mathcal{O}'$-algèbre différentielle graduée
$\mathcal{A}'$ et tout morphisme $\varphi : f^{-1}\mathcal{A} \to \mathcal{A}'$ de
$f^{-1}\mathcal{O}$-algèbres différentielles graduées,
l'image inverse $f^*\mathcal{P}$ est acyclique et bonne.
\end{enumerate}
\end{lemma}

\begin{proof}
Démontrons (1). Par le lemme \ref{lemma-resolve} nous pouvons choisir
un bon module différentiel gradué à gauche $\mathcal{Q}$ et un
quasi-isomorphisme $\mathcal{Q} \to \mathcal{N}$.
Alors $\mathcal{P} \otimes_\mathcal{A} \mathcal{Q}$
est acyclique parce que $\mathcal{Q}$
est bon. Soit $\mathcal{N}'$ le cône du
morphisme $\mathcal{Q} \to \mathcal{N}$.
Ensuite, $\mathcal{P} \otimes_\mathcal{A} \mathcal{N}'$ est acyclique
parce que $\mathcal{P}$ est bon et parce que $\mathcal{N}'$ est
acyclique (comme le cône sur un quasi-isomorphisme). Nous avons un
triangle
distingué $$
\mathcal{Q} \to \mathcal{N} \to \mathcal{N}' \to \mathcal{Q}[1]
$$
dans $K(\textit{Mod}(\mathcal{A}, \text{d}))$ par notre
construction de la structure triangulée. Puisque
$\mathcal{P} \otimes_\mathcal{A} -$ envoie les triangles
distingués sur des triangles distingués, nous obtenons un triangle
distingué
$$
\mathcal{P} \otimes_\mathcal{A} \mathcal{Q} \to
\mathcal{P} \otimes_\mathcal{A} \mathcal{N} \to
\mathcal{P} \otimes_\mathcal{A} \mathcal{N}' \to
\mathcal{P} \otimes_\mathcal{A} \mathcal{Q}[1]
$$ dans
$K(\textit{Mod}(\mathcal{O}))$. C'est ainsi que nous concluons.

\medskip\noindent
Démontrons (2). Notons que $f^*\mathcal{P}$ est bon selon notre définition de
bons modules. Rappelons
que $f^*\mathcal{P} = f^{-1}\mathcal{P} \otimes_{f^{-1}\mathcal{A}} \mathcal{A}'$.
Alors $f^{-1}\mathcal{P}$ est un bon module différentiel gradué acyclique (car $f^{-1}$
est exact) sur $f^{-1}\mathcal{A}$. On voit donc
que $f^*\mathcal{P}$ est acyclique par la partie (1).
\end{proof}






\section{L'enveloppe différentielle graduée d'un module gradué}
\label{section-dg-hull}

\noindent
L'enveloppe différentielle graduée d'un module gradué $\mathcal{N}$
est le résultat de l'application du foncteur $G$ du lemme suivant.

\begin{lemma}
\label{lemma-dg-hull}
Soit $(\mathcal{C}, \mathcal{O})$ un site annelé.
Soit $\mathcal{A}$ un faisceau d'algèbres différentielles graduées
sur $(\mathcal{C}, \mathcal{O})$. Le foncteur d'oubli
$F : \textit{Mod}(\mathcal{A}, \text{d}) \to \textit{Mod}(\mathcal{A})$ a
un adjoint à gauche $G : \textit{Mod}(\mathcal{A}) \to
\textit{Mod}(\mathcal{A}, \text{d})$.
\end{lemma}

\begin{proof}
Pour prouver l'existence de $G$, nous pouvons utiliser le théorème du foncteur
adjoint, voir Catégories, théorème \ref{categories-theorem-adjoint-functor} (noter que
nous avons interverti les rôles de $F$ et $G$). Les conditions d'exactitude sur
$F$ sont satisfaites par le lemme \ref{lemma-dgm-abelian}. La condition
ensembliste se vérifie comme suit : supposons donné un $\mathcal{A}$-module gradué
$\mathcal{N}$. Pour toute application
$$
\varphi : \mathcal{N} \longrightarrow F(\mathcal{M})
$$ nous
pouvons considérer le plus petit sous-module différentiel gradué sur $\mathcal{A}$,
$\mathcal{M}' \subset \mathcal{M}$, tel
que $\Im(\varphi) \subset F(\mathcal{M}')$.
Il est clair que $\mathcal{M}'$ est l'image de l'application
de $\mathcal{A}$-modules gradués
$$
\mathcal{N} \oplus
\mathcal{N}[-1] \otimes_\mathcal{O} \mathcal{A}
\longrightarrow
\mathcal{M}
$$ définie
par
$$ (n,
\sum n_i \otimes a_i) \longmapsto
\varphi(n) + \sum \text{d}(\varphi(n_i)) a_i
$$ car on vérifie
aisément que l'image de cette application est un sous-module différentiel
gradué de $\mathcal{M}$. Le nombre
de classes d'isomorphisme possibles de ces $\mathcal{M}'$ est donc borné, ce
qui conclut la démonstration.
\end{proof}

\noindent
Soit $(\mathcal{C}, \mathcal{O})$ un site annelé. Soit
$\mathcal{A}$ un faisceau d'algèbres différentielles graduées sur
$(\mathcal{C}, \mathcal{O})$. Soit $\mathcal{M}$ un
$\mathcal{A}$-module différentiel gradué et supposons donnée une suite exacte courte
$$ 0
\to \mathcal{N} \to F(\mathcal{M}) \to \mathcal{N}' \to 0
$$ dans
$\textit{Mod}(\mathcal{A})$. Nous obtenons alors un homomorphisme canonique
de modules gradués sur $\mathcal{A}$
$$
\overline{\text{d}} : \mathcal{N} \to \mathcal{N}'[1]
$$ de la manière
suivante : pour une section locale $x$ de $\mathcal{N}$, notons
$\overline{\text{d}}(x)$ l'image dans $\mathcal{N}'$ de
$\text{d}_\mathcal{M}(x)$, où $x$ est considéré comme une section locale
de $\mathcal{M}$.

\begin{lemma}
\label{lemma-dg-hull-acyclic}
Les foncteurs $F, G$ du lemme \ref{lemma-dg-hull} ont les
propriétés suivantes. Pour un $\mathcal{A}$-module gradué
$\mathcal{N}$, nous avons
\begin{enumerate}
\item l'unité $\mathcal{N} \to F(G(\mathcal{N}))$ est injective,
\item l'application $\overline{\text{d}} : \mathcal{N} \to
\Coker(\mathcal{N} \to F(G(\mathcal{N})))[1]$ est un isomorphisme, et
\item $G(\mathcal{N})$ est un $\mathcal{A}$-module différentiel gradué acyclique.
\end{enumerate}
\end{lemma}

\begin{proof}
On observe que la propriété (3) est une conséquence des propriétés (1) et (2).
En effet, si $s$ est une section locale non nulle de $F(G(\mathcal{N}))$
telle que $\text{d}(s) = 0$, alors $s$ ne peut pas
appartenir à l'image de $\mathcal{N} \to F(G(\mathcal{N}))$. Ainsi, nous
pouvons écrire l'image $\overline{s}$ de $s$
dans le conoyau comme $\overline{\text{d}}(s')$ pour une section locale $s'$ de $\mathcal{N}$.
Alors $s = \text{d}(s')$, car la différence
$s - \text{d}(s')$ appartient encore au noyau de $\text{d}$
et est contenue dans l'image de l'unité.

\medskip\noindent
Notons provisoirement $\mathcal{A}_{gr}$, respectivement $\mathcal{A}_{dg}$, le
faisceau $\mathcal{A}$ considéré comme un module gradué à droite sur
lui-même, respectivement comme un module différentiel gradué à droite
sur lui-même. Le cas le plus important du lemme est de comprendre
ce qu'est $G(\mathcal{A}_{gr})$. Bien sûr, $G(\mathcal{A}_{gr})$ est
l'objet de $\textit{Mod}(\mathcal{A}, \text{d})$ représentant
le
foncteur $$
\mathcal{M} \longmapsto
\Hom_{\textit{Mod}(\mathcal{A})}(\mathcal{A}_{gr}, F(\mathcal{M}))
= \Gamma(\mathcal{C}, \mathcal{M})
$$
Par la remarque \ref{remark-cone-identity}, nous voyons que ce
foncteur est représenté par $C[-1]$ où $C$ est le cône sur l'identité de $\mathcal{A}_{dg}$.
Nous avons une
suite exacte courte $$
0 \to \mathcal{A}_{dg}[-1] \to C[-1] \to \mathcal{A}_{dg} \to
0 $$
dans $\textit{Mod}(\mathcal{A}, \text{d})$,
scindée par l'unité $\mathcal{A}_{gr} \to F(C[-1])$ dans $\textit{Mod}(\mathcal{A})$.
Ainsi, $G(\mathcal{A}_{gr})$ vérifie (1) et (2).

\medskip\noindent
Soit $U$ un objet de $\mathcal{C}$. Notons
$j_U : \mathcal{C}/U \to \mathcal{C}$ le morphisme de localisation.
Notons $\mathcal{A}_U$ la restriction de $\mathcal{A}$ à $U$.
Nous noterons $\mathcal{A}_{U, gr}$ le faisceau
$\mathcal{A}_U$ considéré comme un module gradué sur $\mathcal{A}_U$.
Notons $F_U : \textit{Mod}(\mathcal{A}_U, \text{d}) \to
\textit{Mod}(\mathcal{A}_U)$ le foncteur d'oubli et
$G_U$ son adjoint. Nous avons alors les diagrammes commutatifs
$$
\vcenter{
\xymatrix{
\textit{Mod}(\mathcal{A}, \text{d}) \ar[d]_{j_U^*} \ar[r]_F &
\textit{Mod}(\mathcal{A}) \ar[d]^{j_U^*} \\
\textit{Mod}(\mathcal{A}_U, \text{d}) \ar[r]^{F_U} &
\textit{Mod}(\mathcal{A}_U)
}
}
\quad\text{et}\quad
\vcenter{
\xymatrix{
\textit{Mod}(\mathcal{A}_U, \text{d}) \ar[r]_{F_U} \ar[d]_{j_{U!}}
& \textit{Mod}(\mathcal{A}_U) \ar[d]^{j_{U!}} \\
\textit{Mod}(\mathcal{A}, \text{d}) \ar[r]^F
& \textit{Mod}(\mathcal{A})
}
}
$$
par la construction de $j^*_U$ et $j_{U!}$
dans les sections \ref{section-functoriality-graded},
\ref{section-functoriality-dg},
\ref{section-localize-graded}
et \ref{section-localize-dg}.
Par l'unicité des adjoints, on obtient $j_{U!} \circ G_U = G \circ j_{U!}$.
Puisque $j_{U!}$ est exact, les propriétés (1)
et (2) pour l'unité
$\mathcal{A}_{U, gr} \to F_U(G_U(\mathcal{A}_{U, gr}))$, établies
dans la partie précédente de la démonstration,
impliquent les mêmes propriétés pour l'unité
$j_{U!}\mathcal{A}_{U, gr} \to F(G(j_{U!}\mathcal{A}_{U, gr})) =
j_{U!}F_U(G_U(\mathcal{A}_{U, gr}))$.

\medskip\noindent Dans
la démonstration du lemme \ref{lemma-gm-grothendieck-abelian},
nous avons vu que tout objet de
$\textit{Mod}(\mathcal{A})$ est un quotient d'une somme directe de copies
de $j_{U!}\mathcal{A}_{U, gr}$. Puisque $G$ est un adjoint à gauche,
nous voyons que $G$ commute aux sommes directes. Ainsi, les
propriétés (1) et (2) sont vérifiées pour les sommes directes
d'objets qui les vérifient. Par conséquent, tout objet $\mathcal{N}$ de
$\textit{Mod}(\mathcal{A})$ s'insère dans une suite exacte
$$
\mathcal{N}_1 \to \mathcal{N}_0 \to \mathcal{N} \to 0
$$ telle que les propriétés
(1) et (2) soient vérifiées pour $\mathcal{N}_1$ et $\mathcal{N}_0$. Nous
laissons au lecteur le soin de déduire (1) et (2) pour
$\mathcal{N}$ en utilisant que $G$ est exact à droite.
\end{proof}









\section{Modules différentiels gradués K-injectifs}
\label{section-K-injective}

\noindent
Cette section reprend les résultats
d'Injectifs, section \ref{injectives-section-K-injective},
dans le cadre des faisceaux de modules différentiels gradués sur
un faisceau d'algèbres différentielles graduées.

\begin{lemma}
\label{lemma-characterize-injectives} Soit
$(\mathcal{C}, \mathcal{O})$ un site annelé. Soit $\mathcal{A}$ un faisceau
d'algèbres graduées sur $(\mathcal{C}, \mathcal{O})$. Il existe un
ensemble $T$ et, pour chaque $t \in T$, une application injective
$\mathcal{N}_t \to \mathcal{N}'_t$ de $\mathcal{A}$-modules gradués, tels qu'un
objet $\mathcal{I}$ de $\textit{Mod}(\mathcal{A})$ soit injectif si et
seulement si, pour tout diagramme en traits pleins
$$
\xymatrix{
\mathcal{N}_t \ar[r] \ar[d] & \mathcal{I} \\
\mathcal{N}'_t \ar@{..>}[ru]
}
$$ il existe
une flèche pointillée dans $\textit{Mod}(\mathcal{A})$ rendant le diagramme commutatif.
\end{lemma}

\begin{proof}
Cet énoncé vaut dans toute catégorie abélienne de Grothendieck ;
voir Injectifs, lemme \ref{injectives-lemma-characterize-injective}.
Par le lemme \ref{lemma-gm-grothendieck-abelian}, la catégorie
$\textit{Mod}(\mathcal{A})$ est une catégorie abélienne de Grothendieck.
\end{proof}

\begin{definition}
\label{definition-graded-injective}
Soit $(\mathcal{C}, \mathcal{O})$ un site annelé. Soit
$(\mathcal{A}, \text{d})$ un faisceau d'algèbres différentielles graduées sur
$(\mathcal{C}, \mathcal{O})$. Un $\mathcal{A}$-module différentiel gradué
$\mathcal{I}$ est dit {\it injectif gradué}\footnote{Cette terminologie n'est
peut-être pas standard.} si $\mathcal{M}$, considéré comme un
$\mathcal{A}$-module gradué, est un objet injectif de la catégorie
$\textit{Mod}(\mathcal{A})$ des $\mathcal{A}$-modules gradués.
\end{definition}

\begin{remark}
\label{remark-why-graded-injective}
Soit $(\mathcal{C}, \mathcal{O})$ un site annelé. Soit
$(\mathcal{A}, \text{d})$ un faisceau d'algèbres différentielles graduées sur
$(\mathcal{C}, \mathcal{O})$. Soit $\mathcal{I}$ un
$\mathcal{A}$-module différentiel gradué injectif gradué. Soit
$$ 0
\to \mathcal{M}_1 \to \mathcal{M}_2 \to \mathcal{M}_3 \to 0
$$ une
suite exacte courte de $\mathcal{A}$-modules différentiels gradués. Puisque
$\mathcal{I}$ est injectif gradué, nous
obtenons une suite exacte courte de complexes
$$ 0
\to
\Hom_{\textit{Mod}^{dg}(\mathcal{A}, \text{d})}(\mathcal{M}_3, \mathcal{I})
\to
\Hom_{\textit{Mod}^{dg}(\mathcal{A}, \text{d})}(\mathcal{M}_2, \mathcal{I})
\to
\Hom_{\textit{Mod}^{dg}(\mathcal{A}, \text{d})}(\mathcal{M}_1, \mathcal{I})
\to 0
$$ de modules
sur $\Gamma(\mathcal{C}, \mathcal{O})$. En passant à la cohomologie,
nous obtenons une suite exacte
longue $$
\xymatrix{
\Hom_{K(\textit{Mod}(\mathcal{A}, \text{d}))}(\mathcal{M}_3, \mathcal{I})
\ar[d]
& \Hom_{K(\textit{Mod}(\mathcal{A}, \text{d}))}(\mathcal{M}_3, \mathcal{I})[1]
\ar[d] \\
\Hom_{K(\textit{Mod}(\mathcal{A}, \text{d}))}(\mathcal{M}_2, \mathcal{I})
\ar[d] &
\Hom_{K(\textit{Mod}(\mathcal{A}, \text{d}))}(\mathcal{M}_2, \mathcal{I})[1]
\ar[d] \\
\Hom_{K(\textit{Mod}(\mathcal{A}, \text{d}))}(\mathcal{M}_1, \mathcal{I})
\ar[ruu]
&
\Hom_{K(\textit{Mod}(\mathcal{A}, \text{d}))}(\mathcal{M}_1, \mathcal{I})[1]
}
$$ de
groupes d'homomorphismes dans la catégorie homotopique. Le point est que nous
obtenons cela même si nous n'avons pas supposé que notre suite exacte courte
est admissible (donc la suite exacte courte en général ne définit pas
un triangle distingué dans la catégorie homotopique).
\end{remark}

\begin{lemma}
\label{lemma-product-graded-injective}
Soit $(\mathcal{C}, \mathcal{O})$ un site annelé. Soit
$(\mathcal{A}, \text{d})$ un faisceau d'algèbres différentielles graduées sur
$(\mathcal{C}, \mathcal{O})$. Soit $T$ un ensemble et, pour
chaque $t \in T$, soit $\mathcal{I}_t$ un module différentiel
gradué injectif gradué sur $\mathcal{A}$. Alors
$\prod \mathcal{I}_t$ est un
$\mathcal{A}$-module différentiel gradué injectif gradué.
\end{lemma}

\begin{proof}
Cela résulte du fait que les produits d'objets injectifs sont injectifs,
voir Homologie, lemme \ref{homology-lemma-product-injectives},
et que les produits
dans $\textit{Mod}(\mathcal{A}, \text{d})$ correspondent
aux produits dans $\textit{Mod}(\mathcal{A})$ par le foncteur d'oubli.
\end{proof}

\begin{lemma}
\label{lemma-characterize-graded-injectives-in-dg} Soit
$(\mathcal{C}, \mathcal{O})$ un site annelé. Soit
$(\mathcal{A}, \text{d})$ un faisceau
d'algèbres différentielles graduées sur $(\mathcal{C}, \mathcal{O})$. Il existe un ensemble
$T$ et, pour chaque $t \in T$, une application injective
$\mathcal{M}_t \to \mathcal{M}'_t$ de
$\mathcal{A}$-modules différentiels gradués acycliques, tels que, pour
un objet $\mathcal{I}$ de $\textit{Mod}(\mathcal{A}, \text{d})$, les propriétés suivantes
soient équivalentes :
\begin{enumerate}
\item $\mathcal{I}$ est injectif gradué, et
\item pour chaque diagramme en traits pleins
$$
\xymatrix{
\mathcal{M}_t \ar[r] \ar[d] & \mathcal{I} \\
\mathcal{M}'_t \ar@{..>}[ru]
}
$$ il
existe une flèche pointillée dans $\textit{Mod}(\mathcal{A}, \text{d})$
rendant le diagramme commutatif.
\end{enumerate}
\end{lemma}

\begin{proof}
Soient $T$ et $\mathcal{N}_t \to \mathcal{N}'_t$ comme dans
le lemme \ref{lemma-characterize-injectives}.
Notons $F : \textit{Mod}(\mathcal{A}, \text{d}) \to \textit{Mod}(\mathcal{A})$ le
foncteur d'oubli. Soit
$G$ l'adjoint à gauche de
$F$ donné par le lemme \ref{lemma-dg-hull}. Posons
$$
\mathcal{M}_t = G(\mathcal{N}_t) \to
G(\mathcal{N}'_t) = \mathcal{M}'_t
$$ Il s'agit d'un
homomorphisme injectif de modules différentiels gradués acycliques
sur $\mathcal{A}$, par le lemme \ref{lemma-dg-hull-acyclic}.
Puisque $G$ est l'adjoint à gauche de $F$ nous voyons
qu'il existe une flèche pointillée dans le
diagramme $$
\xymatrix{
\mathcal{M}_t \ar[r] \ar[d] & \mathcal{I} \\
\mathcal{M}'_t \ar@{..>}[ru]
}
$$
si et seulement s'il existe une flèche pointillée dans
le diagramme $$
\xymatrix{
\mathcal{N}_t \ar[r] \ar[d] & F(\mathcal{I}) \\
\mathcal{N}'_t \ar@{..>}[ru]
}
$$
L'assertion résulte donc du choix de
notre collection de flèches $\mathcal{N}_t \to \mathcal{N}_t'$.
\end{proof}


\begin{lemma}
\label{lemma-small-acyclics}
Soit $(\mathcal{C}, \mathcal{O})$ un site annelé. Soit
$(\mathcal{A}, \text{d})$ un faisceau d'algèbres différentielles graduées sur
$(\mathcal{C}, \mathcal{O})$. Il existe un ensemble $S$ et, pour chaque $s$,
un $\mathcal{A}$-module différentiel gradué acyclique $\mathcal{M}_s$, tels
que, pour tout $\mathcal{A}$-module différentiel gradué acyclique non nul
$\mathcal{M}$, il existe $s \in S$ et une application injective
$\mathcal{M}_s \to \mathcal{M}$ dans $\textit{Mod}(\mathcal{A}, \text{d})$.
\end{lemma}

\begin{proof}
Avant de commencer, rappelons que nos conventions garantissent que le site $\mathcal{C}$
a un ensemble d'objets et de morphismes et un ensemble $\text{Cov}(\mathcal{C})$
de
couvertures. Si $\mathcal{F}$
est un $\mathcal{A}$-module différentiel gradué, définissons $|\mathcal{F}|$
comme le cardinal de
$$
\coprod\nolimits_{(U, n)} \mathcal{F}^n(U)
$$
lorsque $U$ parcourt les objets de $\mathcal{C}$ et $n \in \mathbf{Z}$.
Choisissons un cardinal infini $\kappa$ plus grand
que les cardinaux $|\Ob(\mathcal{C})|$, $|\text{Arrows}(\mathcal{C})|$,
$|\text{Cov}(\mathcal{C})|$, $\sup |I|$
pour $\{U_i \to U\}_{i \in I} \in \text{Cov}(\mathcal{C})$
et $|\mathcal{A}|$.

\medskip\noindent
Soit $\mathcal{F} \subset \mathcal{M}$ une inclusion de modules
différentiels gradués sur $\mathcal{A}$.
Supposons donnés un ensemble $K$ et, pour chaque $k \in K$, un
triplet $(U_k, n_k, x_k)$ constitué d'un objet $U_k$ de $\mathcal{C}$,
d'un entier $n_k$ et d'une section $x_k \in \mathcal{M}^{n_k}(U_k)$.
Nous pouvons alors considérer le plus petit sous-module
différentiel gradué sur $\mathcal{A}$, $\mathcal{F}' \subset \mathcal{M}$,
contenant $\mathcal{F}$ et les sections $x_k$ pour $k \in K$.
On peut
décrire $$
(\mathcal{F}')^n(U) \subset \mathcal{M}^n(U)
$$
comme l'ensemble des éléments $x \in \mathcal{M}^n(U)$ pour lesquels
il existe un recouvrement $\{f_i : U_i \to U\}_{i \in I} \in \text{Cov}(\mathcal{C})$
tel que, pour tout $i \in I$, il existe un ensemble fini $T_i$
et des morphismes $g_t : U_i \to U_{k_t}$
vérifiant $$
f_i^*x
= y_i + \sum\nolimits_{t \in T_i} a_{it}g_t^*x_{k_t} + b_{it}g_t^*\text{d}(x_{k_t})
$$
pour une section $y_i \in \mathcal{F}^n(U)$
et des sections $a_{it} \in \mathcal{A}^{n - n_{k_t}}(U_i)$
et $b_{it} \in \mathcal{A}^{n - n_{k_t} - 1}(U_i)$.
(Détails omis ; indication : ces sections appartiennent à $\mathcal{F}'$
et, réciproquement, on vérifie que cette règle définit
un sous-module différentiel gradué sur $\mathcal{A}$.)
Cette description
donne $|\mathcal{F}'| \leq \max(|\mathcal{F}|, |K|, \kappa)$.

\medskip\noindent
Soit $\mathcal{M}$ un
$\mathcal{A}$-module différentiel gradué acyclique non nul. Alors il existe un entier $n$ et une section
non nulle $x$ de $\mathcal{M}^n$ sur un objet $U$ de
$\mathcal{C}$. Soit
$$
\mathcal{F}_0 \subset \mathcal{M}
$$ le plus
petit sous-module différentiel gradué sur $\mathcal{A}$ contenant
$x$. Par le paragraphe précédent, nous avons
$|\mathcal{F}_0| \leq \kappa$. Par récurrence, étant donnés
$\mathcal{F}_0, \ldots, \mathcal{F}_n$, définissons
$\mathcal{F}_{n + 1}$ comme suit. Considérons
l'ensemble $$
L = \{(U, n, x)\}
\{U_i \to U\}_{i \in I}, (x_i)_{i \in I})\}
$$
des triples où $U$ est un objet de $\mathcal{C}$, $n \in \mathbf{Z}$
et $x \in \mathcal{F}_n(U)$ avec $\text{d}(x) = 0$. Puisque
$\mathcal{M}$ est acyclique pour chaque triple $l = (U_l, n_l, x_l) \in L$
nous pouvons
choisir $\{(U_{l, i} \to U_l\}_{i \in I_l} \in \text{Cov}(\mathcal{C})$
et $x_{l, i} \in \mathcal{M}^{n_l - 1}(U_{l, i})$ tels
que $\text{d}(x_{l, i}) = x|_{U_{l, i}}$. Posons alors
$$
K = \{(U_{l, i}, n_l - 1, x_{l, i}) \mid l \in L, i \in I_l\}
$$
et prenons pour $\mathcal{F}_{n + 1}$ le plus petit sous-module
différentiel gradué sur $\mathcal{A}$ de $\mathcal{M}$ contenant
$\mathcal{F}_n$ et les sections $x_{l, i}$.
Puisque $|K| \leq \max(\kappa, |\mathcal{F}_n|)$
nous concluons que $|\mathcal{F}_{n + 1}| \leq \kappa$ par induction.

\medskip\noindent
Par construction, l'inclusion $\mathcal{F}_n \to \mathcal{F}_{n + 1}$ induit
l'application nulle sur les faisceaux de cohomologie. Ainsi,
$\mathcal{F} = \bigcup \mathcal{F}_n$ est un sous-module acyclique non nul tel
que $|\mathcal{F}| \leq \kappa$. Puisque les classes d'isomorphisme
de $\mathcal{A}$-modules différentiels gradués
$\mathcal{F}$ dont $|\mathcal{F}|$ est borné forment un ensemble, nous concluons.
\end{proof}

\begin{definition}
\label{definition-K-injective} Soit
$(\mathcal{C}, \mathcal{O})$ un site annelé. Soit
$(\mathcal{A}, \text{d})$ un faisceau d'algèbres différentielles graduées sur
$(\mathcal{C}, \mathcal{O})$. Un $\mathcal{A}$-module différentiel gradué
$\mathcal{I}$ est dit {\it K-injectif} si, pour tout module différentiel
gradué acyclique $\mathcal{M}$, nous
avons $$
\Hom_{K(\textit{Mod}(\mathcal{A}, \text{d}))}(\mathcal{M}, \mathcal{I}) = 0
$$
\end{definition}

\noindent
Notons la similitude avec
Catégories dérivées, définition \ref{derived-definition-K-injective}.

\begin{lemma}
\label{lemma-product-K-injective}
Soit $(\mathcal{C}, \mathcal{O})$ un site annelé. Soit
$(\mathcal{A}, \text{d})$ un faisceau d'algèbres différentielles graduées sur
$(\mathcal{C}, \mathcal{O})$. Soit $T$ un ensemble et, pour
chaque $t \in T$, soit $\mathcal{I}_t$ un module différentiel
gradué K-injectif sur $\mathcal{A}$. Alors
$\prod \mathcal{I}_t$ est un
$\mathcal{A}$-module différentiel gradué K-injectif.
\end{lemma}

\begin{proof}
Soit $\mathcal{K}$ un module différentiel gradué acyclique sur $\mathcal{A}$. Alors
nous avons
$$
\Hom_{\textit{Mod}^{dg}(\mathcal{A}, \text{d})}(\mathcal{K},
\prod\nolimits_{t \in T} \mathcal{I}_t)
=
\prod\nolimits_{t \in T}
\Hom_{\textit{Mod}^{dg}(\mathcal{A}, \text{d})}(\mathcal{K}, \mathcal{I}_t)
$$ car la formation
des produits dans $\textit{Mod}(\mathcal{A}, \text{d})$ commute au
foncteur d'oubli vers les modules gradués sur $\mathcal{A}$. Puisque la
prise de produits est un foncteur exact sur la catégorie des groupes abéliens, nous
en concluons.
\end{proof}

\begin{lemma}
\label{lemma-first-property-dg-injective} Soit
$(\mathcal{C}, \mathcal{O})$ un site annelé. Soit
$(\mathcal{A}, \text{d})$ un
faisceau d'algèbres différentielles graduées sur $(\mathcal{C}, \mathcal{O})$. Soit
$\mathcal{I}$ un objet K-injectif et injectif gradué de
$\textit{Mod}(\mathcal{A}, \text{d})$. Pour chaque diagramme
en traits pleins dans $\textit{Mod}(\mathcal{A}, \text{d})$
$$
\xymatrix{
\mathcal{M} \ar[r]_a \ar[d]_b & \mathcal{I} \\
\mathcal{M}' \ar@{..>}[ru]
}
$$ où
$b$ est injectif et $\mathcal{M}$ acyclique, il existe une
flèche pointillée rendant le diagramme commutatif.
\end{lemma}

\begin{proof}
Puisque $\mathcal{M}$ est acyclique et $\mathcal{I}$ est K-injectif, il
existe un homomorphisme de $\mathcal{A}$-modules gradués
$h : \mathcal{M} \to \mathcal{I}$ de degré $-1$ tel
que $a = \text{d}(h)$. Puisque $\mathcal{I}$ est injectif gradué et
$b$ est injectif, il existe un homomorphisme de $\mathcal{A}$-modules gradués
$h' : \mathcal{M}' \to \mathcal{I}$ de degré $-1$ tel que
$h = h' \circ b$. Ensuite, nous pouvons prendre $a' = \text{d}(h')$ comme
la flèche pointillée.
\end{proof}

\begin{lemma}
\label{lemma-second-property-dg-injective} Soit
$(\mathcal{C}, \mathcal{O})$ un site annelé. Soit
$(\mathcal{A}, \text{d})$ un faisceau d'algèbres différentielles graduées sur
$(\mathcal{C}, \mathcal{O})$. Soit $\mathcal{I}$ un objet K-injectif
et injectif gradué de
$\textit{Mod}(\mathcal{A}, \text{d})$. Pour chaque diagramme
en traits pleins dans $\textit{Mod}(\mathcal{A}, \text{d})$
$$
\xymatrix{
\mathcal{M} \ar[r]_a \ar[d]_b & \mathcal{I} \\
\mathcal{M}' \ar@{..>}[ru]
}
$$ où
$b$ est un quasi-isomorphisme, il existe une flèche pointillée rendant le diagramme
commutatif à homotopie près.
\end{lemma}

\begin{proof}
Après avoir remplacé $\mathcal{M}'$ par la somme directe de $\mathcal{M}'$
et du cône de l'identité de $\mathcal{M}$ (qui est acyclique), nous
pouvons supposer que $b$ est également injectif. Alors le conoyau $\mathcal{Q}$
de $b$ est acyclique. Ainsi,
nous voyons que $$
\Hom_{K(\textit{Mod}(\mathcal{A}, \text{d}))}(\mathcal{Q}, \mathcal{I})
= \Hom_{K(\textit{Mod}(\mathcal{A}, \text{d}))}(\mathcal{Q}, \mathcal{I})[1]
= 0 $$
comme $\mathcal{I}$ est K-injectif. Puisque $\mathcal{I}$ est injectif
gradué, la remarque \ref{remark-why-graded-injective}
montre
que $$
\Hom_{K(\textit{Mod}(\mathcal{A}, \text{d}))}(\mathcal{M}', \mathcal{I})
\longrightarrow
\Hom_{K(\textit{Mod}(\mathcal{A}, \text{d}))}(\mathcal{M}, \mathcal{I})
$$
est bijectif et la preuve est complète.
\end{proof}

\begin{lemma}
\label{lemma-better-set-of-monos} Soit
$(\mathcal{C}, \mathcal{O})$ un site annelé. Soit
$(\mathcal{A}, \text{d})$ un faisceau
d'algèbres différentielles graduées sur $(\mathcal{C}, \mathcal{O})$. Il existe un
ensemble $R$ et, pour chaque $r \in R$, une application injective
$\mathcal{M}_r \to \mathcal{M}'_r$ de
$\mathcal{A}$-modules différentiels gradués acycliques, tels que,
pour un objet $\mathcal{I}$ de $\textit{Mod}(\mathcal{A}, \text{d})$, les propriétés
suivantes soient équivalentes :
\begin{enumerate}
\item $\mathcal{I}$ est K-injectif et injectif gradué, et
\item pour chaque diagramme en traits pleins
$$
\xymatrix{
\mathcal{M}_r \ar[r] \ar[d] & \mathcal{I} \\
\mathcal{M}'_r \ar@{..>}[ru]
}
$$ il
existe une flèche pointillée dans $\textit{Mod}(\mathcal{A}, \text{d})$
rendant le diagramme commutatif.
\end{enumerate}
\end{lemma}

\begin{proof}
Soient $T$ et $\mathcal{M}_t \to \mathcal{M}'_t$ comme dans le
lemme \ref{lemma-characterize-graded-injectives-in-dg}. Soient
$S$ et $\mathcal{M}_s$ comme dans le lemme
\ref{lemma-small-acyclics}. Choisissons
une application injective $\mathcal{M}_s \to \mathcal{M}'_s$
de modules différentiels gradués acycliques sur $\mathcal{A}$
qui soit homotope à zéro. Cela est possible parce que nous
pouvons considérer que $\mathcal{M}'_s$ est le cône de l'identité ;
dans ce cas, il est même vrai que l'identité sur
$\mathcal{M}'_s$ est homotope à zéro, voir Algèbre
différentielle graduée, lemme \ref{dga-lemma-id-cone-null}, qui s'applique
par la discussion de la section \ref{section-conclude-triangulated}. Nous affirmons
que $R = T \coprod S$ avec les applications données fonctionne.

\medskip\noindent
L'implication (1) $\Rightarrow$ (2) résulte du
lemme \ref{lemma-first-property-dg-injective}.

\medskip\noindent
Supposons (2). Tout d'abord, par le lemme \ref{lemma-characterize-graded-injectives-in-dg}
nous voyons que $\mathcal{I}$ est injectif gradué.
Soit maintenant $\mathcal{M}$ un module différentiel gradué acyclique sur $\mathcal{A}$.
Nous
devons montrer que $$
\Hom_{K(\textit{Mod}(\mathcal{A}, \text{d}))}(\mathcal{M}, \mathcal{I})
= 0 $$
La démonstration sera identique à celle
d'Injectifs, lemme \ref{injectives-lemma-characterize-K-injective}.

\medskip\noindent
Nous allons construire par récurrence transfinie sur l'ordinal $\alpha$
un sous-module différentiel gradué acyclique
$\mathcal{K}_\alpha \subset \mathcal{M}$ comme suit.
Pour $\alpha = 0$, posons $\mathcal{K}_0 = 0$. Pour $\alpha > 0$,
on procède comme suit :
\begin{enumerate}
\item Si $\alpha = \beta + 1$ et $\mathcal{K}_\beta = \mathcal{M}$,
alors nous choisissons $\mathcal{K}_\alpha = \mathcal{K}_\beta$.
\item Si $\alpha = \beta + 1$ et $\mathcal{K}_\beta \not = \mathcal{M}$,
alors $\mathcal{M}/\mathcal{K}_\beta$ est un module différentiel gradué
acyclique non nul sur $\mathcal{A}$.
Nous choisissons un sous-$\mathcal{A}$-module différentiel gradué
$\mathcal{N}_\alpha \subset \mathcal{M}/\mathcal{K}_\beta$
isomorphe à $\mathcal{M}_s$ pour un certain $s \in S$ ; voir le
lemme \ref{lemma-small-acyclics}.
Enfin, nous définissons $\mathcal{K}_\alpha \subset \mathcal{M}$
comme l'image inverse de $\mathcal{N}_\alpha$.
\item Si $\alpha$ est un ordinal limite, posons
$\mathcal{K}_\beta = \colim \mathcal{K}_\alpha$.
\end{enumerate}
Il est clair que $\mathcal{M} = \mathcal{K}_\alpha$ pour un
ordinal $\alpha$ suffisamment grand. Nous allons montrer que
$$
\Hom_{K(\textit{Mod}(\mathcal{A}, \text{d}))}(\mathcal{K}_\alpha, \mathcal{I})
$$
est nul par récurrence transfinie sur $\alpha$. C'est vrai pour $\alpha = 0$,
car $\mathcal{K}_0$ est nul. Supposons-le vrai pour $\beta$ et
$\alpha = \beta + 1$. Dans le cas (1) de la liste ci-dessus, le résultat est clair.
Dans le cas (2), il existe une suite exacte courte
$$
0 \to \mathcal{K}_\beta \to \mathcal{K}_\alpha \to \mathcal{N}_\alpha \to 0
$$
D'après la remarque \ref{remark-why-graded-injective}
et puisque nous avons vu que $\mathcal{I}$ est injectif
gradué, nous obtenons une suite exacte
$$
\Hom_{K(\textit{Mod}(\mathcal{A}, \text{d}))}(\mathcal{K}_\beta, \mathcal{I})
\to
\Hom_{K(\textit{Mod}(\mathcal{A}, \text{d}))}(\mathcal{K}_\alpha, \mathcal{I})
\to
\Hom_{K(\textit{Mod}(\mathcal{A}, \text{d}))}(\mathcal{N}_\alpha, \mathcal{I})
$$
Par récurrence, le terme de gauche est nul. Par l'hypothèse (2),
celui de droite est nul : toute application $\mathcal{M}_s \to \mathcal{I}$
se factorise par $\mathcal{M}'_s$ et est donc homotope à zéro.
Le groupe du milieu est donc nul lui aussi. Enfin, supposons que $\alpha$ soit un
ordinal limite. Puisque nous avons aussi
$\mathcal{K}_\alpha = \colim \mathcal{K}_\alpha$ comme
$\mathcal{A}$-modules gradués, nous voyons que
$$
\Hom_{\textit{Mod}^{dg}(\mathcal{A}, \text{d})}
(\mathcal{K}_\alpha, \mathcal{I})
= \lim_{\beta < \alpha}
\Hom_{\textit{Mod}^{dg}(\mathcal{A}, \text{d})}
(\mathcal{K}_\beta, \mathcal{I})
$$
comme complexes de groupes abéliens. Les groupes de cohomologie de ces
complexes calculent les morphismes dans $K(\textit{Mod}(\mathcal{A}, \text{d}))$
entre décalages. Les morphismes de transition du système de complexes
sont surjectifs d'après la remarque \ref{remark-why-graded-injective},
car $\mathcal{I}$ est injectif gradué.
De plus, pour tout ordinal limite $\beta \leq \alpha$,
la limite coïncide avec la valeur. Nous pouvons donc appliquer le
lemme \ref{homology-lemma-ML-over-ordinals} d'Homologie
pour conclure.
\end{proof}

\begin{lemma}
\label{lemma-functor-set-of-monos} Soit
$(\mathcal{C}, \mathcal{O})$ un site annelé. Soit
$(\mathcal{A}, \text{d})$ un faisceau
d'algèbres différentielles graduées sur $(\mathcal{C}, \mathcal{O})$. Soit
$R$ un ensemble et, pour chaque $r \in R$, soit donnée une application
injective $\mathcal{M}_r \to \mathcal{M}'_r$
de $\mathcal{A}$-modules différentiels gradués acycliques.
Il existe un foncteur $M : \textit{Mod}(\mathcal{A}, \text{d}) \to
\textit{Mod}(\mathcal{A}, \text{d})$ et une transformation naturelle
$j : \text{id} \to M$ telle que
\begin{enumerate}
\item $j_\mathcal{M} : \mathcal{M} \to M(\mathcal{M})$ est injectif et est un
quasi-isomorphisme,
\item pour tout diagramme en traits pleins
$$
\xymatrix{
\mathcal{M}_r \ar[r] \ar[d] & \mathcal{M} \ar[d]^{j_\mathcal{M}} \\
\mathcal{M}'_r \ar@{..>}[r] & M(\mathcal{M})
}
$$ il existe
une flèche pointillée dans $\textit{Mod}(\mathcal{A}, \text{d})$
rendant le diagramme commutatif.
\end{enumerate}
\end{lemma}

\begin{proof}
Définissons $M(\mathcal{M})$ comme la somme amalgamée du diagramme suivant
$$
\xymatrix{
\bigoplus_{(r, \varphi)} \mathcal{M}_r \ar[r] \ar[d] &
\mathcal{M} \ar[d] \\
\bigoplus_{(r, \varphi)} \mathcal{M}'_r \ar[r] &
M(\mathcal{M})
}
$$ où la somme
directe porte sur toutes les paires $(r, \varphi)$ avec
$r \in R$ et $\varphi \in
\Hom_{\textit{Mod}(\mathcal{A}, \text{d})}(\mathcal{M}_r, \mathcal{M})$. Puisque tout changement de cobase
d'un morphisme injectif est injectif, nous voyons que
$\mathcal{M} \to M(\mathcal{M})$ est injectif. Puisque le
conoyau de la flèche verticale gauche est acyclique, nous voyons
que le conoyau, qui lui est isomorphe, de $\mathcal{M} \to M(\mathcal{M})$ est
acyclique ; ainsi $\mathcal{M} \to M(\mathcal{M})$ est un
quasi-isomorphisme. La propriété (2) est vérifiée par construction.
Nous omettons la vérification que
cette procédure peut être transformée en foncteur.
\end{proof}

\begin{theorem}
\label{theorem-qis-into-dg-injective}
Soit $(\mathcal{C}, \mathcal{O})$ un site annelé. Soit
$(\mathcal{A}, \text{d})$ un
faisceau d'algèbres différentielles graduées sur $(\mathcal{C}, \mathcal{O})$.
Pour tout $\mathcal{A}$-module différentiel gradué $\mathcal{M}$, il
existe un quasi-isomorphisme $\mathcal{M} \to \mathcal{I}$
où $\mathcal{I}$ est un module différentiel gradué injectif
gradué et K-injectif sur $\mathcal{A}$. En outre, la
construction est fonctorielle en $\mathcal{M}$.
\end{theorem}

\begin{proof}
Soit $R$ l'ensemble indexant les morphismes $\mathcal{M}_r \to \mathcal{M}'_r$ de
$\textit{Mod}(\mathcal{A}, \text{d})$
considérés dans le lemme \ref{lemma-better-set-of-monos}. Soit
$M$, muni de la transformation $\text{id} \to M$, le
foncteur construit au lemme \ref{lemma-functor-set-of-monos} à
partir de $R$ et des $\mathcal{M}_r \to \mathcal{M}'_r$.
Par récurrence transfinie, nous définissons une suite de foncteurs
$M_\alpha$ et des transformations naturelles $M_\beta \to M_\alpha$
pour $\alpha < \beta$, en posant
\begin{enumerate}
\item $M_0 = \text{id}$,
\item $M_{\alpha + 1} = M \circ M_\alpha$, avec la transformation naturelle
$M_\beta \to M_{\alpha + 1}$ pour $\beta < \alpha + 1$
provenant des transformations déjà construites $M_\beta \to M_\alpha$ et des
morphismes $M_\alpha \to M \circ M_\alpha$ induits par $\text{id} \to M$, et
\item $M_\alpha = \colim_{\beta < \alpha} M_\beta$ si $\alpha$ est un
ordinal limite, avec les coprojections comme transformations
$M_\beta \to M_\alpha$ pour $\alpha < \beta$.
\end{enumerate} Notons que
pour chaque $\mathcal{A}$-module différentiel gradué les applications
$\mathcal{M} \to M_\beta(\mathcal{M}) \to M_\alpha(\mathcal{M})$ sont des
quasi-isomorphismes injectifs (comme les limites inductives filtrantes sont exactes).

\medskip\noindent
Rappelons que $\textit{Mod}(\mathcal{A}, \text{d})$ est une catégorie
abélienne de Grothendieck. Ainsi,
par Injectifs, proposition \ref{injectives-proposition-objects-are-small}
(appliquée à la somme directe des $\mathcal{M}_r$ pour tous les $r \in R$),
il existe un ordinal limite $\alpha$ tel que $\mathcal{M}_r$ soit $\alpha$-petit
par rapport aux injections pour tout $r \in R$.
Nous affirmons que $\mathcal{M} \to M_\alpha(\mathcal{M})$
est le plongement fonctoriel souhaité de $\mathcal{M}$ dans un
module injectif gradué et K-injectif.

\medskip\noindent
En effet, tout morphisme $\mathcal{M}_r \to M_\alpha(\mathcal{M})$
se factorise par $M_\beta(\mathcal{M})$ pour un certain $\beta < \alpha$.
Par la construction de $M$, cela signifie que
$\mathcal{M}_r \to M_{\beta + 1}(\mathcal{M}) = M(M_\beta(\mathcal{M}))$ se factorise
par $\mathcal{M}'_r$. Puisque
$M_\beta(\mathcal{M}) \subset M_{\beta + 1}(\mathcal{M})
\subset M_\alpha(\mathcal{M})$, on obtient la factorisation souhaitée dans
$M_\alpha(\mathcal{M})$. Nous concluons par notre choix de
$R$ et $\mathcal{M}_r \to \mathcal{M}'_r$ dans le
lemme \ref{lemma-better-set-of-monos}.
\end{proof}










\section{La catégorie dérivée}
\label{section-derived}

\noindent
Cette section est l'analogue de l'Algèbre différentielle graduée, section
\ref{dga-section-derived}.

\medskip\noindent
Soit $(\mathcal{C}, \mathcal{O})$ un site annelé. Soit
$(\mathcal{A}, \text{d})$ un faisceau d'algèbres différentielles graduées
sur $(\mathcal{C}, \mathcal{O})$. Nous construirons la catégorie
dérivée $D(\mathcal{A}, \text{d})$ en inversant les quasi-isomorphismes
dans $K(\textit{Mod}(\mathcal{A}, \text{d}))$.

\begin{lemma}
\label{lemma-cohomology-homological}
Soit $(\mathcal{C}, \mathcal{O})$ un site annelé. Soit
$(\mathcal{A}, \text{d})$ un faisceau d'algèbres différentielles graduées
sur $(\mathcal{C}, \mathcal{O})$. Le foncteur
$H^0 : \textit{Mod}(\mathcal{A}, \text{d}) \to \textit{Mod}(\mathcal{O})$
de la section \ref{section-modules} se factorise par
un
foncteur $$
H^0 : K(\textit{Mod}(\mathcal{A}, \text{d})) \to \textit{Mod}(\mathcal{O})
$$
qui est homologique au sens des
Catégories dérivées, définition \ref{derived-definition-homological}.
\end{lemma}

\begin{proof}
Il découle immédiatement des définitions qu'il existe un
diagramme commutatif
$$
\xymatrix{
\textit{Mod}(\mathcal{A}, \text{d}) \ar[r] \ar[d] &
K(\textit{Mod}(\mathcal{A}, \text{d})) \ar[d] \\
\text{Comp}(\mathcal{O}) \ar[r] &
K(\textit{Mod}(\mathcal{O}))
}
$$
Puisque $H^0(\mathcal{M})$ est le faisceau de cohomologie de
degré zéro du complexe de $\mathcal{O}$-modules sous-jacent à $\mathcal{M}$,
le lemme découle du cas des complexes de $\mathcal{O}$-modules, qui
est un cas particulier de
Catégories dérivées, lemme \ref{derived-lemma-cohomology-homological}.
\end{proof}

\begin{lemma}
\label{lemma-acyclics}
Soit $(\mathcal{C}, \mathcal{O})$ un site annelé. Soit
$(\mathcal{A}, \text{d})$ un faisceau d'algèbres différentielles graduées
sur $(\mathcal{C}, \mathcal{O})$. La sous-catégorie pleine $\text{Ac}$
de la catégorie homotopique $K(\textit{Mod}(\mathcal{A}, \text{d}))$
composée de modules acycliques est une sous-catégorie triangulée
strictement pleine et saturée de $K(\textit{Mod}(\mathcal{A}, \text{d}))$.
\end{lemma}

\begin{proof}
Bien sûr, un objet $\mathcal{M}$ de $K(\textit{Mod}(\mathcal{A}, \text{d}))$ est
dans $\text{Ac}$ si et seulement si $H^i(\mathcal{M}) = H^0(\mathcal{M}[i])$ est
nul pour tout $i$. Le lemme résulte de cette
observation, du lemme \ref{lemma-cohomology-homological}
et de Catégories dérivées, lemme \ref{derived-lemma-homological-functor-kernel}.
Voir aussi Catégories dérivées, définitions \ref{derived-definition-saturated}
et \ref{derived-definition-triangulated-subcategory}
et lemme \ref{derived-lemma-triangulated-subcategory}.
\end{proof}

\begin{lemma}
\label{lemma-qis}
Soit $(\mathcal{C}, \mathcal{O})$ un site annelé. Soit
$(\mathcal{A}, \text{d})$ un faisceau d'algèbres différentielles graduées
sur $(\mathcal{C}, \mathcal{O})$.
Considérons la
sous-classe $\text{Qis} \subset \text{Arrows}(K(\textit{Mod}(\mathcal{A}, \text{d})))$
composée de quasi-isomorphismes. Il s'agit d'un système multiplicatif
saturé compatible avec la structure triangulée sur
$K(\textit{Mod}(\mathcal{A}, \text{d}))$.
\end{lemma}

\begin{proof}
Observons que si $f , g : \mathcal{M} \to \mathcal{N}$ sont des morphismes
de $\textit{Mod}(\mathcal{A}, \text{d})$ qui sont homotopes,
alors $f$ est un quasi-isomorphisme si et seulement si $g$ est un quasi-isomorphisme.
À savoir, les applications $H^i(f) = H^0(f[i])$ et $H^i(g) = H^0(g[i])$
sont les mêmes par le lemme \ref{lemma-cohomology-homological}.
Il est donc sans ambiguïté de dire qu'un morphisme de la
catégorie homotopique $K(\textit{Mod}(\mathcal{A}, \text{d}))$ est un
quasi-isomorphisme. Pour les définitions de « système multiplicatif »,
« saturé » et « compatible avec la structure triangulée
», voir Catégories dérivées, définition \ref{derived-definition-localization},
et
Catégories, définitions \ref{categories-definition-multiplicative-system}
et \ref{categories-definition-saturated-multiplicative-system}.

\medskip\noindent
Pour démontrer le lemme, considérons la composition de
foncteurs exacts de catégories triangulées
$$
K(\textit{Mod}(\mathcal{A}, \text{d}))
\longrightarrow
K(\textit{Mod}(\mathcal{O}))
\longrightarrow
D(\mathcal{O})
$$ et
remarquons qu'un morphisme $f : \mathcal{M} \to \mathcal{N}$ de
$K(\textit{Mod}(\mathcal{A}, \text{d}))$ appartient à $\text{Qis}$ si
et seulement si son image est un isomorphisme dans $D(\mathcal{O})$.
Ainsi, le lemme découle de Catégories dérivées, lemme
\ref{derived-lemma-triangle-functor-localize}.
\end{proof}

\noindent
Dans le cas du lemme \ref{lemma-qis}, nous pouvons appliquer
Catégories dérivées, proposition
\ref{derived-proposition-construct-localization} pour
obtenir un foncteur exact de catégories triangulées
$$ Q
:
K(\textit{Mod}(\mathcal{A}, \text{d}))
\longrightarrow
\text{Qis}^{-1}K(\textit{Mod}(\mathcal{A}, \text{d}))
$$ Cependant,
comme $\textit{Mod}(\mathcal{A}, \text{d})$ est une « grande »
catégorie, c'est-à-dire que ses objets forment une classe propre, il n'est pas
immédiatement clair que, pour $\mathcal{M}$ et $\mathcal{N}$ donnés, la
construction de $\text{Qis}^{-1}K(\textit{Mod}(\mathcal{A}, \text{d}))$
produise un {\it ensemble}
$$
\Mor_{\text{Qis}^{-1}K(\textit{Mod}(\mathcal{A}, \text{d}))}
(\mathcal{M}, \mathcal{N})
$$
de morphismes. La construction de complexes K-injectifs permet
toutefois de le démontrer. En effet, par le théorème \ref{theorem-qis-into-dg-injective},
choisissons un quasi-isomorphisme $s_0 : \mathcal{N} \to \mathcal{I}$ où
$\mathcal{I}$ est un module différentiel gradué injectif
gradué et K-injectif sur $\mathcal{A}$. Rappelons ensuite que les éléments
de l'ensemble affiché sont des classes d'équivalence des
paires $(f : \mathcal{M} \to \mathcal{N}', s : \mathcal{N} \to \mathcal{N}')$
où $f$ est un morphisme arbitraire de
$K(\textit{Mod}(\mathcal{A}, \text{d}))$ et $s$ est un quasi-isomorphisme, voir la
description du calcul gauche des fractions dans Catégories,
section \ref{categories-section-localization}. Par le lemme
\ref{lemma-second-property-dg-injective} nous pouvons
choisir la flèche pointillée
$$
\xymatrix{
\mathcal{M} \ar[rd]^f & &
\mathcal{N} \ar[ld]_s \ar[rd]^{s_0} \\ &
\mathcal{N}' \ar@{..>}[rr]^{s'} & & \mathcal{I}
}
$$ rendant
le diagramme commutatif (dans la catégorie homotopique).
Ainsi, la paire $(f, s)$ est équivalente à la paire $(s' \circ f, s_0)$ et
nous trouvons que la collection de classes d'équivalence forme un ensemble.

\begin{definition}
\label{definition-derived-category} Soit
$(\mathcal{C}, \mathcal{O})$ un site annelé. Soit
$(\mathcal{A}, \text{d})$ un faisceau d'algèbres différentielles graduées sur
$(\mathcal{C}, \mathcal{O})$. Soit $\text{Qis}$ comme dans le lemme
\ref{lemma-qis}. La
{\it catégorie dérivée de $(\mathcal{A}, \text{d})$} est la catégorie
triangulée
$$
D(\mathcal{A}, \text{d}) =
\text{Qis}^{-1}K(\textit{Mod}(\mathcal{A}, \text{d}))
$$ dont la construction a été
discutée plus en détail ci-dessus.
\end{definition}

\noindent
Nous prouvons quelques faits sur cette construction.

\begin{lemma}
\label{lemma-kernel-localization}
Dans la définition \ref{definition-derived-category},
le noyau du foncteur de localisation
$Q : K(\textit{Mod}(\mathcal{A}, \text{d})) \to D(\mathcal{A}, \text{d})$
est la catégorie $\text{Ac}$ du lemme \ref{lemma-acyclics}.
\end{lemma}

\begin{proof}
Cela découle immédiatement
de Catégories dérivées, lemme \ref{derived-lemma-kernel-localization},
et du fait que $0 \to \mathcal{M}$ est un quasi-isomorphisme si
et seulement si $\mathcal{M}$ est acyclique.
\end{proof}

\begin{lemma}
\label{lemma-H0-over-D}
Dans la définition \ref{definition-derived-category}, le foncteur
$H^0 : K(\textit{Mod}(\mathcal{A}, \text{d})) \to
\textit{Mod}(\mathcal{O})$ se factorise par un foncteur homologique
$H^0 : D(\mathcal{A}, \text{d}) \to \textit{Mod}(\mathcal{O})$.
\end{lemma}

\begin{proof}
Cela découle immédiatement
de Catégories dérivées, lemme \ref{derived-lemma-universal-property-localization}.
\end{proof}

\noindent
Voici le lemme annoncé qui calcule les ensembles de morphismes dans la
catégorie dérivée.

\begin{lemma}
\label{lemma-hom-derived}
Soit $(\mathcal{C}, \mathcal{O})$ un site annelé. Soit
$(\mathcal{A}, \text{d})$ un faisceau d'algèbres différentielles graduées
sur $(\mathcal{C}, \mathcal{O})$.
Soient $\mathcal{M}$ et $\mathcal{N}$ des
$\mathcal{A}$-modules différentiels gradués. Soit $\mathcal{N} \to \mathcal{I}$
un quasi-isomorphisme, où $\mathcal{I}$ est un
$\mathcal{A}$-module différentiel gradué injectif gradué et K-injectif. Alors
$$
\Hom_{D(\mathcal{A}, \text{d})}(\mathcal{M}, \mathcal{N}) =
\Hom_{K(\textit{Mod}(\mathcal{A}, \text{d}))}(\mathcal{M}, \mathcal{I})
$$
\end{lemma}

\begin{proof}
Puisque $\mathcal{N} \to \mathcal{I}$ est un quasi-isomorphisme
nous voyons
que $$
\Hom_{D(\mathcal{A}, \text{d})}(\mathcal{M}, \mathcal{N})
= \Hom_{D(\mathcal{A}, \text{d})}(\mathcal{M}, \mathcal{I})
$$
Dans la discussion précédant la définition \ref{definition-derived-category},
nous avons constaté, en utilisant le lemme \ref{lemma-second-property-dg-injective},
que tout morphisme $\mathcal{M} \to \mathcal{I}$
dans $D(\mathcal{A}, \text{d})$ peut être représenté
par un morphisme $f : \mathcal{M} \to \mathcal{I}$
dans $K(\textit{Mod}(\mathcal{A}, \text{d}))$.
Si $f, f' :  \mathcal{M} \to \mathcal{I}$ sont deux
morphismes de $K(\textit{Mod}(\mathcal{A}, \text{d}))$, ils
définissent le même morphisme dans $D(\mathcal{A}, \text{d})$ si et
seulement s'il existe un quasi-isomorphisme $g : \mathcal{I} \to \mathcal{K}$
dans $K(\textit{Mod}(\mathcal{A}, \text{d}))$
tel que $g \circ f = g \circ f'$ ; voir
Catégories, lemme \ref{categories-lemma-equality-morphisms-left-localization}.
Cependant, par le lemme \ref{lemma-second-property-dg-injective}, il
existe une
application $h : \mathcal{K} \to \mathcal{I}$
telle que $h \circ g = \text{id}_\mathcal{I}$
dans $K(\textit{Mod}(\mathcal{A}, \text{d}))$.
Ainsi, $g \circ f = g \circ f'$ signifie $f = f'$ et
la preuve est complète.
\end{proof}

\begin{lemma}
\label{lemma-derived-products} Soit
$(\mathcal{C}, \mathcal{O})$ un site annelé. Soit
$(\mathcal{A}, \text{d})$ un faisceau d'algèbres différentielles graduées sur
$(\mathcal{C}, \mathcal{O})$. Alors
\begin{enumerate}
\item $D(\mathcal{A}, \text{d})$ possède des sommes directes et des produits,
\item les sommes directes se calculent en prenant les sommes directes de
$\mathcal{A}$-modules différentiels gradués,
\item les produits se calculent en prenant les produits
de modules différentiels gradués K-injectifs.
\end{enumerate}
\end{lemma}

\begin{proof}
Nous utiliserons que $\textit{Mod}(\mathcal{A}, \text{d})$ est une catégorie abélienne
avec des sommes directes et des produits arbitraires, et que ceux-ci donnent
lieu à des sommes directes et des produits dans $K(\textit{Mod}(\mathcal{A}, \text{d}))$.
Voir lemmes \ref{lemma-dgm-abelian} et \ref{lemma-homotopy-direct-sums}.

\medskip\noindent
Soit $\mathcal{M}_j$ une famille de $\mathcal{A}$-modules différentiels gradués.
Considérons la somme directe $\mathcal{M} = \bigoplus \mathcal{M}_j$
comme un $\mathcal{A}$-module différentiel gradué.
Pour un $\mathcal{A}$-module différentiel gradué
$\mathcal{N}$, choisissons un quasi-isomorphisme
$\mathcal{N} \to \mathcal{I}$
où $\mathcal{I}$ est injectif gradué et K-injectif comme
$\mathcal{A}$-module différentiel gradué. Voir le
théorème \ref{theorem-qis-into-dg-injective}.
En utilisant le lemme \ref{lemma-hom-derived},
nous avons \begin{align*}
\Hom_{D(\mathcal{A}, \text{d})}(\mathcal{M}, \mathcal{N})
&
= \Hom_{K(\mathcal{A}, \text{d})}(\mathcal{M}, \mathcal{I}) \\
&
= \prod \Hom_{K(\mathcal{A}, \text{d})}(\mathcal{M}_j, \mathcal{I}) \\
&
= \prod \Hom_{D(\mathcal{A}, \text{d})}(\mathcal{M}_j, \mathcal{I})
\end{align*}
ce qui prouve l'existence des sommes directes dans $D(A, \text{d})$,
comme indiqué en (2).

\medskip\noindent
Soit $\mathcal{M}_j$ une famille de $\mathcal{A}$-modules différentiels gradués.
Pour chaque $j$, choisissons un
quasi-isomorphisme $\mathcal{M} \to \mathcal{I}_j$
où $\mathcal{I}_j$ est injectif gradué et K-injectif
comme $\mathcal{A}$-module différentiel
gradué. Considérons le produit $\mathcal{I} = \prod \mathcal{I}_j$
de modules différentiels gradués sur $\mathcal{A}$.
Par les lemmes \ref{lemma-product-K-injective}
et \ref{lemma-product-graded-injective},
$\mathcal{I}$ est injectif gradué et K-injectif
comme $\mathcal{A}$-module différentiel
gradué. Pour un $\mathcal{A}$-module différentiel
gradué $\mathcal{N}$, le lemme \ref{lemma-hom-derived} donne
\begin{align*}
\Hom_{D(\mathcal{A}, \text{d})}(\mathcal{N}, \mathcal{I})
&
= \Hom_{K(\mathcal{A}, \text{d})}(\mathcal{N}, \mathcal{I}) \\
&
= \prod \Hom_{K(\mathcal{A}, \text{d})}(\mathcal{N}, \mathcal{I}_j) \\
&
= \prod \Hom_{D(\mathcal{A}, \text{d})}(\mathcal{N}, \mathcal{M}_j)
\end{align*}
ce qui prouve l'existence des produits dans $D(\mathcal{A}, \text{d})$,
comme indiqué en (3).
\end{proof}







\section{Le foncteur delta canonique}
\label{section-canonical-delta-functor}

\noindent
Soit $(\mathcal{C}, \mathcal{O})$ un site annelé. Soit
$(\mathcal{A}, \text{d})$ un faisceau d'algèbres différentielles graduées
sur $(\mathcal{C}, \mathcal{O})$.
Considérons le foncteur
$\textit{Mod}(\mathcal{A}, \text{d}) \to
K(\textit{Mod}(\mathcal{A}, \text{d}))$.
Ce foncteur n'est {\bf pas} un foncteur $\delta$ en
général. Cependant, il s'avère que le foncteur
$\textit{Mod}(\mathcal{A}, \text{d}) \to D(A, \text{d})$ est
un foncteur $\delta$. Pour voir cela, nous devons définir
les morphismes $\delta$ associés à une suite exacte courte
$$
0 \to \mathcal{K} \xrightarrow{a} \mathcal{L} \xrightarrow{b} \mathcal{M} \to 0
$$
dans la catégorie abélienne $\textit{Mod}(\mathcal{A}, \text{d})$.
Considérons le cône $C(a)$ du morphisme $a$ avec ses morphismes
canoniques $i : \mathcal{L} \to C(a)$ et
$p : C(a) \to \mathcal{K}[1]$, voir définition \ref{definition-cone}. Il existe un
homomorphisme de modules différentiels gradués sur $\mathcal{A}$
$$ q :
C(a) \longrightarrow \mathcal{M}
$$ par
Algèbre différentielle graduée, lemme \ref{dga-lemma-cone}
(que nous pouvons utiliser dans la discussion
de la section \ref{section-conclude-triangulated})
appliqué au
diagramme $$
\xymatrix{
\mathcal{K} \ar[r]_a \ar[d]
& \mathcal{L} \ar[d]^b \\
0 \ar[r]
& \mathcal{M}
}
$$
Le morphisme $q$ est un quasi-isomorphisme, par exemple parce
que c'est le cas pour les complexes sous-jacents de $\mathcal{O}$-modules ;
voir
Catégories dérivées, section \ref{derived-section-canonical-delta-functor}.
D'après Algèbre différentielle graduée, lemme \ref{dga-lemma-cone-homotopy} (applicable
par la discussion de la section
\ref{section-conclude-triangulated}), le
triangle
$$
(\mathcal{K}, \mathcal{L}, C(a), a, i, -p)
$$ est un
triangle distingué en $K(\textit{Mod}(\mathcal{A}, \text{d}))$. Comme le foncteur
de localisation
$K(\textit{Mod}(\mathcal{A}, \text{d})) \to D(\mathcal{A}, \text{d})$ est exact,
nous voyons que $(\mathcal{K}, \mathcal{L}, C(a), a, i, -p)$ est un
triangle distingué
dans $D(\mathcal{A}, \text{d})$. Puisque $q$ est un quasi-isomorphisme,
nous voyons que $q$ est un isomorphisme dans $D(\mathcal{A}, \text{d})$.
Nous déduisons
donc que $$
(\mathcal{K}, \mathcal{L}, \mathcal{M}, a, b, -p \circ q^{-1})
$$
est un triangle distingué de $D(\mathcal{A}, \text{d})$.
Cela conduit au lemme suivant.

\begin{lemma}
\label{lemma-derived-canonical-delta-functor}
Soit $(\mathcal{C}, \mathcal{O})$ un site annelé. Soit
$(\mathcal{A}, \text{d})$ un faisceau d'algèbres différentielles graduées
sur $(\mathcal{C}, \mathcal{O})$. Le foncteur de localisation
$\textit{Mod}(\mathcal{A}, \text{d}) \to D(\mathcal{A}, \text{d})$ a
la structure naturelle d'un foncteur $\delta$, avec
$$
\delta_{\mathcal{K} \to \mathcal{L} \to \mathcal{M}} = - p \circ q^{-1}
$$ avec
$p$ et $q$ comme expliqué ci-dessus.
\end{lemma}

\begin{proof}
Nous avons déjà vu que ce choix conduit à un triangle distingué
chaque fois qu'une suite exacte courte de complexes est donnée. Il
reste à démontrer la fonctorialité de cette construction ;
voir Catégories dérivées, définition \ref{derived-definition-delta-functor}.
Cela découle, avec quelques vérifications, d'Algèbre différentielle graduée, lemme \ref{dga-lemma-cone},
applicable
par la discussion de la section \ref{section-conclude-triangulated}.
Comparer avec
Catégories dérivées, lemme \ref{derived-lemma-derived-canonical-delta-functor}.
\end{proof}

\begin{lemma}
\label{lemma-homotopy-colimit}
Soit $(\mathcal{C}, \mathcal{O})$ un site annelé. Soit
$(\mathcal{A}, \text{d})$ un faisceau d'algèbres différentielles graduées sur
$(\mathcal{C}, \mathcal{O})$. Soit
$\mathcal{M}_n$ un système de modules différentiels gradués sur $\mathcal{A}$.
Ensuite, la limite inductive dérivée $\text{hocolim} \mathcal{M}_n$
dans $D(\mathcal{A}, \text{d})$ est représentée
par le module différentiel gradué $\colim \mathcal{M}_n$.
\end{lemma}

\begin{proof}
Posons $\mathcal{M} = \colim \mathcal{M}_n$.
Nous avons une suite exacte de $\mathcal{A}$-modules différentiels gradués
$$
0 \to \bigoplus \mathcal{M}_n \to \bigoplus \mathcal{M}_n \to \mathcal{M} \to 0
$$
par Catégories dérivées, lemme \ref{derived-lemma-compute-colimit},
appliqué aux complexes sous-jacents de $\mathcal{O}$-modules. Les
sommes directes sont les sommes directes en $D(\mathcal{A}, \text{d})$
par le lemme \ref{lemma-derived-products}.
Ainsi, le résultat résulte de la
définition des limites inductives dérivées dans Catégories
dérivées, définition \ref{derived-definition-derived-colimit}
et du fait qu'une suite exacte courte de
complexes donne un triangle distingué
(lemme \ref{lemma-derived-canonical-delta-functor}).
\end{proof}






\section{Image inverse dérivée}
\label{section-derived-pullback}

\noindent
Soit $(f, f^\sharp) : (\Sh(\mathcal{C}), \mathcal{O}_\mathcal{C})
\to (\Sh(\mathcal{D}), \mathcal{O}_\mathcal{D})$ un morphisme
de topos annelés. Soit $\mathcal{A}$ une
$\mathcal{O}_\mathcal{C}$-algèbre différentielle graduée. Soit $\mathcal{B}$ une
$\mathcal{O}_\mathcal{D}$-algèbre différentielle graduée. Supposons donné un
morphisme
$$
\varphi : f^{-1}\mathcal{B} \to \mathcal{A}
$$ de
$f^{-1}\mathcal{O}_\mathcal{D}$-algèbres différentielles graduées. Par l'adjonction
entre restriction et extension des scalaires, cela revient à se donner
un morphisme $\varphi : f^*\mathcal{B} \to \mathcal{A}$ de
$\mathcal{O}_\mathcal{C}$-algèbres différentielles graduées ou, de manière équivalente,
$\varphi$ peut être considéré comme un morphisme
$$
\varphi : \mathcal{B} \to f_*\mathcal{A}
$$
de $\mathcal{O}_\mathcal{D}$-algèbres différentielles graduées.
Voir la remarque \ref{remark-functoriality-dga}.

\medskip\noindent
Supposons de plus que $\mathcal{A}'$ soit une seconde
$\mathcal{O}_\mathcal{C}$-algèbre différentielle graduée et que $\mathcal{N}$ soit un
$(\mathcal{A}, \mathcal{A}')$-bimodule différentiel gradué. Dans
ce contexte, nous pouvons considérer le foncteur
$$
\textit{Mod}(\mathcal{B}, \text{d})
\longrightarrow
\textit{Mod}(\mathcal{A}', \text{d}),\quad
\mathcal{M} \longmapsto f^*\mathcal{M} \otimes_{\mathcal{A}} \mathcal{N}
$$
Notons que cela s'étend à un foncteur
$$
\textit{Mod}^{dg}(\mathcal{B}, \text{d})
\longrightarrow
\textit{Mod}^{dg}(\mathcal{A}', \text{d}),\quad
\mathcal{M} \longmapsto f^*\mathcal{M} \otimes_{\mathcal{A}} \mathcal{N}
$$ des
catégories différentielles graduées par la discussion des
sections \ref{section-functoriality-dg} et \ref{section-dg-bimodules}. Il
s'ensuit formellement que nous obtenons également un foncteur
exact \begin{equation}
\label{equation-pullback}
K(\textit{Mod}(\mathcal{B}, \text{d}))
\longrightarrow
K(\textit{Mod}(\mathcal{A}', \text{d})),\quad
\mathcal{M} \longmapsto f^*\mathcal{M} \otimes_{\mathcal{A}} \mathcal{N}
\end{equation}
de catégories triangulées.

\begin{lemma}
\label{lemma-derived-tensor-product} Dans
la situation ci-dessus, le foncteur (\ref{equation-pullback}),
composé avec le foncteur de localisation
$K(\textit{Mod}(\mathcal{A}', \text{d})) \to D(\mathcal{A}', \text{d})$, possède
une extension dérivée à
gauche $D(\mathcal{B}, \text{d}) \to D(\mathcal{A}', \text{d})$ dont
la valeur sur un bon $\mathcal{B}$-module différentiel gradué à droite
$\mathcal{P}$ est $f^*\mathcal{P} \otimes_\mathcal{A} \mathcal{N}$.
\end{lemma}

\begin{proof}
Rappelons que pour tout $\mathcal{B}$-module différentiel gradué à
droite $\mathcal{M}$ il existe un quasi-isomorphisme $\mathcal{P} \to \mathcal{M}$
avec $\mathcal{P}$ un bon $\mathcal{B}$-module différentiel gradué. Voir
lemme \ref{lemma-resolve}. Par
conséquent, d'après Catégories dérivées, lemme \ref{derived-lemma-find-existence-computes},
il suffit de montrer que, pour tout
quasi-isomorphisme $\mathcal{P} \to \mathcal{P}'$ de bons
$\mathcal{B}$-modules différentiels gradués, l'application
induite $$
f^*\mathcal{P} \otimes_\mathcal{A} \mathcal{N}
\longrightarrow
f^*\mathcal{P}' \otimes_\mathcal{A} \mathcal{N}
$$
est un quasi-isomorphisme. Le cône $\mathcal{P}''$
sur $\mathcal{P} \to \mathcal{P}'$ est un bon
$\mathcal{A}$-module différentiel gradué par le lemme \ref{lemma-good-admissible-ses}.
Puisque nous avons un triangle
distingué $$
\mathcal{P} \to \mathcal{P}' \to \mathcal{P}'' \to \mathcal{P}[1]
$$
dans $K(\textit{Mod}(\mathcal{B}, \text{d}))$ nous obtenons un
triangle
distingué $$
f^*\mathcal{P} \otimes_\mathcal{A} \mathcal{N} \to
f^*\mathcal{P}' \otimes_\mathcal{A} \mathcal{N} \to
f^*\mathcal{P}'' \otimes_\mathcal{A} \mathcal{N} \to
f^*\mathcal{P}[1] \otimes_\mathcal{A} \mathcal{N}
$$
dans $K(\textit{Mod}(\mathcal{A}', \text{d}))$. Par
le lemme \ref{lemma-acyclic-good}
le module différentiel gradué
$f^*\mathcal{P}'' \otimes_\mathcal{A} \mathcal{N}$
est acyclique et la preuve est complète.
\end{proof}

\begin{definition}
\label{definition-pullback}
Produit tensoriel dérivé et image inverse dérivée.
\begin{enumerate}
\item Soit $(\mathcal{C}, \mathcal{O})$ un site annelé.
Soient $\mathcal{A}$ et $\mathcal{B}$ des $\mathcal{O}$-algèbres différentielles
graduées. Soit $\mathcal{N}$ un
$(\mathcal{A}, \mathcal{B})$-bimodule différentiel
gradué. Le foncteur $D(\mathcal{A}, \text{d}) \to D(\mathcal{B}, \text{d})$
construit dans le lemme \ref{lemma-derived-tensor-product}
est appelé le {\it produit tensoriel dérivé} et noté
$- \otimes_\mathcal{A}^\mathbf{L} \mathcal{N}$.
\item Soit $(f, f^\sharp) : (\Sh(\mathcal{C}), \mathcal{O}_\mathcal{C})
\to (\Sh(\mathcal{D}), \mathcal{O}_\mathcal{D})$ un
morphisme de topos annelés. Soit $\mathcal{A}$ une
$\mathcal{O}_\mathcal{C}$-algèbre différentielle graduée. Soit $\mathcal{B}$ une
$\mathcal{O}_\mathcal{D}$-algèbre différentielle graduée. Soit
$\varphi : \mathcal{B} \to f_*\mathcal{A}$ un homomorphisme de
$\mathcal{O}_\mathcal{D}$-algèbres différentielles graduées. Le
foncteur $D(\mathcal{B}, \text{d}) \to D(\mathcal{A}, \text{d})$
construit dans le lemme \ref{lemma-derived-tensor-product} est
appelé {\it image inverse dérivée}
et noté $Lf^*$.
\end{enumerate}
\end{definition}

\noindent
Avec ce langage en place, nous pouvons exprimer certaines compatibilités évidentes.

\begin{lemma}
\label{lemma-compose-pullback-tensor}
Dans le lemme \ref{lemma-derived-tensor-product}, le foncteur
$D(\mathcal{B}, \text{d}) \to D(\mathcal{A}', \text{d})$ est égal
à $\mathcal{M} \mapsto
Lf^*\mathcal{M} \otimes_\mathcal{A}^\mathbf{L} \mathcal{N}$.
\end{lemma}

\begin{proof}
Cela résulte immédiatement du fait que nous pouvons calculer ces
foncteurs en représentant les objets par de bons modules différentiels
gradués et parce que $f^*\mathcal{P}$ est un bon
$\mathcal{A}$-module différentiel gradué si $\mathcal{P}$ est un bon
$\mathcal{B}$-module différentiel gradué.
\end{proof}

\begin{lemma}
\label{lemma-compose-pullback}
Soient $(f, f^\sharp) : (\Sh(\mathcal{C}), \mathcal{O})
\to (\Sh(\mathcal{C}'), \mathcal{O}')$ et
$(g, g^\sharp) : (\Sh(\mathcal{C}'), \mathcal{O}')
\to (\Sh(\mathcal{C}''), \mathcal{O}'')$ des
morphismes de topos annelés. Soient $\mathcal{A}$, $\mathcal{A}'$ et
$\mathcal{A}''$ des algèbres différentielles graduées respectivement sur $\mathcal{O}$,
$\mathcal{O}'$ et $\mathcal{O}''$. Soient
$\varphi : \mathcal{A}' \to f_*\mathcal{A}$ et
$\varphi' : \mathcal{A}'' \to g_*\mathcal{A}'$ des
homomorphismes d'algèbres différentielles graduées respectivement sur $\mathcal{O}'$
et $\mathcal{O}''$.
Alors nous avons $L(g \circ f)^* = Lf^* \circ Lg^*
: D(\mathcal{A}'', \text{d}) \to D(\mathcal{A}, \text{d})$.
\end{lemma}

\begin{proof}
Cela résulte immédiatement du calcul de ces foncteurs au
moyen de bons modules différentiels gradués, et du fait
que $f^*\mathcal{P}$ est un bon
$\mathcal{A}'$-module différentiel gradué si $\mathcal{P}$ est un bon
$\mathcal{A}$-module différentiel gradué.
\end{proof}

\noindent
Soit $(\mathcal{C}, \mathcal{O})$ un site annelé. Soient
$\mathcal{A}$ et $\mathcal{B}$ des $\mathcal{O}$-algèbres différentielles graduées.
Soit $\mathcal{N} \to \mathcal{N}'$ un homomorphisme de bimodules différentiels
gradués sur $(\mathcal{A}, \mathcal{B})$. On obtient alors des applications
canoniques
$$
t
: \mathcal{M} \otimes_\mathcal{A}^\mathbf{L} \mathcal{N}
\longrightarrow
\mathcal{M} \otimes_\mathcal{A}^\mathbf{L} \mathcal{N}'
$$
fonctorielles en $\mathcal{M}$ dans $D(\mathcal{A}, \text{d})$,
qui définissent une transformation naturelle entre les
foncteurs exacts $D(\mathcal{A}, \text{d}) \to D(\mathcal{B}, \text{d})$
de catégories triangulées. La valeur de $t$ sur un bon
$\mathcal{A}$-module différentiel gradué $\mathcal{P}$ est l'application
évidente $$
\mathcal{P} \otimes_\mathcal{A}^\mathbf{L} \mathcal{N}
= \mathcal{P} \otimes_\mathcal{A} \mathcal{N}
\longrightarrow
\mathcal{P} \otimes_\mathcal{A} \mathcal{N}'
= \mathcal{P} \otimes_\mathcal{A}^\mathbf{L} \mathcal{N}'
$$

\begin{lemma}
\label{lemma-tensor-symmetry} Dans
la situation ci-dessus, si $\mathcal{N} \to \mathcal{N}'$ est un isomorphisme sur
les faisceaux de cohomologie, alors $t$ est un isomorphisme des
foncteurs $(- \otimes_\mathcal{A}^\mathbf{L} \mathcal{N}) \to
(- \otimes_\mathcal{A}^\mathbf{L} \mathcal{N}')$.
\end{lemma}

\begin{proof}
Il suffit de démontrer que
$\mathcal{P} \otimes_\mathcal{A} \mathcal{N} \to
\mathcal{P} \otimes_\mathcal{A} \mathcal{N}'$ est
un isomorphisme sur les faisceaux de cohomologie pour tout bon
$\mathcal{A}$-module différentiel gradué $\mathcal{P}$.
Pour cela, soit $\mathcal{N}''$ le cône du morphisme
$\mathcal{N} \to \mathcal{N}'$ comme
$\mathcal{A}$-module différentiel gradué à gauche, voir la définition \ref{definition-cone}.
(Bien sûr, $\mathcal{N}''$ est aussi un bimodule, mais
nous n'en avons pas besoin.) Par fonctorialité de la construction
tensorielle (c'est un foncteur de catégories
différentielles graduées), $\mathcal{P} \otimes_\mathcal{A} \mathcal{N}''$
est le cône, comme complexe de $\mathcal{O}$-modules, du
morphisme $\mathcal{P} \otimes_\mathcal{A} \mathcal{N} \to
\mathcal{P} \otimes_\mathcal{A} \mathcal{N}'$.
Il suffit donc de démontrer
que $\mathcal{P} \otimes_\mathcal{A} \mathcal{N}''$
est acyclique. Cela découle du fait que $\mathcal{P}$ est bon
et du fait que $\mathcal{N}''$ est acyclique en tant
que cône sur un quasi-isomorphisme.
\end{proof}

\begin{lemma}
\label{lemma-good-on-other-side}
Soit $(\mathcal{C}, \mathcal{O})$ un site annelé. Soient
$\mathcal{A}$ et $\mathcal{B}$ des $\mathcal{O}$-algèbres différentielles
graduées. Soit $\mathcal{N}$ un bimodule différentiel
gradué sur $(\mathcal{A}, \mathcal{B})$. Si $\mathcal{N}$ est
bon comme $\mathcal{A}$-module différentiel gradué à
gauche, alors $\mathcal{M} \otimes_\mathcal{A}^\mathbf{L} \mathcal{N} =
\mathcal{M} \otimes_\mathcal{A} \mathcal{N}$ pour tout
$\mathcal{A}$-module différentiel gradué $\mathcal{M}$.
\end{lemma}

\begin{proof}
Soit $\mathcal{P} \to \mathcal{M}$ un quasi-isomorphisme où
$\mathcal{P}$ est un bon module différentiel gradué à droite sur $\mathcal{A}$. Pour
prouver le lemme, nous devons montrer que
$\mathcal{P} \otimes_\mathcal{A} \mathcal{N} \to
\mathcal{M} \otimes_\mathcal{A} \mathcal{N}$ est un
quasi-isomorphisme. Le cône $C$ du morphisme
$\mathcal{P} \to \mathcal{M}$ est un module différentiel gradué acyclique
à droite sur $\mathcal{A}$. Ainsi,
$C \otimes_\mathcal{A} \mathcal{N}$ est acyclique, car $\mathcal{N}$
est supposé bon comme module différentiel gradué à gauche sur $\mathcal{A}$.
Puisque $C \otimes_\mathcal{A} \mathcal{N}$ est le cône du
morphisme $\mathcal{P} \otimes_\mathcal{A} \mathcal{N} \to
\mathcal{M} \otimes_\mathcal{A} \mathcal{N}$ comme complexe de
$\mathcal{O}$-modules, nous concluons.
\end{proof}

\begin{lemma}
\label{lemma-compose-tensor} Soit
$(\mathcal{C}, \mathcal{O})$ un site annelé. Soient
$\mathcal{A}$, $\mathcal{A}'$ et $\mathcal{A}''$ des
$\mathcal{O}$-algèbres différentielles graduées. Soient $\mathcal{N}$ et $\mathcal{N}'$ des bimodules
différentiels gradués respectivement sur $(\mathcal{A}, \mathcal{A}')$
et $(\mathcal{A}', \mathcal{A}'')$. Supposons
que l'application canonique
$$
\mathcal{N} \otimes_{\mathcal{A}'}^\mathbf{L} \mathcal{N}'
\longrightarrow
\mathcal{N} \otimes_{\mathcal{A}'} \mathcal{N}'
$$
dans $D(\mathcal{A}'', \text{d})$ soit un quasi-isomorphisme.
Alors nous
avons $$
(\mathcal{M}
\otimes_\mathcal{A}^\mathbf{L} \mathcal{N})
\otimes_{\mathcal{A}'}^\mathbf{L} \mathcal{N}'
=
\mathcal{M}
\otimes_\mathcal{A}^\mathbf{L}
(\mathcal{N} \otimes_{\mathcal{A}'} \mathcal{N}')
$$ en
tant que foncteurs $D(\mathcal{A}, \text{d}) \to D(\mathcal{A}'', \text{d})$.
\end{lemma}

\begin{proof}
Choisissons un bon $\mathcal{A}$-module différentiel gradué
$\mathcal{P}$ et un quasi-isomorphisme $\mathcal{P} \to \mathcal{M}$, voir
lemme \ref{lemma-resolve}. Alors
$$
\mathcal{M}
\otimes_\mathcal{A}^\mathbf{L}
(\mathcal{N} \otimes_{\mathcal{A}'} \mathcal{N}') =
\mathcal{P} \otimes_\mathcal{A} \mathcal{N}
\otimes_{\mathcal{A}'} \mathcal{N}'
$$ et nous
avons
$$
(\mathcal{M}
\otimes_\mathcal{A}^\mathbf{L} \mathcal{N})
\otimes_{\mathcal{A}'}^\mathbf{L} \mathcal{N}' =
(\mathcal{P} \otimes_\mathcal{A} \mathcal{N})
\otimes_{\mathcal{A}'}^\mathbf{L} \mathcal{N}'
$$ Ainsi, nous
devons montrer que l'application canonique
$$
(\mathcal{P} \otimes_\mathcal{A} \mathcal{N})
\otimes_{\mathcal{A}'}^\mathbf{L} \mathcal{N}'
\longrightarrow
\mathcal{P} \otimes_\mathcal{A} \mathcal{N}
\otimes_{\mathcal{A}'} \mathcal{N}'
$$ est un
quasi-isomorphisme. Choisissons un quasi-isomorphisme
$\mathcal{Q} \to \mathcal{N}'$ où $\mathcal{Q}$ est un
bon module différentiel gradué à gauche sur $\mathcal{A}'$
(lemme \ref{lemma-resolve}). Par
le lemme \ref{lemma-good-on-other-side}, le
morphisme ci-dessus, considéré dans la catégorie dérivée des $\mathcal{O}$-modules,
est le morphisme $$
\mathcal{P} \otimes_\mathcal{A} \mathcal{N}
\otimes_{\mathcal{A}'} \mathcal{Q}
\longrightarrow
\mathcal{P} \otimes_\mathcal{A} \mathcal{N}
\otimes_{\mathcal{A}'} \mathcal{N}'
$$
Puisque $\mathcal{N} \otimes_{\mathcal{A}'} \mathcal{Q} \to
\mathcal{N} \otimes_{\mathcal{A}'} \mathcal{N}'$ est un quasi-isomorphisme
par hypothèse et $\mathcal{P}$ est un bon
$\mathcal{A}$-module différentiel gradué, cette application est un quasi-isomorphisme
par le lemme \ref{lemma-tensor-symmetry} (le côté gauche et
droit calculent $\mathcal{P} \otimes_\mathcal{A}^\mathbf{L} (\mathcal{N}
\otimes_{\mathcal{A}'} \mathcal{Q})$
et $\mathcal{P} \otimes_\mathcal{A}^\mathbf{L} (\mathcal{N}
\otimes_{\mathcal{A}'} \mathcal{N}')$ ou l'on peut simplement
répéter l'argument dans la preuve du lemme).
\end{proof}






\section{Image directe dérivée}
\label{section-derived-pushforward}

\noindent
L'existence de suffisamment d'objets K-injectifs garantit que nous pouvons prendre le foncteur
dérivé à droite de tout foncteur exact défini sur la catégorie homotopique.

\begin{lemma}
\label{lemma-right-derived}
Soit $(\mathcal{C}, \mathcal{O})$ un site annelé. Soit
$(\mathcal{A}, \text{d})$ un faisceau d'algèbres différentielles graduées sur
$(\mathcal{C}, \mathcal{O})$. Alors tout foncteur exact
$$
T : K(\textit{Mod}(\mathcal{A}, \text{d})) \longrightarrow \mathcal{D}
$$ de
catégories triangulées possède une extension dérivée à droite
$RT : D(\mathcal{A}, \text{d}) \to \mathcal{D}$
dont la valeur sur un
$\mathcal{A}$-module différentiel gradué $\mathcal{I}$ injectif gradué et K-injectif
est $T(\mathcal{I})$.
\end{lemma}

\begin{proof}
D'après le théorème \ref{theorem-qis-into-dg-injective},
pour tout module différentiel gradué à droite sur $\mathcal{A}$,
$\mathcal{M}$, il existe un
quasi-isomorphisme $\mathcal{M} \to \mathcal{I}$ où $\mathcal{I}$
est un module différentiel
gradué injectif gradué et K-injectif sur $\mathcal{A}$.
Par conséquent, d'après Catégories dérivées, lemme \ref{derived-lemma-find-existence-computes},
il suffit de montrer
que, pour tout quasi-isomorphisme $\mathcal{I} \to \mathcal{I}'$
de $\mathcal{A}$-modules
différentiels gradués tous deux
injectifs gradués et K-injectifs, $T(\mathcal{I}) \to T(\mathcal{I}')$
est un isomorphisme. En effet, $\mathcal{I} \to \mathcal{I}'$
est un isomorphisme dans $K(\textit{Mod}(\mathcal{A}, \text{d}))$,
comme on le voit par exemple par le lemme \ref{lemma-hom-derived},
ou par le lemme \ref{lemma-second-property-dg-injective}.
\end{proof}

\noindent
Nous avons déjà vu un certain nombre de foncteurs auxquels cela s'applique.
Voici deux exemples.

\begin{definition}
\label{definition-pushforward}
Hom interne dérivé et image directe dérivée.
\begin{enumerate}
\item Soit $(\mathcal{C}, \mathcal{O})$ un site annelé. Soient
$\mathcal{A}$ et $\mathcal{B}$ des $\mathcal{O}$-algèbres différentielles graduées.
Soit $\mathcal{N}$ un
$(\mathcal{A}, \mathcal{B})$-bimodule différentiel gradué. L'extension dérivée à droite
$$
R\SheafHom_\mathcal{B}(\mathcal{N}, -) :
D(\mathcal{B}, \text{d})
\longrightarrow
D(\mathcal{A}, \text{d})
$$
du foncteur Hom interne $\SheafHom_\mathcal{B}^{dg}(\mathcal{N}, -)$ est
appelée {\it Hom interne dérivé}.
\item Soit $(f, f^\sharp) : (\Sh(\mathcal{C}), \mathcal{O}_\mathcal{C})
\to (\Sh(\mathcal{D}), \mathcal{O}_\mathcal{D})$ un morphisme
de topos annelés. Soit $\mathcal{A}$ une
$\mathcal{O}_\mathcal{C}$-algèbre différentielle graduée. Soit $\mathcal{B}$ une
$\mathcal{O}_\mathcal{D}$-algèbre différentielle graduée. Soit
$\varphi : \mathcal{B} \to f_*\mathcal{A}$ un homomorphisme de
$\mathcal{O}_\mathcal{D}$-algèbres différentielles graduées. L'extension
dérivée à droite
$$ Rf_*
:
D(\mathcal{A}, \text{d})
\longrightarrow
D(\mathcal{B}, \text{d})
$$ de
l'image directe $f_*$ est appelée {\it image directe dérivée}.
\end{enumerate}
\end{definition}

\noindent
Le foncteur
$Rf_* : D(\mathcal{A}, \text{d}) \to D(\mathcal{B}, \text{d})$ coïncide
avec l'image directe dérivée des complexes sous-jacents de
$\mathcal{O}$-modules ; voir le lemme \ref{lemma-pushforward-agrees}.

\medskip\noindent Ces
foncteurs sont les adjoints de l'image inverse dérivée et du
produit tensoriel dérivé.

\begin{lemma}
\label{lemma-derived-adjoint-tensor-hom}
Soit $(\mathcal{C}, \mathcal{O})$ un site annelé. Soient
$\mathcal{A}$ et $\mathcal{B}$ des $\mathcal{O}$-algèbres différentielles graduées.
Soit $\mathcal{N}$ un
$(\mathcal{A}, \mathcal{B})$-bimodule différentiel gradué. Alors
$$
R\SheafHom_\mathcal{B}(\mathcal{N}, -) :
D(\mathcal{B}, \text{d})
\longrightarrow
D(\mathcal{A}, \text{d})
$$ est
adjoint à droite à
$$
- \otimes_\mathcal{A}^\mathbf{L} \mathcal{N} :
D(\mathcal{A}, \text{d})
\longrightarrow
D(\mathcal{B}, \text{d})
$$
\end{lemma}

\begin{proof}
Cela découle de Catégories dérivées, lemme
\ref{derived-lemma-pre-derived-adjoint-functors-general}
et lemme \ref{lemma-tensor-hom-adjunction-dg}.
\end{proof}

\begin{lemma}
\label{lemma-derived-adjoint-push-pull}
Soit $(f, f^\sharp) : (\Sh(\mathcal{C}), \mathcal{O}_\mathcal{C})
\to (\Sh(\mathcal{D}), \mathcal{O}_\mathcal{D})$ un
morphisme de topos annelés. Soit $\mathcal{A}$ une
$\mathcal{O}_\mathcal{C}$-algèbre différentielle graduée. Soit $\mathcal{B}$ une
$\mathcal{O}_\mathcal{D}$-algèbre différentielle graduée. Soit
$\varphi : \mathcal{B} \to f_*\mathcal{A}$ un homomorphisme de
$\mathcal{O}_\mathcal{D}$-algèbres différentielles graduées.
Alors
$$ Rf_*
:
D(\mathcal{A}, \text{d})
\longrightarrow
D(\mathcal{B}, \text{d})
$$ est
adjoint à droite à
$$
Lf^* :
D(\mathcal{B}, \text{d})
\longrightarrow
D(\mathcal{A}, \text{d})
$$
\end{lemma}

\begin{proof}
Cela découle de Catégories dérivées, lemme
\ref{derived-lemma-pre-derived-adjoint-functors-general}
et lemme \ref{lemma-adjunction-push-pull-dg}.
\end{proof}

\noindent
Ensuite, nous discuterons de ce qui se passe dans la situation considérée
à la section \ref{section-derived-pullback}.

\medskip\noindent
Soit $(f, f^\sharp) : (\Sh(\mathcal{C}), \mathcal{O}_\mathcal{C})
\to (\Sh(\mathcal{D}), \mathcal{O}_\mathcal{D})$ un morphisme
de topos annelés. Soit $\mathcal{A}$ une
$\mathcal{O}_\mathcal{C}$-algèbre différentielle graduée. Soit $\mathcal{B}$ une
$\mathcal{O}_\mathcal{D}$-algèbre différentielle graduée. Supposons donné un
morphisme
$$
\varphi : f^{-1}\mathcal{B} \to \mathcal{A}
$$ de
$f^{-1}\mathcal{O}_\mathcal{D}$-algèbres différentielles graduées. Par l'adjonction
entre restriction et extension des scalaires, cela revient à se donner
un morphisme $\varphi : f^*\mathcal{B} \to \mathcal{A}$ de
$\mathcal{O}_\mathcal{C}$-algèbres différentielles graduées ou, de manière équivalente,
$\varphi$ peut être considéré comme un morphisme
$$
\varphi : \mathcal{B} \to f_*\mathcal{A}
$$
de $\mathcal{O}_\mathcal{D}$-algèbres différentielles graduées.
Voir la remarque \ref{remark-functoriality-dga}.

\medskip\noindent
Supposons de plus que $\mathcal{A}'$ soit une seconde
$\mathcal{O}_\mathcal{C}$-algèbre différentielle graduée et que $\mathcal{N}$ soit un
$(\mathcal{A}, \mathcal{A}')$-bimodule différentiel gradué. Dans ce
contexte, nous pouvons considérer le foncteur
$$
\textit{Mod}(\mathcal{A}', \text{d})
\longrightarrow
\textit{Mod}(\mathcal{B}, \text{d}),\quad
\mathcal{M} \longmapsto
f_*\SheafHom^{dg}_{\mathcal{A}'}(\mathcal{N}, \mathcal{M})
$$
Notons que cela s'étend à un foncteur
$$
\textit{Mod}^{dg}(\mathcal{A}', \text{d})
\longrightarrow
\textit{Mod}^{dg}(\mathcal{B}, \text{d}),\quad
\mathcal{M} \longmapsto
f_*\SheafHom^{dg}_{\mathcal{A}'}(\mathcal{N}, \mathcal{M})
$$ des
catégories différentielles graduées par la discussion des
sections \ref{section-functoriality-dg} et \ref{section-dg-bimodules}. Il
s'ensuit formellement que nous obtenons également un foncteur
exact \begin{equation}
\label{equation-pushforward}
K(\textit{Mod}(\mathcal{A}', \text{d}))
\longrightarrow
K(\textit{Mod}(\mathcal{B}, \text{d})),\quad
\mathcal{M} \longmapsto
f_*\SheafHom^{dg}_{\mathcal{A}'}(\mathcal{N}, \mathcal{M})
\end{equation}
de catégories triangulées.

\begin{lemma}
\label{lemma-compose-pushforward-hom} Dans
la situation ci-dessus, notons $RT : D(\mathcal{A}', \text{d}) \to
D(\mathcal{B}, \text{d})$ l'extension dérivée à droite de
(\ref{equation-pushforward}). Alors nous avons
$$
RT(\mathcal{M}) = Rf_* R\SheafHom(\mathcal{N}, \mathcal{M})
$$
fonctoriellement en $\mathcal{M}$.
\end{lemma}

\begin{proof}
Par les lemmes \ref{lemma-tensor-hom-adjunction-dg}
et \ref{lemma-adjunction-push-pull-dg}, le foncteur
(\ref{equation-pushforward}) est adjoint à droite
du foncteur (\ref{equation-pullback}). Par Catégories dérivées, lemme
\ref{derived-lemma-pre-derived-adjoint-functors-general},
le foncteur $RT$ est adjoint à droite au foncteur
du lemme \ref{lemma-derived-tensor-product}, lequel est égal à
$Lf^*(-) \otimes_\mathcal{A}^\mathbf{L} \mathcal{N}$ par le
lemme \ref{lemma-compose-pullback-tensor}.
Par les lemmes \ref{lemma-derived-adjoint-tensor-hom}
et \ref{lemma-derived-adjoint-push-pull}, le foncteur
$Lf^*(-) \otimes_\mathcal{A}^\mathbf{L} \mathcal{N}$
est adjoint à gauche
à $Rf_* R\SheafHom(\mathcal{N}, -)$.
Nous concluons par unicité des adjoints.
\end{proof}

\begin{lemma}
\label{lemma-compose-pushforward}
Soient $(f, f^\sharp) : (\Sh(\mathcal{C}), \mathcal{O})
\to (\Sh(\mathcal{C}'), \mathcal{O}')$ et
$(g, g^\sharp) : (\Sh(\mathcal{C}'), \mathcal{O}')
\to (\Sh(\mathcal{C}''), \mathcal{O}'')$ des
morphismes de topos annelés. Soient $\mathcal{A}$, $\mathcal{A}'$ et
$\mathcal{A}''$ des algèbres différentielles graduées respectivement sur $\mathcal{O}$,
$\mathcal{O}'$ et $\mathcal{O}''$. Soient
$\varphi : \mathcal{A}' \to f_*\mathcal{A}$ et
$\varphi' : \mathcal{A}'' \to g_*\mathcal{A}'$ des
homomorphismes d'algèbres différentielles graduées respectivement sur $\mathcal{O}'$
et $\mathcal{O}''$.
Alors nous avons $R(g \circ f)_* = Rg_* \circ Rf_*
: D(\mathcal{A}, \text{d}) \to D(\mathcal{A}'', \text{d})$.
\end{lemma}

\begin{proof}
Cela découle des lemmes \ref{lemma-compose-pullback} et
\ref{lemma-derived-adjoint-push-pull} et
de l'unicité des adjoints.
\end{proof}

\begin{lemma}
\label{lemma-compose-hom} Soit
$(\mathcal{C}, \mathcal{O})$ un site annelé. Soient
$\mathcal{A}$, $\mathcal{A}'$ et $\mathcal{A}''$ des
$\mathcal{O}$-algèbres différentielles graduées. Soient $\mathcal{N}$ et $\mathcal{N}'$ des bimodules
différentiels gradués respectivement sur $(\mathcal{A}, \mathcal{A}')$
et $(\mathcal{A}', \mathcal{A}'')$. Supposons
que l'application canonique
$$
\mathcal{N} \otimes_{\mathcal{A}'}^\mathbf{L} \mathcal{N}'
\longrightarrow
\mathcal{N} \otimes_{\mathcal{A}'} \mathcal{N}'
$$
dans $D(\mathcal{A}'', \text{d})$ soit un quasi-isomorphisme.
Alors nous
avons $$
R\SheafHom_{\mathcal{A}''}
(\mathcal{N} \otimes_{\mathcal{A}'} \mathcal{N}', -)
=
R\SheafHom_{\mathcal{A}'}(\mathcal{N},
R\SheafHom_{\mathcal{A}''}(\mathcal{N}', -))
$$ en
tant que foncteurs $D(\mathcal{A}'', \text{d}) \to D(\mathcal{A}, \text{d})$.
\end{lemma}

\begin{proof}
Cela découle des lemmes \ref{lemma-compose-tensor} et
\ref{lemma-derived-adjoint-tensor-hom} et
de l'unicité des adjoints.
\end{proof}

\begin{lemma}
\label{lemma-pushforward-agrees}
Soit $(f, f^\sharp) : (\Sh(\mathcal{C}), \mathcal{O}_\mathcal{C})
\to (\Sh(\mathcal{D}), \mathcal{O}_\mathcal{D})$ un
morphisme de topos annelés. Soit $\mathcal{A}$ une
$\mathcal{O}_\mathcal{C}$-algèbre différentielle graduée. Soit $\mathcal{B}$ une
$\mathcal{O}_\mathcal{D}$-algèbre différentielle graduée. Soit
$\varphi : \mathcal{B} \to f_*\mathcal{A}$ un homomorphisme de
$\mathcal{O}_\mathcal{D}$-algèbres différentielles graduées. Le
diagramme
$$
\xymatrix{
D(\mathcal{A}, \text{d}) \ar[d]_{Rf_*} \ar[rr]_{forget} & &
D(\mathcal{O}_\mathcal{C}) \ar[d]^{Rf_*} \\
D(\mathcal{B}, \text{d}) \ar[rr]^{forget} & &
D(\mathcal{O}_\mathcal{D})
}
$$
est commutatif.
\end{lemma}

\begin{proof}
Une fois certaines catégories identifiées, ce lemme découle immédiatement
du lemme \ref{lemma-compose-pushforward}.

\medskip\noindent
Nous pouvons considérer $(\mathcal{O}_\mathcal{C}, 0)$ comme une
$\mathcal{O}_\mathcal{C}$-algèbre différentielle graduée en plaçant $\mathcal{O}_\mathcal{C}$
dans le degré $0$ et en lui conférant le différentiel zéro.
Il est clair
que nous avons $$
\textit{Mod}(\mathcal{O}_\mathcal{C},
0) = \text{Comp}(\mathcal{O}_\mathcal{C})
\quad\text{et}\quad
D(\mathcal{O}_\mathcal{C}, 0) = D(\mathcal{O}_\mathcal{C})
$$
Sous cette identification, le foncteur
d'oubli $\textit{Mod}(\mathcal{A}, \text{d}) \to
\text{Comp}(\mathcal{O}_\mathcal{C})$
est l'« image directe » $\text{id}_{\mathcal{C}, *}$
définie à la section \ref{section-functoriality-dg},
correspondant au morphisme identité
$\text{id}_\mathcal{C} : (\mathcal{C}, \mathcal{O}_\mathcal{C}) \to
(\mathcal{C}, \mathcal{O}_\mathcal{C})$ de topos annelés et au morphisme
$(\mathcal{O}_\mathcal{C}, 0) \to (\mathcal{A}, \text{d})$ de
$\mathcal{O}_\mathcal{C}$-algèbres différentielles graduées.
Puisque $\text{id}_{\mathcal{C}, *}$ est exact, nous voyons immédiatement
que $$
R\text{id}_{\mathcal{C}, *} = forget
: D(\mathcal{A}, \text{d}) \longrightarrow
D(\mathcal{O}_\mathcal{C}, 0) = D(\mathcal{O}_\mathcal{C})
$$
Le même raisonnement montre que
$$
R\text{id}_{\mathcal{D}, *} = forget :
D(\mathcal{B}, \text{d}) \longrightarrow
D(\mathcal{O}_\mathcal{D}, 0) = D(\mathcal{O}_\mathcal{D})
$$ En outre,
la construction de
$Rf_* : D(\mathcal{O}_\mathcal{C}) \to D(\mathcal{O}_\mathcal{D})$ de Cohomologie
sur les sites, section \ref{sites-cohomology-section-unbounded} concorde avec
la construction de
$Rf_* : D(\mathcal{O}_\mathcal{C}, 0) \to D(\mathcal{O}_\mathcal{D}, 0)$ dans la définition
\ref{definition-pushforward}, car les deux
foncteurs sont définis comme l'extension dérivée droite de l'image
directe sur les complexes sous-jacents de modules.
Par le lemme \ref{lemma-compose-pushforward}, nous voyons
que $Rf_* \circ R\text{id}_{\mathcal{C}, *}$
et $R\text{id}_{\mathcal{D}, *} \circ Rf_*$ sont les foncteurs dérivés
de $f_* \circ forget = forget \circ f_*$ et sont donc égaux
par unicité des adjoints.
\end{proof}

\begin{lemma}
\label{lemma-cohomology-ext}
Soit $(\mathcal{C}, \mathcal{O})$ un site annelé. Soit
$\mathcal{A}$ une $\mathcal{O}$-algèbre différentielle graduée. Soit
$\mathcal{M}$ un $\mathcal{A}$-module différentiel gradué. Soit
$n \in \mathbf{Z}$. Nous avons
$$
H^n(\mathcal{C}, \mathcal{M}) =
\Hom_{D(\mathcal{A}, \text{d})}(\mathcal{A}, \mathcal{M}[n])
$$ où
sur le côté gauche nous avons la cohomologie de $\mathcal{M}$ considérée
comme un complexe de $\mathcal{O}$-modules.
\end{lemma}

\begin{proof}
Pour prouver la formule, observons que
$$
R\Gamma(\mathcal{C}, \mathcal{M}) = \Gamma(\mathcal{C}, \mathcal{I})
$$ où
$\mathcal{M} \to \mathcal{I}$ est un quasi-isomorphisme vers un
module différentiel gradué injectif gradué et K-injectif sur
$\mathcal{A}$, noté $\mathcal{I}$ (en combinant
les lemmes \ref{lemma-right-derived} et \ref{lemma-pushforward-agrees}).
Par le lemme \ref{lemma-hom-derived} nous
avons $$
\Hom_{D(\mathcal{A}, \text{d})}(\mathcal{A}, \mathcal{M}[n])
= \Hom_{K(\textit{Mod}(\mathcal{A}, \text{d}))}(\mathcal{M}, \mathcal{I}[n]) =
H^0(\Gamma(\mathcal{C}, \mathcal{I}[n])) =
H^n(\Gamma(\mathcal{C}, \mathcal{I}))
$$ En
combinant ces deux résultats, nous obtenons notre égalité.
\end{proof}








\section{Équivalences de catégories dérivées}
\label{section-equivalence}

\noindent
Cette section est l'analogue de
l'Algèbre différentielle graduée, section \ref{dga-section-equivalence}.

\begin{lemma}
\label{lemma-qis-equivalence}
Soit $(\mathcal{C}, \mathcal{O})$ un site annelé. Si
$\varphi : \mathcal{A} \to \mathcal{B}$ est un homomorphisme de
$\mathcal{O}$-algèbres différentielles graduées qui induit un isomorphisme
sur les faisceaux de cohomologie,
alors $$
D(\mathcal{A}, \text{d}) \longrightarrow D(\mathcal{B}, \text{d}), \quad
\mathcal{M}
\longmapsto
\mathcal{M} \otimes_\mathcal{A}^\mathbf{L} \mathcal{B}
$$
est une équivalence de catégories.
\end{lemma}

\begin{proof}
Rappelons que le foncteur de restriction
$$
\textit{Mod}^{dg}(\mathcal{B}, \text{d}) \to
\textit{Mod}^{dg}(\mathcal{A}, \text{d}),\quad
\mathcal{N} \mapsto res_\varphi \mathcal{N}
$$ est
un adjoint à droite à
$$
\textit{Mod}^{dg}(\mathcal{A}, \text{d}) \to
\textit{Mod}^{dg}(\mathcal{B}, \text{d}),\quad
\mathcal{M} \mapsto \mathcal{M} \otimes_\mathcal{A} \mathcal{B}
$$ Voir la
section \ref{section-dg-bimodules}. Puisque la restriction
envoie des quasi-isomorphismes sur des quasi-isomorphismes, nous
voyons qu'elle a triviellement une extension dérivée à
gauche (donnée par restriction). Ce foncteur sera adjoint
à droite à $- \otimes_\mathcal{A}^\mathbf{L} \mathcal{B}$
par Catégories dérivées,
lemme \ref{derived-lemma-pre-derived-adjoint-functors-general}.
L'application
de l'adjonction $$
\mathcal{M} \to
res_\varphi(\mathcal{M} \otimes_\mathcal{A}^\mathbf{L} \mathcal{B})
$$
est un isomorphisme dans $D(\mathcal{A}, \text{d})$ par
notre hypothèse que $\mathcal{A} \to \mathcal{B}$ est un
quasi-isomorphisme de $\mathcal{A}$-modules différentiels gradués à gauche.
En particulier, le foncteur du lemme est pleinement
fidèle, voir Catégories, lemme \ref{categories-lemma-adjoint-fully-faithful}.
Il est clair que le noyau du foncteur de restriction
$D(\mathcal{B}, \text{d}) \to D(\mathcal{A}, \text{d})$
est nul. Nous concluons ainsi par Catégories dérivées,
lemme \ref{derived-lemma-fully-faithful-adjoint-kernel-zero}.
\end{proof}













\section{Résolutions d'algèbres différentielles graduées}
\label{section-resolution-dgas}

\noindent
Cette section est l'analogue de
l'Algèbre différentielle graduée, section \ref{dga-section-resolution-dgas}.

\medskip\noindent
Soit $(\mathcal{C}, \mathcal{O})$ un site annelé. Comme dans la
remarque \ref{remark-sheaf-graded-sets}, considérons
un faisceau d'ensembles gradués $\mathcal{S}$ sur $\mathcal{C}$.
Considérons le produit de $r$ copies
$\mathcal{S} \times \ldots \times \mathcal{S}$ comme
un faisceau d'ensembles gradués avec la
règle $\deg(s_1 \cdot \ldots \cdot s_r) = \sum \deg(s_i)$.
Pour des sections locales $s_i \in \mathcal{S}(U)$, $i = 1, \ldots, r$, nous
utilisons $s_1 \cdot \ldots \cdot s_r$
pour désigner la section correspondante
de $\mathcal{S} \times \ldots \times \mathcal{S}$
sur $U$. Notons $\mathcal{O}\langle \mathcal{S} \rangle$ la
$\mathcal{O}$-algèbre graduée libre sur $\mathcal{S}$. Plus précisément,
posons
$$
\mathcal{O}\langle \mathcal{S} \rangle =
\mathcal{O} \oplus
\bigoplus\nolimits_{r \geq 1}
\mathcal{O}[\mathcal{S} \times \ldots \times \mathcal{S}]
$$ avec une notation comme dans
la remarque \ref{remark-sheaf-graded-sets}. Cette construction
devient un faisceau de $\mathcal{O}$-algèbres graduées par
concaténation
$$
(s_1 \cdot \ldots \cdot s_r) (s'_1 \cdot \ldots \cdot s'_{r'}) = s_1
\cdot \ldots s_r \cdot s'_1 \cdot \ldots \cdot s'_{r'}
$$ Nous pouvons
munir $\mathcal{O}\langle \mathcal{S} \rangle$ d'un
différentiel en posant $\text{d}(s) = 0$ pour toute section
locale $s$ de $\mathcal{S}$, puis en prolongeant
de manière unique par la règle de Leibniz ; il
est toutefois important de considérer aussi d'autres différentiels.

\medskip\noindent En
effet, supposons qu'on nous donne un système du genre suivant :
\begin{enumerate}
\item pour $i = 0, 1, 2, \ldots$, des faisceaux d'ensembles gradués $\mathcal{S}_i$,
\item pour $i = 0, 1, 2, \ldots$, des
applications $$
\delta_{i + 1} : \mathcal{S}_{i + 1}
\longrightarrow
\mathcal{A}_i
= \mathcal{O}\langle \mathcal{S}_0 \amalg \ldots \amalg \mathcal{S}_i\rangle
$$
de faisceaux d'ensembles gradués de degré $1$ dont l'image est
contenue dans le noyau du différentiel défini par induction sur la cible.
\end{enumerate} Plus
précisément, posons d'abord
$\mathcal{A}_0 = \mathcal{O}\langle \mathcal{S}_0 \rangle$ et
munissons-la de l'unique différentiel satisfaisant à la règle de Leibniz et tel que
$\text{d}(s) = 0$ pour toute section locale $s$ de $\mathcal{S}$. Par
récurrence, supposons qu'un différentiel $\text{d}$ soit donné
sur $\mathcal{A}_i$. Nous le prolongeons alors en
l'unique différentiel sur $\mathcal{A}_{i + 1}$ satisfaisant à la règle de
Leibniz
et tel que $$
\text{d}(s) = \delta(s)
$$
où $\delta(s) = \delta_j(s)$ si $s$ appartient à la composante $\mathcal{S}_j$
de $\mathcal{S}_0 \amalg \ldots \amalg \mathcal{S}_{i + 1}$.
Cela a du sens précisément parce que $\delta(s)$ est dans le
noyau du différentiel défini par induction.

\begin{lemma}
\label{lemma-special-good}
Dans la situation ci-dessus, la $\mathcal{O}$-algèbre différentielle graduée
$$
\mathcal{A} = \colim \mathcal{A}_i
$$
possède la propriété suivante : pour tout morphisme
$(f, f^\sharp) : (\Sh(\mathcal{C}'), \mathcal{O}')
\to (\Sh(\mathcal{C}), \mathcal{O})$
de topos annelés, l'image inverse $f^*\mathcal{A}$
est plate comme module gradué sur $\mathcal{O}'$
et K-plate comme complexe de $\mathcal{O}'$-modules.
\end{lemma}

\begin{proof}
Notons que $f^*\mathcal{A} = \colim f^*\mathcal{A}_i$
et que
$$
f^*\mathcal{A}_i = \mathcal{O}'\langle
f^{-1}\mathcal{S}_0 \amalg \ldots \amalg f^{-1}\mathcal{S}_i\rangle
$$
avec un différentiel donné par la procédure inductive ci-dessus en
utilisant $f^{-1}\delta_{i + 1}$. Il suffit donc de prouver que $\mathcal{A}$
est plate comme $\mathcal{O}$-module gradué
et K-plate comme complexe de $\mathcal{O}$-modules.
Il suffit pour cela de montrer que chaque $\mathcal{A}_i$
est plate comme module gradué sur $\mathcal{O}$
et K-plate comme complexe de $\mathcal{O}$-modules ; comparer au
lemme \ref{lemma-good-direct-sum}.


\medskip\noindent
Pour $i \geq 1$,
écrivons $\mathcal{S} = \mathcal{S}_0 \amalg \ldots \amalg \mathcal{S}_i$,
de sorte que $\mathcal{A}_i = \mathcal{O}\langle \mathcal{S} \rangle$
comme $\mathcal{O}$-algèbre graduée. Nous allons filtrer cette algèbre
par des sous-modules différentiels gradués sur $\mathcal{O}$.

\medskip\noindent
Munissons $W = \mathbf{Z}_{\geq 0}^{i + 1}$ de l'ordre lexicographique.
Autrement dit, étant donnés $w = (w_0, \ldots w_i)$ et
$w' = (w'_0, \ldots, w'_i)$ dans $W$, nous posons
$$
w > w' \Leftrightarrow
\exists j,\ 0 \leq j \leq i :
w_i = w'_i,\ w_{i - 1} = w'_{i - 1},\ \ldots ,
\ w_{j + 1} = w'_{j + 1},\ w_j > w'_j
$$
et ainsi de suite. Supposons donnée une section
$s = s_1 \cdot \ldots \cdot s_r$ de
$\mathcal{S} \times \ldots \times \mathcal{S}$
sur $U$. Nous disons que le {\it poids de $s$ est défini}
si $s_a \in \mathcal{S}_{j_a}(U)$ pour un unique
$0 \leq j_a \leq i$. Dans ce cas, nous définissons le poids
$$
w(s) = (w_0(s), \ldots, w_i(s)) \in W,\quad
w_j(s) = |\{a \mid j_a = j\}|
$$
Le poids de toute section de $\mathcal{S} \times \ldots \times \mathcal{S}$
est défini localement. On vérifie aisément que l'on obtient une décomposition
en union disjointe
$$
\mathcal{S} \times \ldots \times \mathcal{S} =
\coprod\nolimits_{w \in W} \left(
\mathcal{S} \times \ldots \times \mathcal{S}\right)_w
$$
en sous-faisceaux de sections de poids donné. Bien entendu,
seuls les $w \in W$ tels que $\sum_{0 \leq j \leq i} w_j = r$ apparaissent pour un $r$ donné.
Nous obtenons en conséquence une décomposition
$$
\mathcal{A}_i = \mathcal{O} \oplus
\bigoplus\nolimits_{r \geq 1}
\bigoplus\nolimits_{w \in W}
\mathcal{O}[\left(\mathcal{S} \times \ldots \times \mathcal{S}\right)_w]
$$
Le reste de la preuve repose sur l'observation immédiate suivante :
étant donnés $r$, $w$ et une section locale $s = s_1 \cdot \ldots \cdot s_r$ de
$\left(\mathcal{S} \times \ldots \times \mathcal{S}\right)_w$,
nous avons
$$
\text{d}(s) \text{ est une section locale de } \mathcal{O} \oplus
\bigoplus\nolimits_{r' \geq 1}
\bigoplus\nolimits_{w' \in W,\ w' < w}
\mathcal{O}[\left(\mathcal{S} \times \ldots \times \mathcal{S}\right)_{w'}]
$$
En effet, dans chacune des expressions
$$
(-1)^{\deg(s_1) + \ldots + \deg(s_{a - 1})}
s_1 \cdot \ldots s_{a - 1} \cdot \delta(s_a) \cdot
s_{a + 1} \cdot \ldots \cdot s_r
$$
dont la somme donne l'élément $\text{d}(s)$, l'élément $\delta(s_a)$
est localement une combinaison $\mathcal{O}$-linéaire d'éléments
$s'_1 \cdot \ldots \cdot s'_{r'}$, avec $s'_{a'}$ dans
$\mathcal{S}_{j'_a}$ pour un certain $0 \leq j'_{a'} < j_a$, où $j_a$ est tel que
$s_a$ soit une section de $\mathcal{S}_{j_a}$.

\medskip\noindent
Cela signifie ce qui suit. Pour $w \in W$, posons
$$
F_w \mathcal{A}_i = \mathcal{O} \oplus \bigoplus\nolimits_{r \geq 1}
\bigoplus\nolimits_{w' \in W,\ w' \leq w}
\mathcal{O}[\left(\mathcal{S} \times \ldots \times \mathcal{S}\right)_{w'}]
$$
D'après l'observation ci-dessus, c'est un sous-$\mathcal{O}$-module différentiel gradué.
Nous obtenons des suites exactes courtes admissibles
$$
0 \to \colim_{w' < w} F_{w'}\mathcal{A}_i \to
F_w\mathcal{A}_i \to
\bigoplus\nolimits_{r \geq 1}
\mathcal{O}[\left(\mathcal{S} \times \ldots \times \mathcal{S}\right)_w]
\to 0
$$
de $\mathcal{A}$-modules différentiels gradués dont le différentiel
du membre de droite est nul.

\medskip\noindent
Maintenant, nous terminons la preuve par récurrence transfinie sur
l'ensemble ordonné $W$. Le complexe différentiel gradué $F_0\mathcal{A}_0$
est le facteur direct $\mathcal{O}$, qui est K-plat et plat comme module gradué.
Pour $w \in W$, si le résultat est vrai pour $F_{w'}\mathcal{A}_i$
pour $w' < w$, alors par les lemmes \ref{lemma-good-direct-sum},
\ref{lemma-good-admissible-ses} et \ref{lemma-free-graded-module-good}
nous obtenons le résultat pour $w$. Enfin, nous avons
$\mathcal{A}_i = \colim_{w \in W} F_w\mathcal{A}_i$ et
nous concluons.
\end{proof}

\begin{lemma}
\label{lemma-good-dga} Soit
$(\mathcal{C}, \mathcal{O})$ un site annelé. Soit
$(\mathcal{B}, \text{d})$ une $\mathcal{O}$-algèbre différentielle graduée. Il existe
un quasi-isomorphisme de $\mathcal{O}$-algèbres différentielles graduées
$(\mathcal{A}, \text{d}) \to (\mathcal{B}, \text{d})$ tel que
$\mathcal{A}$ soit plate comme module gradué et K-plate comme complexe de $\mathcal{O}$-modules, et que ces
propriétés restent vraies après image inverse par tout morphisme de
topos annelés.
\end{lemma}

\begin{proof} La
démonstration est identique à la première démonstration du
lemme \ref{lemma-resolve}, en utilisant cette fois des algèbres
graduées libres au lieu de modules gradués libres.

\medskip\noindent
Nous construirons $\mathcal{A} = \colim \mathcal{A}_i$ comme dans
le lemme \ref{lemma-special-good}, en construisant
$$
\mathcal{A}_0 \to \mathcal{A}_1 \to \mathcal{A}_2 \to \ldots \to \mathcal{B}
$$
Soit $\mathcal{S}_0$ le faisceau d'ensembles gradués
(remarque \ref{remark-sheaf-graded-sets})
dont la composante de degré $n$ est $\Ker(\text{d}_\mathcal{B}^n)$.
Considérons l'homomorphisme de
modules différentiels
gradués $$
\mathcal{A}_0 = \mathcal{O}\langle \mathcal{S}_0 \rangle
\longrightarrow
\mathcal{B}
$$
où le morphisme envoie une section locale $s$ de $\mathcal{S}_0$
sur la section locale correspondante de $\mathcal{A}^{\deg(s)}$ ;
celle-ci appartient au noyau du différentiel,
ce qui assure la compatibilité du morphisme avec les différentiels.
Par construction, les applications induites en cohomologie $H^n(\mathcal{A}_0) \to H^n(\mathcal{B})$
sont surjectives et le resteront pour tout $i$.

\medskip\noindent
Étape d'induction de la construction. Étant donné $\mathcal{A}_i \to \mathcal{B}$,
notons $\mathcal{S}_{i + 1}$ le faisceau d'ensembles gradués dont la composante de degré $n$
est $$
\Ker(\text{d}_{\mathcal{A}_i}^{n + 1})
\times_{\mathcal{B}^{n + 1}, \text{d}}
\mathcal{B}^n
$$
Ce faisceau est muni d'une application
canonique $$
\delta_{i + 1} : \mathcal{S}_{i + 1} \longrightarrow
\mathcal{A}_i
$$
dont l'image est contenue dans le noyau de $\text{d}_{\mathcal{A}_i}$
par construction. Par conséquent, $\mathcal{A}_{i + 1} = \mathcal{O}\langle
\mathcal{S}_0 \amalg \ldots \mathcal{S}_{i + 1}\rangle$ a un différentiel
étendant le différentiel sur $\mathcal{A}_i$, voir la discussion au début
de cette section. Le morphisme de $\mathcal{A}_{i + 1}$ vers
$\mathcal{B}$ est l'unique morphisme d'algèbres graduées qui
coïncide avec le morphisme donné sur $\mathcal{A}_i$ et envoie une
section locale $s = (a, b)$ de $\mathcal{S}_{i + 1}$ sur
$b$ dans $\mathcal{B}$. Cela est compatible avec les différentiels
précisément parce que $\text{d}(b)$ est l'image de $a$ dans $\mathcal{B}$.

\medskip\noindent
L'application $\mathcal{A} \to \mathcal{B}$ est un quasi-isomorphisme :
nous avons $H^n(\mathcal{A}) = \colim H^n(\mathcal{A}_i)$
et pour chaque $i$ l'application $H^n(\mathcal{A}_i) \to H^n(\mathcal{B})$
est surjective avec le noyau anéanti par l'application
$H^n(\mathcal{A}_i) \to H^n(\mathcal{A}_{i + 1})$ par construction.
Enfin, la condition de platitude pour $\mathcal{A}$ a été établie dans
le lemme \ref{lemma-special-good}.
\end{proof}






\section{Compléments}
\label{section-misc}

\noindent
Soit $(f, f^\sharp) : (\Sh(\mathcal{C}), \mathcal{O})
\to (\Sh(\mathcal{C}'), \mathcal{O}')$ un morphisme de topos
annelés. Soit $\mathcal{A}$ un faisceau de
$\mathcal{O}$-algèbres différentielles graduées. En utilisant la composition\footnote{Il serait
plus précis d'écrire
$F(\mathcal{A}) \otimes_\mathcal{O}^\mathbf{L} F(\mathcal{A})
\to F(\mathcal{A} \otimes_\mathcal{O} \mathcal{A}) \to F(\mathcal{A})$ si
$F$ désigne le foncteur d'oubli pour les complexes de modules sur $\mathcal{O}$. Il
convient également de noter que $\mathcal{A} \otimes_\mathcal{O} \mathcal{A}$
indique le produit tensoriel de la section \ref{section-tensor-product-dg},
de sorte que $F(\mathcal{A} \otimes_\mathcal{O} \mathcal{A})
= \text{Tot}(F(\mathcal{A}) \otimes_\mathcal{O} F(\mathcal{A}))$.
La première flèche est l'application canonique du produit tensoriel
dérivé de deux complexes de $\mathcal{O}$-modules vers leur produit tensoriel
usuel comme complexes de $\mathcal{O}$-modules.}
$$
\mathcal{A} \otimes_\mathcal{O}^\mathbf{L} \mathcal{A}
\longrightarrow
\mathcal{A} \otimes_\mathcal{O} \mathcal{A}
\longrightarrow
\mathcal{A}
$$
et le cup-produit relatif (voir Cohomologie sur les sites, remarque
\ref{sites-cohomology-remark-cup-product} et
section \ref{sites-cohomology-section-cup-product}),
nous obtenons une multiplication\footnote{Ici et dans
la suite, $Rf_* : D(\mathcal{O}) \to D(\mathcal{O}')$ est le
foncteur dérivé étudié
dans Cohomologie sur les sites, section \ref{sites-cohomology-section-unbounded} et suivantes.}
$$
\mu
: Rf_*\mathcal{A} \otimes_{\mathcal{O}'}^\mathbf{L} Rf_*\mathcal{A}
\longrightarrow
Rf_*\mathcal{A}
$$
dans $D(\mathcal{O}')$. Cette multiplication est
associative dans le
sens où le diagramme $$
\xymatrix{
Rf_*\mathcal{A} \otimes_{\mathcal{O}'}^\mathbf{L} Rf_*\mathcal{A}
\otimes_{\mathcal{O}'}^\mathbf{L} Rf_*\mathcal{A}
\ar[rr]_-{\mu \otimes 1} \ar[d]_{1 \otimes \mu}
& & Rf_*\mathcal{A} \otimes_{\mathcal{O}'}^\mathbf{L} Rf_*\mathcal{A}
\ar[d]^\mu \\
Rf_*\mathcal{A} \otimes_{\mathcal{O}'}^\mathbf{L} Rf_*\mathcal{A}
\ar[rr]^-\mu
& & Rf_*\mathcal{A}
}
$$
est commutatif dans $D(\mathcal{O}')$ ; cela
résulte de Cohomologie sur les sites, lemme \ref{sites-cohomology-lemma-cup-product-associative}.
De même, étant donné
un $\mathcal{A}$-module différentiel gradué à droite $\mathcal{M}$,
nous obtenons une
multiplication $$
\mu_\mathcal{M}
: Rf_*\mathcal{M} \otimes_{\mathcal{O}'}^\mathbf{L} Rf_*\mathcal{A}
\longrightarrow
Rf_*\mathcal{M}
$$
dans $D(\mathcal{O}')$. Cette multiplication est compatible avec $\mu$
ci-dessus au sens où le
diagramme $$
\xymatrix{
Rf_*\mathcal{M} \otimes_{\mathcal{O}'}^\mathbf{L} Rf_*\mathcal{A}
\otimes_{\mathcal{O}'}^\mathbf{L} Rf_*\mathcal{A}
\ar[rr]_-{\mu_\mathcal{M} \otimes 1} \ar[d]_{1 \otimes \mu} & &
Rf_*\mathcal{M} \otimes_{\mathcal{O}'}^\mathbf{L} Rf_*\mathcal{A}
\ar[d]^{\mu_\mathcal{M}} \\
Rf_*\mathcal{M} \otimes_{\mathcal{O}'}^\mathbf{L} Rf_*\mathcal{A}
\ar[rr]^-{\mu_\mathcal{M}} & &
Rf_*\mathcal{M}
}
$$ est
commutatif dans $D(\mathcal{O}')$ ; cela résulte encore de
Cohomologie sur les sites, lemme \ref{sites-cohomology-lemma-cup-product-associative}.

\medskip\noindent
Un cas particulier s'obtient en prenant pour $f$ le morphisme vers le topos
ponctuel $\Sh(pt)$. Dans ce cas, $\mu$ est simplement l'application du
cup-produit
$$
R\Gamma(\mathcal{C}, \mathcal{A})
\otimes_{\Gamma(\mathcal{C}, \mathcal{O})}^\mathbf{L}
R\Gamma(\mathcal{C}, \mathcal{A})
\longrightarrow
R\Gamma(\mathcal{C}, \mathcal{A}),
\quad \eta \otimes \theta \mapsto \eta \cup \theta
$$ et, de même,
$\mu_\mathcal{M}$ est l'application du cup-produit
$$
R\Gamma(\mathcal{C}, \mathcal{M})
\otimes_{\Gamma(\mathcal{C}, \mathcal{O})}^\mathbf{L}
R\Gamma(\mathcal{C}, \mathcal{A})
\longrightarrow
R\Gamma(\mathcal{C}, \mathcal{M}),
\quad \eta \otimes \theta \mapsto \eta \cup \theta
$$ En général,
via l'identification
$$
R\Gamma(\mathcal{C}, \mathcal{A}) =
R\Gamma(\mathcal{C}', Rf_*\mathcal{A})
\quad\text{et}\quad
R\Gamma(\mathcal{C}, \mathcal{M}) =
R\Gamma(\mathcal{C}', Rf_*\mathcal{M})
$$ de Cohomologie sur
les sites, remarque \ref{sites-cohomology-remark-before-Leray}, l'application
$\mu_\mathcal{M}$ induit le cup-produit en cohomologie. Pour le voir,
utilisons
Cohomologie sur les sites, lemme \ref{sites-cohomology-lemma-compose-cup-product},
avec, pour second morphisme de topos, le morphisme de $\Sh(\mathcal{C}')$ vers le
topos ponctuel comme ci-dessus.

\medskip\noindent Si
$\mathcal{M}_1 \to \mathcal{M}_2$ est un homomorphisme de
$\mathcal{A}$-modules différentiels gradués, alors le diagramme
$$
\xymatrix{
Rf_*\mathcal{M}_1 \otimes_{\mathcal{O}'}^\mathbf{L} Rf_*\mathcal{A}
\ar[rr]_-{\mu_{\mathcal{M}_1}} \ar[d] & &
Rf_*\mathcal{M}_1 \ar[d] \\
Rf_*\mathcal{M}_2 \otimes_{\mathcal{O}'}^\mathbf{L} Rf_*\mathcal{A}
\ar[rr]^-{\mu_{\mathcal{M}_2}} & &
Rf_*\mathcal{M}_2
}
$$ est commutatif dans
$D(\mathcal{O}')$ ; cela résulte de la fonctorialité du cup-produit relatif. Supposons
que nous ayons une suite exacte courte
$$ 0
\to \mathcal{M}_1 \xrightarrow{a} \mathcal{M}_2 \to \mathcal{M}_3 \to 0
$$ de modules différentiels gradués à
droite sur $\mathcal{A}$. Alors, nous affirmons que le
diagramme
$$
\xymatrix{
Rf_*\mathcal{M}_3 \otimes_{\mathcal{O}'}^\mathbf{L} Rf_*\mathcal{A}
\ar[rr]_-{\mu_{\mathcal{M}_3}} \ar[d]_{Rf_*\delta \otimes \text{id}} & &
Rf_*\mathcal{M}_3 \ar[d]^{Rf_*\delta} \\
Rf_*\mathcal{M}_1[1] \otimes_{\mathcal{O}'}^\mathbf{L} Rf_*\mathcal{A}
\ar[rr]^-{\mu_{\mathcal{M}_1[1]}} & &
Rf_*\mathcal{M}_1[1]
}
$$ est commutatif dans
$D(\mathcal{O}')$, où
$\delta : \mathcal{M}_3 \to \mathcal{M}_1[1]$ est le morphisme de
$D(\mathcal{O})$ provenant de la suite exacte courte donnée (voir Catégories
dérivées, section
\ref{derived-section-canonical-delta-functor}). Cela est clair si la suite est scindée comme
suite de modules gradués à droite sur
$\mathcal{A}$ : dans ce cas, $\delta$ peut être représenté par un homomorphisme de modules à droite sur
$\mathcal{A}$ et la discussion ci-dessus s'applique. En général, nous argumentons en utilisant
le cône sur $a$ et le diagramme
$$
\xymatrix{
\mathcal{M}_1 \ar[r]_a \ar[d] &
\mathcal{M}_2 \ar[r]_i \ar[d] & C(a)
\ar[r]_{-p} \ar[d]^q &
\mathcal{M}_1[1] \ar[d] \\
\mathcal{M}_1 \ar[r] &
\mathcal{M}_2 \ar[r] &
\mathcal{M}_3 \ar[r]^\delta &
\mathcal{M}_1[1]
}
$$ où le carré de droite est commutatif
dans $D(\mathcal{O})$, par la définition de
$\delta$ dans Catégories dérivées, lemme
\ref{derived-lemma-derived-canonical-delta-functor}. Le cône
$C(a)$ est muni d'une structure de
$\mathcal{A}$-module différentiel gradué à droite telle que $i$, $p$, $q$ soient des homomorphismes de
$\mathcal{A}$-modules différentiels gradués à droite ; voir la définition
\ref{definition-cone}. Nous savons donc, par ce qui précède,
que les diagrammes correspondants sont commutatifs
pour les morphismes $q$ et $-p$. Puisque
$q$ est un isomorphisme dans $D(\mathcal{O})$, nous concluons que la même chose
est vraie pour $\delta$ comme souhaité.

\medskip\noindent
Dans la situation ci-dessus, étant donné un $\mathcal{A}$-module différentiel gradué
à droite $\mathcal{M}$,
soit $$
\xi \in H^n(\mathcal{C}, \mathcal{M})
$$
Autrement dit, $\xi$ est une classe de cohomologie de degré $n$ de
$\mathcal{M}$, considéré comme un complexe de $\mathcal{O}$-modules.
Par le lemme \ref{lemma-cohomology-ext}, nous pouvons
construire des morphismes $$
x : \mathcal{A} \rightarrow \mathcal{M}'[n]
\quad\text{et}\quad
s : \mathcal{M} \to \mathcal{M}'
$$
de $\mathcal{A}$-modules différentiels gradués où $s$
est un quasi-isomorphisme et tels que $\xi$ soit l'image
de $1 \in H^0(\mathcal{C}, \mathcal{A})$ par le morphisme
$s[n]^{-1} \circ x$ dans la catégorie dérivée
$D(\mathcal{A}, \text{d})$ et a fortiori dans la catégorie
dérivée $D(\mathcal{O})$. Il s'ensuit que l'application correspondante
$$
\xi' = (s[n])^{-1} \circ x : \mathcal{A} \longrightarrow \mathcal{M}[n]
$$
dans $D(\mathcal{O})$ est uniquement caractérisée par les deux
propriétés suivantes :
\begin{enumerate}
\item $\xi'$ se relève en un morphisme dans $D(\mathcal{A}, \text{d})$, et
\item $\xi = \xi'(1)$ dans
$H^0(\mathcal{C}, \mathcal{M}[n]) = H^n(\mathcal{C}, \mathcal{M})$.
\end{enumerate} En
utilisant les compatibilités de $x$ et $s$ avec le cup-produit relatif
considéré ci-dessus, il s'ensuit que, pour tout\footnote{Par exemple, le
morphisme identité.} morphisme de topos annelés
$(f, f^\sharp) : (\Sh(\mathcal{C}), \mathcal{O})
\to (\Sh(\mathcal{C}'), \mathcal{O}')$, l'image directe dérivée
$$
Rf_*\xi' : Rf_*\mathcal{A} \longrightarrow Rf_*\mathcal{M}[n]
$$ de
$\xi'$ est compatible avec les applications $\mu$ et $\mu_{\mathcal{M}[n]}$ construites
ci-dessus dans le sens où le diagramme
$$
\xymatrix{
Rf_*\mathcal{A} \otimes_{\mathcal{O}'}^\mathbf{L} Rf_*\mathcal{A}
\ar[rr]_-\mu \ar[d]_{Rf_*\xi' \otimes \text{id}} & &
Rf_*\mathcal{A} \ar[d]^{Rf_*\xi'} \\
Rf_*\mathcal{M}[n] \otimes_{\mathcal{O}'}^\mathbf{L} Rf_*\mathcal{A}
\ar[rr]^-{\mu_{\mathcal{M}[n]}} & &
Rf_*\mathcal{M}[n]
}
$$ est commutatif dans
$D(\mathcal{O}')$. En utilisant cette compatibilité pour l'application
aux topos ponctuels, nous voyons en particulier que
$$
\xymatrix{
R\Gamma(\mathcal{C}, \mathcal{A})
\otimes_{\Gamma(\mathcal{C}, \mathcal{O})}^\mathbf{L}
R\Gamma(\mathcal{C}, \mathcal{A})
\ar[d]_{\xi' \otimes \text{id}} \ar[r] &
R\Gamma(\mathcal{C}, \mathcal{A}) \ar[d]^{\xi'} \\
R\Gamma(\mathcal{C}, \mathcal{M}[n])
\otimes_{\Gamma(\mathcal{C}, \mathcal{O})}^\mathbf{L}
R\Gamma(\mathcal{C}, \mathcal{A})
\ar[r] &
R\Gamma(\mathcal{C}, \mathcal{M}[n])
}
$$ est commutatif.
Combiné à $\xi'(1) = \xi$, cela implique que l'application
induite sur la cohomologie
$$
\xi' : R\Gamma(\mathcal{C}, \mathcal{A}) \to
R\Gamma(\mathcal{C}, \mathcal{M}[n]), \quad \eta \mapsto \xi \cup \eta
$$ est donnée
par le cup-produit à gauche par $\xi$, comme indiqué.














\section{Modules différentiels gradués sur une catégorie}
\label{section-modules-cohomology}

\noindent
Cette section prolonge Cohomologie
sur les sites, section \ref{sites-cohomology-section-modules-cohomology}.

\medskip\noindent
Soit $\mathcal{C}$ une catégorie. Nous pensons à $\mathcal{C}$ comme un site avec
la topologie chaotique. Soit $\mathcal{O}$ un faisceau d'anneaux sur $\mathcal{C}$.
Soit $(\mathcal{A}, \text{d})$ un faisceau de
$\mathcal{O}$-algèbres différentielles graduées. En d'autres termes, $\mathcal{O}$ est
un préfaisceau d'anneaux de la catégorie $\mathcal{C}$ et
$(\mathcal{A}, \text{d})$ est un préfaisceau de
$\mathcal{O}$-algèbres différentielles graduées sur $\mathcal{C}$,
voir Catégories, définition \ref{categories-definition-presheaf}.

\begin{definition}
\label{definition-cartesian} Dans
la situation ci-dessus, nous notons
{\it $\mathit{QC}(\mathcal{A}, \text{d})$} la
sous-catégorie pleine de $D(\mathcal{A}, \text{d})$
formée des objets $M$ tels que, pour tout
morphisme $U \to V$ de $\mathcal{C}$, l'application
canonique $$
R\Gamma(V, M) \otimes_{\mathcal{A}(V)}^\mathbf{L} \mathcal{A}(U)
\longrightarrow
R\Gamma(U, M)
$$
est un isomorphisme dans $D(\mathcal{A}(U), \text{d})$.
\end{definition}

\begin{lemma}
\label{lemma-cartesian-good-sub} Dans
la situation ci-dessus, la sous-catégorie $\mathit{QC}(\mathcal{A}, \text{d})$ est
une sous-catégorie strictement pleine, saturée et triangulée de
$D(\mathcal{A}, \text{d})$, stable par sommes directes arbitraires.
\end{lemma}

\begin{proof}
Soit $U$ un objet de $\mathcal{C}$. Puisque la topologie sur
$\mathcal{C}$ est chaotique, le foncteur $\mathcal{F} \mapsto \mathcal{F}(U)$ est
exact et commute aux sommes directes. Par conséquent, le foncteur
exact $M \mapsto R\Gamma(U, M)$ est calculé en représentant $K$ par
un $\mathcal{A}$-module différentiel
gradué quelconque $\mathcal{M}$ et en prenant $\mathcal{M}(U)$.
Ainsi, $R\Gamma(U, -)$ commute aux sommes
directes, voir le lemme \ref{lemma-derived-products}.
De même, pour un morphisme $U \to V$ de $\mathcal{C}$,
le foncteur de produit
tensoriel dérivé $- \otimes_{\mathcal{O}(A)}^\mathbf{L} \mathcal{A}(U)
: D(\mathcal{A}(V)) \to D(\mathcal{A}(U))$
est exact et commute aux sommes directes. Le
lemme découle directement de ces observations ; les détails sont omis.
\end{proof}

\begin{remark}
\label{remark-cartesian-brown}
Comme ci-dessus, soient $\mathcal{C}$ une catégorie considérée comme
un site muni de la topologie chaotique, $\mathcal{O}$ un faisceau d'anneaux sur $\mathcal{C}$
et $(\mathcal{A}, \text{d})$ un faisceau de
$\mathcal{O}$-algèbres différentielles graduées. Alors l'analogue
de Cohomologie sur les sites,
proposition \ref{sites-cohomology-proposition-cartesian-brown},
vaut pour $\mathit{QC}(\mathcal{A}, \text{d})$ avec
presque exactement la même
démonstration : \begin{enumerate}
\item tout foncteur cohomologique
contravariant $H : \mathit{QC}(\mathcal{A}, \text{d}) \to \textit{Ab}$
qui transforme les sommes directes en produits est
représentable, \item tout
foncteur exact $F : \mathit{QC}(\mathcal{A}, \text{d}) \to \mathcal{D}$
de catégories triangulées qui transforme les sommes directes en
sommes directes possède un adjoint
à droite exact, et \item le
foncteur d'inclusion $\mathit{QC}(\mathcal{A}, \text{d}) \to D(\mathcal{A}, \text{d})$
possède un
adjoint à droite exact. \end{enumerate}
Si nous en avons besoin, nous le démontrerons ici.
\end{remark}

\noindent
Soit $u : \mathcal{C}' \to \mathcal{C}$ un foncteur entre catégories. Si nous
considérons $\mathcal{C}$ et $\mathcal{C}'$ comme des sites avec la
topologie chaotique, alors $u$ est un foncteur continu et cocontinu. On
obtient donc un morphisme $g : \Sh(\mathcal{C}') \to \Sh(\mathcal{C})$ de
topos, voir Sites, lemme \ref{sites-lemma-cocontinuous-morphism-topoi}. En
outre, supposons donnés des faisceaux d'anneaux $\mathcal{O}$ sur $\mathcal{C}$ et
$\mathcal{O}'$ sur $\mathcal{C}'$ et une application
$g^\sharp : g^{-1}\mathcal{O} \to \mathcal{O}'$.
Notons simplement le morphisme correspondant de topos annelés
$g : (\Sh(\mathcal{C}'), \mathcal{O}') \to (\Sh(\mathcal{C}), \mathcal{O})$, voir
Modules sur les sites, section \ref{sites-modules-section-ringed-topoi}.
Enfin, supposons que $(\mathcal{A}, \text{d})$ soit un faisceau de
$\mathcal{O}$-algèbres différentielles graduées, que
$(\mathcal{A}', \text{d})$ soit un faisceau de
$\mathcal{O}'$-algèbres différentielles graduées, et que nous ayons un morphisme
$\varphi : g^*\mathcal{A} \to \mathcal{A}'$ d'algèbres différentielles
graduées sur $\mathcal{O}'$ (voir la section \ref{section-functoriality-dg}).

\begin{lemma}
\label{lemma-cartesian-functorial}
Soient $g : (\Sh(\mathcal{C}'), \mathcal{O}') \to (\Sh(\mathcal{C}), \mathcal{O})$
et $\varphi : g^*\mathcal{A} \to \mathcal{A}'$ comme
ci-dessus. Le foncteur
$Lg^* : D(\mathcal{A}, \text{d}) \to D(\mathcal{A}', \text{d})$ envoie
alors $\mathit{QC}(\mathcal{A}, \text{d})$ dans
$\mathit{QC}(\mathcal{A}', \text{d})$.
\end{lemma}

\begin{proof}
Soit $U' \in \Ob(\mathcal{C}')$, d'image $U = u(U')$ dans $\mathcal{C}$.
Notons $pt$ la catégorie ayant un seul objet et un seul morphisme. Considérons
les topos annelés $(\Sh(pt), \mathcal{O}'(U'))$ et $(\Sh(pt), \mathcal{O}(U))$,
munis des algèbres différentielles graduées
$\mathcal{A}'(U)$ et $\mathcal{A}(U)$. Bien sûr, nous identifions
la catégorie dérivée des modules différentiels gradués sur ceux-ci avec
$D(\mathcal{A}'(U'), \text{d})$ et
$D(\mathcal{A}(U), \text{d})$. Ensuite, nous
avons un diagramme commutatif de topos annelés
$$
\xymatrix{
(\Sh(pt), \mathcal{O}'(U')) \ar[rr]_{U'} \ar[d] & &
(\Sh(\mathcal{C}'), \mathcal{O}') \ar[d]^g \\
(\Sh(pt), \mathcal{O}(U)) \ar[rr]^U & &
(\Sh(\mathcal{C}), \mathcal{O})
}
$$ Chacun
est doté d'algèbres différentielles graduées correspondantes.
L'image inverse par le morphisme horizontal inférieur envoie $M$
dans $D(\mathcal{A}, \text{d})$ sur $R\Gamma(U, K)$,
considéré comme un objet de $D(\mathcal{A}(U), \text{d})$.
L'image inverse par la flèche verticale gauche envoie $M$
sur $M \otimes_{\mathcal{A}(U)}^\mathbf{L} \mathcal{A}'(U')$.
Les deux parcours du diagramme donnent le
même résultat (lemme \ref{lemma-compose-pullback}), d'où
$$
R\Gamma(U', Lg^*K) =
R\Gamma(U, K) \otimes_{\mathcal{A}(U)}^\mathbf{L} \mathcal{A}'(U')
$$
Enfin, soit $f' : U' \to V'$ un morphisme de $\mathcal{C}'$, et notons
$f = u(f') : U = u(U') \to V = u(V')$ son image dans $\mathcal{C}$. Si
$K$ est dans $\mathit{QC}(\mathcal{A}, \text{d})$ alors nous avons
\begin{align*}
R\Gamma(V', Lg^*K) \otimes_{\mathcal{A}'(V')}^\mathbf{L} \mathcal{A}'(U') &
=
R\Gamma(V, K) \otimes_{\mathcal{A}(V)}^\mathbf{L} \mathcal{A}'(V')
\otimes_{\mathcal{A}'(V')}^\mathbf{L} \mathcal{A}'(U') \\ &
=
R\Gamma(V, K) \otimes_{\mathcal{A}(V)}^\mathbf{L} \mathcal{A}'(U') \\ &
=
R\Gamma(V, K) \otimes_{\mathcal{A}(V)}^\mathbf{L} \mathcal{A}(U)
\otimes_{\mathcal{A}(U)}^\mathbf{L} \mathcal{A}'(U') \\ &
=
R\Gamma(U, K) \otimes_{\mathcal{A}(U)}^\mathbf{L} \mathcal{A}'(U') \\ &
=
R\Gamma(U', Lg^*K)
\end{align*} comme
souhaité. Ici, nous avons utilisé l'observation ci-dessus pour $U'$ et $V'$.
\end{proof}






\section{Modules différentiels gradués sur une catégorie, bis}
\label{section-modules-cohomology-bis}

\noindent
Nous développons quelques résultats supplémentaires sur la notion de modules
quasi cohérents introduite dans la section \ref{section-modules-cohomology}.

\begin{lemma}
\label{lemma-cartesian-cofinal-subcategory}
Soient $\mathcal{C}, \mathcal{O}, \mathcal{A}$ comme
dans la section \ref{section-modules-cohomology}.
Soit $\mathcal{C}' \subset \mathcal{C}$ une sous-catégorie
pleine telle que, pour tout $U \in \Ob(\mathcal{C})$, la catégorie
$U/\mathcal{C}'$ des flèches $U \to U'$ soit cofiltrée. Notons
$\mathcal{O}', \mathcal{A}'$ les restrictions de
$\mathcal{O}, \mathcal{A}$ à $\mathcal{C}'$. Les
restrictions induisent alors une équivalence
$\mathit{QC}(\mathcal{A}, \text{d}) \to \mathit{QC}(\mathcal{A}', \text{d})$.
\end{lemma}

\begin{proof}
Construisons un quasi-inverse de ce foncteur. Soit
$M'$ un objet de $\mathit{QC}(\mathcal{A}', \text{d})$. Nous pouvons représenter
$M'$ par un bon module différentiel gradué $\mathcal{M}'$, voir lemme
\ref{lemma-resolve}. Ensuite, pour chaque $U' \in \Ob(\mathcal{C}')$, le
$\mathcal{A}'(U')$-module différentiel gradué $\mathcal{M}'(U)$ est K-plat et plat comme module
gradué et, pour chaque morphisme $U'_1 \to U'_2$ de
$\mathcal{C}'$, l'application
$$
\mathcal{M}'(U'_2) \otimes_{\mathcal{A}'(U'_2)} \mathcal{A}'(U'_1)
\longrightarrow
\mathcal{M}'(U'_1)
$$ est un
quasi-isomorphisme (puisque la source représente le produit tensoriel dérivé). Considérons
le
$\mathcal{A}$-module différentiel gradué $\mathcal{M}$ défini par
$$
\mathcal{M}(U) =
\colim_{U \to U' \in U/\mathcal{C}'}
\mathcal{M}'(U') \otimes_{\mathcal{A}'(U')} \mathcal{A}(U)
$$ Il s'agit d'une limite inductive
filtrante de complexes d'après l'hypothèse du lemme. Puisque
$M'$ est dans $\mathit{QC}(\mathcal{A}', \text{d})$ tous les morphismes de transition
dans le système sont des quasi-isomorphismes. Puisque les limites inductives
filtrantes sont exactes, nous voyons
que $\mathcal{M}(U)$ dans $D(\mathcal{A}(U), \text{d})$
est isomorphe à $\mathcal{M}'(U') \otimes_{\mathcal{A}'(U')} \mathcal{A}(U)$
pour tout morphisme $U \to U'$ avec $U' \in \Ob(\mathcal{C}')$.

\medskip\noindent
Nous affirmons que $\mathcal{M}$ est dans
$\mathit{QC}(\mathcal{A}, \text{d})$ : à savoir, étant donné $U \to V$
dans $\mathcal{C}$, nous choisissons une application $V \to V'$ avec $V' \in \Ob(\mathcal{C}')$.
Par ce qui précède, nous voyons que l'application $\mathcal{M}(V) \to \mathcal{M}(U)$
est
identifiée avec l'application $$
\mathcal{M}'(V') \otimes_{\mathcal{A}'(V')} \mathcal{A}(V)
\longrightarrow
\mathcal{M}'(V') \otimes_{\mathcal{A}'(V')} \mathcal{A}(U)
$$
Puisque $\mathcal{M'}(V')$ est
K-plat comme $\mathcal{A}'(V')$-module différentiel gradué, nous
obtenons l'affirmation.

\medskip\noindent
Le morphisme naturel $\mathcal{M}|_{\mathcal{C}'} \to \mathcal{M}'$
est un isomorphisme dans $D(\mathcal{A}', d)$, d'après ce qui précède.

\medskip\noindent À
l'inverse, si nous avons un objet $E$ de $\mathit{QC}(\mathcal{A}, \text{d})$, alors
nous le représentons par un bon module différentiel gradué $\mathcal{E}$. En
définissant $\mathcal{M}' = \mathcal{E}|_{\mathcal{C}'}$ (c'est un autre bon module
différentiel gradué) nous voyons qu'il y a une application
$$
\mathcal{E} \to \mathcal{M}
$$ qui,
sur $U$ dans $\mathcal{C}$, est donnée par le morphisme
$$
\mathcal{E}(U)
\longrightarrow
\colim_{U \to U' \in U/\mathcal{C}'}
\mathcal{E}(U') \otimes_{\mathcal{A}'(U')} \mathcal{A}(U)
$$ qui est
un quasi-isomorphisme pour la même raison. Ainsi, la restriction et la
construction ci-dessus sont des foncteurs quasi inverses, comme souhaité.
\end{proof}

\begin{lemma}
\label{lemma-cartesian-qis-equivalence}
Soient $\mathcal{C}, \mathcal{O}$ comme dans
la section \ref{section-modules-cohomology}. Si
$\varphi : \mathcal{A} \to \mathcal{B}$ est un homomorphisme de
$\mathcal{O}$-algèbres différentielles graduées qui induit un
isomorphisme sur les faisceaux de cohomologie, alors
l'équivalence $D(\mathcal{A}, \text{d}) \to D(\mathcal{B}, \text{d})$
du lemme \ref{lemma-qis-equivalence} induit une équivalence
$\mathit{QC}(\mathcal{A}, \text{d}) \to
\mathit{QC}(\mathcal{B}, \text{d})$.
\end{lemma}

\begin{proof}
Il suffit de montrer que, pour un morphisme $U \to V$ de $\mathcal{C}$ et un objet
$M$ de $D(\mathcal{A}, \text{d})$, les conditions suivantes sont équivalentes :
\begin{enumerate}
\item $R\Gamma(V, M) \otimes_{\mathcal{A}(V)}^\mathbf{L} \mathcal{A}(U)
\to \Gamma(U, M)$ est un isomorphisme dans $D(\mathcal{A}(U), \text{d})$,
et \item $R\Gamma(V, M \otimes_\mathcal{A}^\mathbf{L} \mathcal{B})
\otimes_{\mathcal{B}(V)}^\mathbf{L} \mathcal{B}(U)
\to \Gamma(U, M \otimes_\mathcal{A}^\mathbf{L} \mathcal{B})$ est un
isomorphisme dans $D(\mathcal{B}(U), \text{d})$.
\end{enumerate} Comme
la topologie sur $\mathcal{C}$ est chaotique, cela se résume simplement au
fait que $\mathcal{A}(U) \to \mathcal{B}(U)$ et
$\mathcal{A}(V) \to \mathcal{B}(V)$ sont des quasi-isomorphismes. Les
détails sont omis.
\end{proof}






\section{Systèmes projectifs d'algèbres différentielles graduées}
\label{section-qc-inverse-system}

\noindent
Dans cette section, nous examinons le cas particulier suivant de la situation
décrite dans la section \ref{section-modules-cohomology} :
\begin{enumerate}
\item $\mathcal{C}$ est la catégorie $\mathbf{N}$ avec
un morphisme unique $i \to j$ si et seulement si $i \leq j$,
\item $\mathcal{O}$ est le (pré)faisceau constant d'anneaux de valeur l'anneau
donné $R$.
\end{enumerate} Dans
ce cadre, un faisceau $\mathcal{A}$ de
$\mathcal{O}$-algèbres différentielles graduées est la même chose qu'un système projectif
$(A_n)$ de $R$-algèbres différentielles graduées. Un
faisceau $\mathcal{M}$ de $\mathcal{A}$-modules différentiels gradués est
la même chose qu'un système projectif $(M_n)$ où $M_n$ est un
$A_n$-module différentiel gradué et les morphismes de transition
$M_{n + 1} \to M_n$ sont des morphismes de $A_{n + 1}$-modules.
Dans cette section, nous omettons le différentiel de la
notation et écrivons $\mathit{QC}(\mathcal{A})$ pour la catégorie
notée $\mathit{QC}(\mathcal{A}, \text{d})$ dans
la définition \ref{definition-cartesian}.

\medskip\noindent
Supposons que $\mathcal{B} = (B_n)$ soit un deuxième système projectif de
$R$-algèbres différentielles graduées. Étant donné un morphisme $\varphi : (A_n) \to (B_n)$
de pro-objets, construisons un foncteur
exact de $\mathit{QC}(\mathcal{A})$ vers $\mathit{QC}(\mathcal{B})$.
D'après Catégories, exemple
\ref{categories-example-pro-morphism-inverse-systems}, le
morphisme $\varphi$ est donné par une suite d'entiers
$\ldots \geq m(3) \geq m(2) \geq m(1)$ et un
diagramme commutatif
$$
\xymatrix{
\ldots \ar[r] &
A_{m(3)} \ar[d]^{\varphi_3} \ar[r] &
A_{m(2)} \ar[d]^{\varphi_2} \ar[r] &
A_{m(1)} \ar[d]^{\varphi_1} \\
\ldots \ar[r] & B_3
\ar[r] & B_2
\ar[r] &
B_1
}
$$ de
$R$-algèbres différentielles graduées. Étant donné un bon faisceau de
$\mathcal{A}$-modules différentiels gradués $\mathcal{M} = (M_n)$
représentant un objet de $\mathit{QC}(\mathcal{A})$,
posons
$$ N_n
= M_{m(n)} \otimes_{A_{m(n)}} B_n
$$ Ce
système projectif détermine un objet de
$\mathit{QC}(\mathcal{B})$ parce que les
$A_{m(n)}$-modules $M_{m(n)}$ sont K-plats ; détails omis. Nous
laissons également au lecteur le soin de montrer que le foncteur
résultant est indépendant des choix faits dans sa construction.

\begin{lemma}
\label{lemma-pro-isomorphic} Dans
la situation ci-dessus, supposons que $\mathcal{A} = (A_n)$ et
$\mathcal{B} = (B_n)$ soient des systèmes projectifs de
$R$-algèbres différentielles graduées. Si $\varphi : (A_n) \to (B_n)$ est un isomorphisme
de pro-objets, alors le foncteur $\mathit{QC}(\mathcal{A}) \to
\mathit{QC}(\mathcal{B})$ construit ci-dessus est
une équivalence.
\end{lemma}

\begin{proof}
Soit $\psi : (B_n) \to (A_n)$ un morphisme de pro-objets inverse de
$\varphi$. D'après la discussion de Catégories,
exemple
\ref{categories-example-pro-morphism-inverse-systems}, nous pouvons supposer que
$\varphi$ est donné par un système de morphismes comme ci-dessus,
tandis que $\psi$ est donné par une suite $n(1) < n(2) < \ldots$
et un diagramme
commutatif $$
\xymatrix{
\ldots \ar[r]
& B_{n(3)} \ar[d]^{\psi_3} \ar[r]
& B_{n(2)} \ar[d]^{\psi_2} \ar[r]
& B_{n(1)} \ar[d]^{\psi_1} \\
\ldots \ar[r] &
A_3 \ar[r] &
A_2 \ar[r] &
A_1
}
$$
de $R$-algèbres différentielles graduées.
Puisque $\varphi \circ \psi = \text{id}$, nous pouvons, après avoir
éventuellement augmenté les valeurs des fonctions $n(\cdot)$
et $m(\cdot)$, supposer que $B_{n(m(i))} \to A_{m(i)} \to B_i$
est l'application de transition. Il s'ensuit que la composition
des
foncteurs $$
\mathit{QC}(\mathcal{B}) \to
\mathit{QC}(\mathcal{A}) \to
\mathit{QC}(\mathcal{B})
$$
envoie un bon faisceau
de $\mathcal{B}$-modules différentiels gradués $\mathcal{N} = (N_n)$
au système projectif $\mathcal{N}' = (N'_i)$ à valeurs
$$
N'_i = N_{n(m(i))} \otimes_{B_{n(m(i))}} B_i
$$ qui
est canoniquement quasi-isomorphe à $\mathcal{N}$ précisément
parce que $\mathcal{N}$ est un objet de
$\mathit{QC}(\mathcal{B})$ et que
$N_j$ est un module différentiel gradué K-plat pour
tout $j$. Le même argument s'applique à la composition
dans l'autre sens, ce qui conclut la démonstration.
\end{proof}

\noindent
Soient $\mathcal{C} = \mathbf{N}$ et $\mathcal{O}$ le faisceau constant de
valeur un anneau $R$ ; soit $\mathcal{A}$ donné par un système projectif
$(A_n)$ de $R$-algèbres différentielles graduées. Supposons donnés deux
modules différentiels gradués à gauche sur $\mathcal{A}$, $\mathcal{N}$
et $\mathcal{N}'$, par des systèmes projectifs $(N_n)$ et $(N'_n)$.
Ainsi, chaque $N_n$ et $N'_n$ est un module différentiel gradué à gauche sur $A_n$.
Disons provisoirement que $(N_n)$ et $(N'_n)$ sont
{\it pro-isomorphes dans la catégorie dérivée}
s'il existe une suite d'entiers
$$ 1 = n_0
< n_1 < n_2 < n_3 < \ldots
$$ et des
morphismes
$$
N_{n_{2i}} \to N'_{n_{2i - 1}}
\quad\text{dans}\quad
D(A_{n_{2i}}^{opp}, \text{d})
$$
et
$$
N'_{n_{2i + 1}} \to N'_{n_{2i}}
\quad\text{dans}\quad
D(A_{n_{2i + 1}}^{opp}, \text{d})
$$ de
sorte que les compositions $N_{n_{2i}} \to N_{n_{2i - 2}}$ et
$N'_{n_{2i + 1}} \to N'_{2i - 1}$ sont données par les morphismes de
transition des systèmes respectifs.

\begin{lemma}
\label{lemma-pro-isomorphic-systems-tensor} Si
$(N_n)$ et $(N'_n)$ sont pro-isomorphes dans la catégorie dérivée au sens
défini ci-dessus,
alors pour chaque objet $(M_n)$ de $D(\mathbf{N}, \mathcal{A})$ nous avons
$$
R\lim (M_n \otimes_{A_n}^\mathbf{L} N_n) =
R\lim (M_n \otimes_{A_n}^\mathbf{L} N'_n)
$$
dans $D(R)$.
\end{lemma}

\begin{proof}
L'hypothèse implique que le système projectif
$(M_n \otimes_{A_n}^\mathbf{L} N_n)$ de $D(R)$ est pro-isomorphe (au sens
habituel) au système projectif
$(M_n \otimes_{A_n}^\mathbf{L} N'_n)$ de $D(R)$. Le résultat découle
donc de Compléments
sur l'algèbre, lemme \ref{more-algebra-lemma-pro-isomorphism-Rlim}.
\end{proof}

\begin{lemma}
\label{lemma-different-tensors}
\begin{reference} Il
s'agit d'une variante de \cite[Lemma 3.5.4]{BS}.
\end{reference}
Soit $R$ un anneau. Soient $f_1, \ldots, f_r \in R$.
Soit $K_n$ le complexe de Koszul sur $f_1^n, \ldots, f_r^n$,
considéré comme une $R$-algèbre différentielle
graduée. Soit $(M_n)$ un objet de $D(\mathbf{N}, (K_n))$.
Alors, pour tout $t \geq 1$, on
a $$
R\lim (M_n \otimes_R^\mathbf{L} K_t) =
R\lim (M_n \otimes_{K_n}^\mathbf{L} K_t)
$$
dans $D(R)$.
\end{lemma}

\begin{proof}
Fixons $t \geq 1$. Pour $n \geq t$, notons ${}_nK_t$
la $R$-algèbre différentielle graduée $K_t$ considérée comme un module différentiel gradué
à gauche sur $K_n$. Observons que
$$
M_n \otimes_R^\mathbf{L} K_t =
M_n \otimes_{K_n}^\mathbf{L} (K_n \otimes_R^\mathbf{L} K_t) =
M_n \otimes_{K_n}^\mathbf{L} (K_n \otimes_R K_t)
$$
Ainsi, d'après le lemme \ref{lemma-pro-isomorphic-systems-tensor},
il suffit de montrer que
$({}_nK_t)$ et $(K_n \otimes_R K_t)$ sont pro-isomorphes dans la
catégorie dérivée. Les applications de multiplication
$$
K_n \otimes_R K_t \longrightarrow {}_nK_t
$$
sont des morphismes de modules différentiels gradués à gauche sur $K_n$. Ainsi, pour conclure,
il suffit de montrer que, pour tout $n \geq 1$, il existe $N > n$
et un morphisme
$$
{}_NK_t \longrightarrow {}_NK_n \otimes_R K_t
$$
dans $D(K_N^{opp}, \text{d})$ dont la composée avec l'application de multiplication
est l'application de transition (dans l'un ou l'autre sens). Cela est construit dans
Algèbre à puissances divisées, lemme \ref{dpa-lemma-construct-some-maps},
par une construction explicite.
\end{proof}

\begin{proposition}
\label{proposition-derived-complete} Soit
$R$ un anneau noethérien. Soit $I \subset R$ un idéal. Les trois
catégories suivantes sont canoniquement équivalentes :
\begin{enumerate}
\item Soit $\mathcal{A}$ le faisceau de $R$-algèbres sur $\mathbf{N}$ correspondant
au système projectif de $R$-algèbres $A_n = R/I^n$. La catégorie
$\mathit{QC}(\mathcal{A})$.
\item Choisissons des générateurs $f_1, \ldots, f_r$ de $I$. Soit
$\mathcal{B}$ le faisceau de $R$-algèbres différentielles graduées sur
$\mathbf{N}$ correspondant au système projectif des algèbres de
Koszul sur $f_1^n, \ldots, f_r^n$. La catégorie
$\mathit{QC}(\mathcal{B})$.
\item La sous-catégorie pleine $D_{comp}(R, I) \subset D(R)$ des objets
complets au sens dérivé, voir Compléments sur l'algèbre, définition
\ref{more-algebra-definition-derived-complete} et le texte qui suit.
\end{enumerate}
\end{proposition}

\begin{proof}
Considérons le morphisme évident
$f : (\Sh(\mathbf{N}), \mathcal{A}) \to (\Sh(pt), R)$ de
topos annelés
et considérons les foncteurs adjoints $Lf^*$ et $Rf_*$. Le premier
se restreint à un foncteur
$$ F
: D_{comp}(R, I) \longrightarrow \mathit{QC}(\mathcal{A})
$$ qui
envoie un objet $K$ de $D_{comp}(R, I)$ représenté par un complexe
K-plat $K^\bullet$ à l'objet $(K^\bullet \otimes_R R/I^n)$ de
$\mathit{QC}(\mathcal{A})$. Le second se restreint à un foncteur
$$ G :
\mathit{QC}(\mathcal{A}) \longrightarrow D_{comp}(R, I)
$$ qui envoie un objet
$(M_n^\bullet)$ de $\mathit{QC}(\mathcal{A})$ à
$R\lim M_n^\bullet$. L'objet obtenu est complet au sens dérivé, par exemple d'après Compléments
sur l'algèbre, lemme \ref{more-algebra-lemma-naive-derived-completion}. En
outre, il résulte de Compléments sur l'algèbre, proposition
\ref{more-algebra-proposition-noetherian-naive-completion-is-completion},
que $G \circ F = \text{id}$. Ainsi, pour voir que $F$ et $G$
sont des équivalences quasi-inverses, il suffit de voir que le noyau
de $G$ est nul (voir Catégories dérivées, lemme
\ref{derived-lemma-fully-faithful-adjoint-kernel-zero}).
Cependant, il ne semble pas facile de le montrer directement !

\medskip\noindent
Montrons dans ce paragraphe que $\mathit{QC}(\mathcal{A})$ et
$\mathit{QC}(\mathcal{B})$ sont équivalentes. Écrivons
$\mathcal{B} = (B_n)$, où $B_n$ est le complexe de Koszul considéré
comme un complexe de cochaînes en degrés $-r, -r + 1, \ldots, 0$. Par
D'après Algèbre à puissances divisées, remarque \ref{dpa-remark-pro-system-koszul}
(mais avec des complexes
de chaînes transformés en complexes
de cochaînes) nous pouvons trouver $1 < n_1 < n_2 < \ldots$
et des applications de $R$-algèbres
différentielles graduées $B_{n_i} \to E_i \to R/(f_1^{n_i}, \ldots, f_r^{n_i})$
et $E_i \to B_{n_{i - 1}}$
telles que $$
\xymatrix{
B_{n_1} \ar[d]
& B_{n_2} \ar[d] \ar[l]
& B_{n_3} \ar[d] \ar[l]
& \ldots \ar[l] \\
E_1 \ar[d]
& E_2 \ar[l] \ar[d]
& E_3 \ar[l] \ar[d]
& \ldots \ar[l] \\
B_1
& B_{n_1} \ar[l]
& B_{n_2} \ar[l]
& \ldots \ar[l]
}
$$
soit un diagramme commutatif de $R$-algèbres
différentielles graduées et que $E_i \to R/(f_1^{n_i}, \ldots, f_r^{n_i})$
soit un quasi-isomorphisme.
Nous concluons \begin{enumerate}
\item il existe une équivalence entre $\mathit{QC}(\mathcal{B})$
et $\mathit{QC}((E_i))$,
\item il existe une
équivalence entre $\mathit{QC}((E_i))$ et $\mathit{QC}((R/(f_1^{n_i}, \ldots, f_r^{n_i})))$,
\item il existe une équivalence entre
$\mathit{QC}((R/(f_1^{n_i}, \ldots, f_r^{n_i})))$ et
$\mathit{QC}(\mathcal{A})$.
\end{enumerate} À
savoir, pour (1) nous pouvons appliquer le lemme \ref{lemma-pro-isomorphic}
au diagramme ci-dessus qui montre que $(E_i)$ et $(B_n)$ sont pro-isomorphes.
Pour (2) nous pouvons appliquer le lemme \ref{lemma-cartesian-qis-equivalence}
au système projectif des quasi-isomorphismes $E_i \to R/(f_1^{n_i}, \ldots, f_r^{n_i})$.
Pour (3) nous pouvons appliquer le lemme \ref{lemma-pro-isomorphic}
et le fait élémentaire que les systèmes projectifs $(R/I^n)$
et $(R/(f_1^{n_i}, \ldots, f_r^{n_i})$ sont pro-isomorphes.

\medskip\noindent
Comme dans le premier alinéa de la preuve, nous pouvons définir les
foncteurs adjoints\footnote{Il peut être démontré que ces foncteurs
sont, par les équivalences ci-dessus, compatibles avec $F$ et $G$ définis précédemment.}
$$
F' : D_{comp}(R, I) \longrightarrow \mathit{QC}(\mathcal{B})
\quad\text{et}\quad
G' : \mathit{QC}(\mathcal{B}) \longrightarrow D_{comp}(R, I).
$$
Le premier envoie un objet $K$ de $D_{comp}(R, I)$ représenté par un
complexe K-plat $K^\bullet$ à l'objet $(K^\bullet \otimes_R B_n)$ de
$\mathit{QC}(\mathcal{B})$. Le second envoie un objet
$(M_n)$ de $\mathit{QC}(\mathcal{B})$ à $R\lim M_n$. En raisonnant
comme ci-dessus, il suffit de montrer que le noyau de $G'$ est nul. Soit donc
$\mathcal{M} = (M_n)$ un bon faisceau de modules différentiels gradués
sur $\mathcal{B}$ représentant un objet de
$\mathit{QC}(\mathcal{B})$ dans le noyau de $G'$. Alors
$$ 0 =
R\lim M_n \Rightarrow 0 =
(R\lim M_n) \otimes_R^\mathbf{L} B_t =
R\lim (M_n \otimes_R^\mathbf{L} B_t)
$$ Par le lemme
\ref{lemma-different-tensors} nous avons
$R\lim (M_n \otimes_R^\mathbf{L} B_t) =
R\lim (M_n \otimes_{B_n}^\mathbf{L} B_t)$. Puisque
$(M_n)$ est un objet de $\mathit{QC}(\mathcal{B})$, nous voyons que le système
projectif $M_n \otimes_{B_n}^\mathbf{L} B_t$ est constant à partir d'un
certain rang, de valeur $M_t$. Par conséquent, $M_t = 0$, comme voulu.
\end{proof}

\begin{remark}
\label{remark-proregular} Soit
$R$ un anneau et soit $f_1, \ldots, f_r \in R$ une suite d'éléments engendrant un idéal
$I$. Soit $K_n$ le complexe de Koszul sur
$f_1^n, \ldots, f_r^n$, considéré comme une $R$-algèbre différentielle graduée. Nous disons que
$f_1, \ldots, f_r$ est une {\it suite faiblement prorégulière} si, pour tout
$n$, il existe $m > n$ tel que $K_m \to K_n$ induise l'application
nulle en cohomologie, sauf en degré $0$. Si tel est le cas, les arguments de la
preuve de la proposition \ref{proposition-derived-complete} continuent à fonctionner
même lorsque $R$ n'est pas noethérien. En particulier, nous
voyons que $\mathit{QC}(\{R/I^n\})$ est équivalente,
comme catégorie triangulée $R$-linéaire, à la catégorie
$D_{comp}(R, I)$ des objets complets au sens dérivé, à condition que
$I$ puisse être engendré par une suite faiblement prorégulière. Si cela s'avère
nécessaire, nous en donnerons ici un énoncé et une démonstration précis.
\end{remark}








\begin{multicols}{2}[\section{Autres chapitres}]
\noindent
Préliminaires
\begin{enumerate}
\item \hyperref[introduction-section-phantom]{Introduction}
\item \hyperref[conventions-section-phantom]{Conventions}
\item \hyperref[sets-section-phantom]{Théorie des ensembles}
\item \hyperref[categories-section-phantom]{Catégories}
\item \hyperref[topology-section-phantom]{Topologie}
\item \hyperref[sheaves-section-phantom]{Faisceaux sur les espaces}
\item \hyperref[sites-section-phantom]{Sites et faisceaux}
\item \hyperref[stacks-section-phantom]{Champs}
\item \hyperref[fields-section-phantom]{Corps}
\item \hyperref[algebra-section-phantom]{Algèbre commutative}
\item \hyperref[brauer-section-phantom]{Groupes de Brauer}
\item \hyperref[homology-section-phantom]{Algèbre homologique}
\item \hyperref[derived-section-phantom]{Catégories dérivées}
\item \hyperref[simplicial-section-phantom]{Méthodes simpliciales}
\item \hyperref[more-algebra-section-phantom]{Compléments d'algèbre}
\item \hyperref[smoothing-section-phantom]{Lissage des morphismes d'anneaux}
\item \hyperref[modules-section-phantom]{Faisceaux de modules}
\item \hyperref[sites-modules-section-phantom]{Modules sur les sites}
\item \hyperref[injectives-section-phantom]{Objets injectifs}
\item \hyperref[cohomology-section-phantom]{Cohomologie des faisceaux}
\item \hyperref[sites-cohomology-section-phantom]{Cohomologie sur les sites}
\item \hyperref[dga-section-phantom]{Algèbre différentielle graduée}
\item \hyperref[dpa-section-phantom]{Algèbre à puissances divisées}
\item \hyperref[sdga-section-phantom]{Faisceaux différentiels gradués}
\item \hyperref[hypercovering-section-phantom]{Hyperrecouvrements}
\end{enumerate}
Schémas
\begin{enumerate}
\setcounter{enumi}{25}
\item \hyperref[schemes-section-phantom]{Schémas}
\item \hyperref[constructions-section-phantom]{Constructions de schémas}
\item \hyperref[properties-section-phantom]{Propriétés des schémas}
\item \hyperref[morphisms-section-phantom]{Morphismes de schémas}
\item \hyperref[coherent-section-phantom]{Cohomologie des schémas}
\item \hyperref[divisors-section-phantom]{Diviseurs}
\item \hyperref[limits-section-phantom]{Limites de schémas}
\item \hyperref[varieties-section-phantom]{Variétés}
\item \hyperref[topologies-section-phantom]{Topologies sur les schémas}
\item \hyperref[descent-section-phantom]{Descente}
\item \hyperref[perfect-section-phantom]{Catégories dérivées de schémas}
\item \hyperref[more-morphisms-section-phantom]{Compléments sur les morphismes}
\item \hyperref[flat-section-phantom]{Compléments sur la platitude}
\item \hyperref[groupoids-section-phantom]{Schémas en groupoïdes}
\item \hyperref[more-groupoids-section-phantom]{Compléments sur les schémas en groupoïdes}
\item \hyperref[etale-section-phantom]{Morphismes étales de schémas}
\end{enumerate}
Thèmes de théorie des schémas
\begin{enumerate}
\setcounter{enumi}{41}
\item \hyperref[chow-section-phantom]{Homologie de Chow}
\item \hyperref[intersection-section-phantom]{Théorie de l'intersection}
\item \hyperref[pic-section-phantom]{Schémas de Picard des courbes}
\item \hyperref[weil-section-phantom]{Théories cohomologiques de Weil}
\item \hyperref[adequate-section-phantom]{Modules adéquats}
\item \hyperref[dualizing-section-phantom]{Complexes dualisants}
\item \hyperref[duality-section-phantom]{Dualité pour les schémas}
\item \hyperref[discriminant-section-phantom]{Discriminants et différentes}
\item \hyperref[derham-section-phantom]{Cohomologie de de Rham}
\item \hyperref[local-cohomology-section-phantom]{Cohomologie locale}
\item \hyperref[algebraization-section-phantom]{Géométrie algébrique et formelle}
\item \hyperref[curves-section-phantom]{Courbes algébriques}
\item \hyperref[resolve-section-phantom]{Résolution des surfaces}
\item \hyperref[models-section-phantom]{Réduction semi-stable}
\item \hyperref[functors-section-phantom]{Foncteurs et morphismes}
\item \hyperref[equiv-section-phantom]{Catégories dérivées de variétés}
\item \hyperref[pione-section-phantom]{Groupes fondamentaux de schémas}
\item \hyperref[etale-cohomology-section-phantom]{Cohomologie étale}
\item \hyperref[crystalline-section-phantom]{Cohomologie cristalline}
\item \hyperref[proetale-section-phantom]{Cohomologie pro-étale}
\item \hyperref[relative-cycles-section-phantom]{Cycles relatifs}
\item \hyperref[more-etale-section-phantom]{Compléments de cohomologie étale}
\item \hyperref[trace-section-phantom]{La formule des traces}
\end{enumerate}
Espaces algébriques
\begin{enumerate}
\setcounter{enumi}{64}
\item \hyperref[spaces-section-phantom]{Espaces algébriques}
\item \hyperref[spaces-properties-section-phantom]{Propriétés des espaces algébriques}
\item \hyperref[spaces-morphisms-section-phantom]{Morphismes d'espaces algébriques}
\item \hyperref[decent-spaces-section-phantom]{Espaces algébriques décents}
\item \hyperref[spaces-cohomology-section-phantom]{Cohomologie des espaces algébriques}
\item \hyperref[spaces-limits-section-phantom]{Limites d'espaces algébriques}
\item \hyperref[spaces-divisors-section-phantom]{Diviseurs sur les espaces algébriques}
\item \hyperref[spaces-over-fields-section-phantom]{Espaces algébriques sur des corps}
\item \hyperref[spaces-topologies-section-phantom]{Topologies sur les espaces algébriques}
\item \hyperref[spaces-descent-section-phantom]{Descente et espaces algébriques}
\item \hyperref[spaces-perfect-section-phantom]{Catégories dérivées d'espaces}
\item \hyperref[spaces-more-morphisms-section-phantom]{Compléments sur les morphismes d'espaces}
\item \hyperref[spaces-flat-section-phantom]{Platitude sur les espaces algébriques}
\item \hyperref[spaces-groupoids-section-phantom]{Groupoïdes dans les espaces algébriques}
\item \hyperref[spaces-more-groupoids-section-phantom]{Compléments sur les groupoïdes dans les espaces}
\item \hyperref[bootstrap-section-phantom]{Amorçage}
\item \hyperref[spaces-pushouts-section-phantom]{Sommes amalgamées d'espaces algébriques}
\end{enumerate}
Thèmes de géométrie
\begin{enumerate}
\setcounter{enumi}{81}
\item \hyperref[spaces-chow-section-phantom]{Groupes de Chow des espaces}
\item \hyperref[groupoids-quotients-section-phantom]{Quotients de groupoïdes}
\item \hyperref[spaces-more-cohomology-section-phantom]{Compléments sur la cohomologie des espaces}
\item \hyperref[spaces-simplicial-section-phantom]{Espaces simpliciaux}
\item \hyperref[spaces-duality-section-phantom]{Dualité pour les espaces}
\item \hyperref[formal-spaces-section-phantom]{Espaces algébriques formels}
\item \hyperref[restricted-section-phantom]{Algébrisation des espaces formels}
\item \hyperref[spaces-resolve-section-phantom]{Résolution des surfaces revisitée}
\end{enumerate}
Théorie des déformations
\begin{enumerate}
\setcounter{enumi}{89}
\item \hyperref[formal-defos-section-phantom]{Théorie formelle des déformations}
\item \hyperref[defos-section-phantom]{Théorie des déformations}
\item \hyperref[cotangent-section-phantom]{Le complexe cotangent}
\item \hyperref[examples-defos-section-phantom]{Problèmes de déformation}
\end{enumerate}
Champs algébriques
\begin{enumerate}
\setcounter{enumi}{93}
\item \hyperref[algebraic-section-phantom]{Champs algébriques}
\item \hyperref[examples-stacks-section-phantom]{Exemples de champs}
\item \hyperref[stacks-sheaves-section-phantom]{Faisceaux sur les champs algébriques}
\item \hyperref[criteria-section-phantom]{Critères de représentabilité}
\item \hyperref[artin-section-phantom]{Axiomes d'Artin}
\item \hyperref[quot-section-phantom]{Espaces de Quot et de Hilbert}
\item \hyperref[stacks-properties-section-phantom]{Propriétés des champs algébriques}
\item \hyperref[stacks-morphisms-section-phantom]{Morphismes de champs algébriques}
\item \hyperref[stacks-limits-section-phantom]{Limites de champs algébriques}
\item \hyperref[stacks-cohomology-section-phantom]{Cohomologie des champs algébriques}
\item \hyperref[stacks-perfect-section-phantom]{Catégories dérivées de champs}
\item \hyperref[stacks-introduction-section-phantom]{Introduction aux champs algébriques}
\item \hyperref[stacks-more-morphisms-section-phantom]{Compléments sur les morphismes de champs}
\item \hyperref[stacks-geometry-section-phantom]{La géométrie des champs}
\end{enumerate}
Thèmes de théorie des espaces de modules
\begin{enumerate}
\setcounter{enumi}{107}
\item \hyperref[moduli-section-phantom]{Champs de modules}
\item \hyperref[moduli-curves-section-phantom]{Espaces de modules de courbes}
\end{enumerate}
Divers
\begin{enumerate}
\setcounter{enumi}{109}
\item \hyperref[examples-section-phantom]{Exemples}
\item \hyperref[exercises-section-phantom]{Exercices}
\item \hyperref[guide-section-phantom]{Guide bibliographique}
\item \hyperref[desirables-section-phantom]{Desiderata}
\item \hyperref[coding-section-phantom]{Style de programmation}
\item \hyperref[obsolete-section-phantom]{Obsolète}
\item \hyperref[fdl-section-phantom]{Licence de documentation libre GNU}
\item \hyperref[index-section-phantom]{Index généré automatiquement}
\end{enumerate}
\end{multicols}


\bibliography{my}
\bibliographystyle{amsalpha}

\end{document}
